% WYMIARY
\documentclass[11pt,a4paper,twoside]{report}
\usepackage[a4paper,top=3cm, bottom=3cm, outer=2.75cm, inner=3.75cm]{geometry}
\usepackage[unicode=true,hidelinks,pdfpagelayout=TwoPageRight]{hyperref}

% JĘZYKI
\usepackage[english,greek,russian,finnish,polish]{babel} %OSTATNI DOMYŚLNYM
\usepackage[euler]{textgreek}
\usepackage[cp1251,utf8x]{inputenc}
\usepackage[T2A, T1]{fontenc}
\fontencoding{T2A}\fontfamily{Droid}\selectfont 
%\usepackage{libertine}
\usepackage[default]{droidserif}
%\newfontfamily\russianfont[Script=Cyrillic]{Linux Libertine}

\usepackage{paratype} % tlmgr install paratype
\usepackage{tgpagella}
%\usepackage{fontspec}
%\usepackage{polyglossia}

\usepackage{substitutefont}
\substitutefont{T2A}{\rmdefault}{PTSerif-TLF}


\usepackage{polski}
\usepackage{microtype}
\usepackage{breakurl}
\usepackage{pifont}
\usepackage{paracol}
\usepackage{parallel}
\usepackage{array}
\usepackage{tabu}

%LISTY DOMEN
\usepackage{scrextend}
%\addtokomafont{labelinglabel}{\sffamily}


\usepackage[final]{pdfpages}
\usepackage{longtable}

\makeatletter
     \renewcommand*\l@figure{\@dottedtocline{1}{1em}{3.2em}}
     \renewcommand*\l@table{\@dottedtocline{1}{1em}{3.2em}}
\makeatother

\usepackage{lipsum}
\usepackage{mathcomp}
\usepackage{textcomp}
\usepackage{nag}
\usepackage{marvosym}
\renewcommand*{\figurename}{Ryc.}
\renewcommand*{\tablename}{Tab.}
\frenchspacing
\raggedbottom

\hypersetup{
pdftitle={Młynek szaleńców: Prototyp i stereotyp trickstera w fińskim eposie narodowym},
pdfauthor={Karol Nowicki},
pdfsubject={Tricksterzy w Kalevali},
pdfkeywords={głupota} {językoznawstwo} {Kalevala} {kognitywizm} {prototyp} {stereotyp} {szaleństwo} {trickster},
pdfstartview={FitH},
bookmarksopen=true
}


\usepackage[shortlabels]{enumitem}


% PRZEPROGRAMOWANIE DOMYŚLNYCH Z PLAIN NA FANCY/EMPTY
\usepackage{etoolbox}
\patchcmd{\chapter}{\thispagestyle{plain}}{\thispagestyle{fancy}}{}{}
\patchcmd{\part}{\thispagestyle{plain}}{\thispagestyle{empty}}{}{}

\usepackage{emptypage} % MAKES THE DOUBLEPAGE BLANK
\usepackage{fancyhdr}
%\fancyfoot{}
\renewcommand{\headrulewidth}{0pt}
\renewcommand{\footrulewidth}{0pt}
\fancyfoot[LE,RO]{\thepage}
\lhead{}
\chead{}
\rhead{}
%\lfoot{} % get rid of the page number 
\cfoot{}
%\rfoot{}

\usepackage{graphicx} % BEFORE EPSTOPDF!!!
\usepackage{epstopdf} % AFTER GRAPHICX!!!

\usepackage{xfrac}
\newcommand*\rfrac[2]{{}^{#1}\!/_{#2}}

\usepackage{wasysym}
\usepackage{multirow}
\usepackage{latexsym}
\usepackage{changepage}
\usepackage{indentfirst}
\usepackage{setspace}
\usepackage{url}
\usepackage{natbib}
\usepackage{tikz}
\usepackage{tikz-qtree}
\usepackage{forest}
\usepackage{epigraph}
\usepackage[justification=centering]{caption}
\usepackage{minitoc}
\usepackage{fancyhdr}

\renewcommand\thesection{\thechapter.\arabic{section}.}
\renewcommand\thesubsection{\thesection\arabic{subsection}.}
\renewcommand\thesubsubsection{\thesubsection\arabic{subsubsection}.}
\renewcommand\theparagraph{\thesubsubsection\arabic{paragraph}.}
\renewcommand\thesubparagraph{\theparagraph\arabic{subparagraph}.}

\makeatletter
\csdef{forest@compute@node@boundary@rounded rectangle}{%
  \forest@mt{east}%
  \forest@lt{north east}%
  \forest@lt{north}%
  \forest@lt{north west}%
  \forest@lt{west}%
  \forest@lt{south west}%
  \forest@lt{south}%
  \forest@lt{south east}%
  \forest@lt{east}%
}
\makeatother

\usetikzlibrary{arrows,positioning,automata,fit,shapes,trees}

\makeatletter
\newcommand{\ChapterOutsidePart}{%
   \def\toclevel@chapter{-1}\def\toclevel@section{0}\def\toclevel@subsection{1}}
\newcommand{\ChapterInsidePart}{%
   \def\toclevel@chapter{0}\def\toclevel@section{1}\def\toclevel@subsection{2}}
\makeatother

%\overfullrule=5pt

\addto\captionspolish{%
      \renewcommand\listfigurename{Spis rycin}}
\addto\captionspolish{%
      \renewcommand\listtablename{Spis tabel}}
\addto\captionspolish{\renewcommand{\figurename}{Ryc.}}
\addto\captionspolish{\renewcommand{\tablename}{Tab.}}

%CZCIONKA
%\renewcommand*\rmdefault{ppl}

%\usepackage{times}
%\usepackage{charter}
\usepackage{XCharter}
\renewcommand{\bibsection}{}

%\usepackage{mathptmx}

\makeatletter

\pdfinfo{
   /Author (Karol Nowicki)
   /Title  ()
   /CreationDate (D:20140320000000)
   /Subject ()
   /Keywords ()
}

\author{Karol Nowicki}
\title{\textbf{{\huge\ MŁYNEK SZALEŃCÓW:} \\ prototyp i stereotyp trickstera \\ w fińskim eposie narodowym}}
\date{Warszawa 2014}

\makeatletter
\newcommand{\linia}{\rule{\linewidth}{0.4mm}}
\renewcommand{\maketitle}{\begin{titlepage}
    %\vspace*{1cm}
        \begin{center}
        \includegraphics[scale=0.4]{LOGOEPS.eps}
        \end{center}
        \vspace{0.1cm}
        \begin{center}\large
    Uniwersytet Warszawski\\
    Wydział Neofilologii\\
    Katedra Hungarystyki
    \end{center}
    \vspace{0.5cm}
    \begin{center}
    \Large {\@author} \par     %\textsc
    \normalsize{Nr albumu: 281432} %\textsc{281432} \par
    \end{center}
    \vspace{0.75cm}
    \noindent\linia
    \begin{center}
      \LARGE \textsc{\@title}
         \end{center}
     \linia
    \vspace{0.75cm}
  %\vspace{.75cm}
  \begin{center}
\large{\textsc{praca magisterska}}\\
\vspace{0.25cm}
\large {Studia drugiego stopnia\\
        na kierunku \textsc{filologia ugrofińska}\\
        specjalność: filologia ugrofińska --- język fiński}
 \end{center}
 %\vspace{1.5cm}
  \vspace*{\stretch{4}}
     \begin{flushleft}
     \large{
     \hspace{8cm} Praca wykonana pod kierunkiem:

     \hspace{8cm} dr Anna Buncler

     \hspace{8cm} Wydział Neofilologii

     \hspace{8cm} Katedra Hungarystyki}
     \end{flushleft}

 \vspace*{\stretch{6}}
 \begin{center}\large
 \@date
 \end{center}

\newpage\null\thispagestyle{empty}\newpage
\newpage
\pagestyle{empty}
\vspace*{2\baselineskip}
\noindent
\textbf{Oświadczenie kierującego pracą} \newline\newline
\noindent
Oświadczam, że niniejsza praca została przygotowana pod moim kierunkiem i~stwierdzam, że spełnia ona warunki do przedstawienia jej w~postępowaniu o~nadanie tytułu zawodowego.\newline\newline
\newline
Data: 20.03.2014 \hspace{4cm} Podpis kierującego pracą: /--/ \newline\newline
\newline
{}\newline\newline
\newline\newline
\newline
\newline
\textbf{Oświadczenie autora pracy} \newline\newline
\noindent
Świadom odpowiedzialności prawnej oświadczam, że niniejsza praca dyplomowa (magisterska) została napisana przeze mnie samodzielnie i~nie zawiera treści uzyskanych w~sposób niezgodny z~obowiązującymi przepisami.\newline\newline
Oświadczam również, że przedstawiona praca nie była wcześniej przedmiotem procedur związanych z~uzyskaniem tytułu zawodowego w~wyższej uczelni.\newline\newline
Oświadczam ponadto, że niniejsza wersja pracy jest identyczna z~załączoną wersją elektroniczną.\newline\newline
\newline
Data: 20.03.2014 \hspace{4cm} Podpis autora pracy: /--/

\newpage\null\thispagestyle{empty}\newpage
\newpage
\pagestyle{empty}
%\vspace*{6\baselineskip}
%\vspace*{2\baselineskip}
%\vspace*{1\baselineskip}
\vspace*{\baselineskip}
\newlength\iiindent \setlength\iiindent{15pt}
\newcommand\secindent{\parshape 3 0pt \linewidth 0pt\dimexpr\linewidth-\iiindent\relax 0pt \linewidth }
\setlength\parindent{0pt}

\begin{center}
\textbf{Streszczenie}\\
\end{center}
\vspace*{-0.5\baselineskip}
\indent
Zaletą podejścia kognitywnego w~językoznawstwie jest możliwość przeprowadzenia badań lingwistycznych nad zjawiskami i~pojęciami, które często są niejednoznaczne i~rozmyte. Do takich pojęć należy \emph{trickster} --- boski szaleniec, którego obecność w~\emph{Kalevali} jest nader dobrze widoczna, choć temat ten nie został dotychczas podjęty na polu lingwistycznym. Dokładna analiza domen tricksterskich, poprzedzona szczegółowym przedstawieniem nauk o~poznaniu oraz zróżnicowanej sylwetki przechery, udowadnia, iż \emph{Kalevala} pełna jest postaci szelmowskich, których cechuje przede wszystkim magia, głupota oraz więź z~bóstwem. \begingroup \leftskip=0cm plus 0.5fil \rightskip=0cm plus -0.5fil \parfillskip=0cm plus 1fil
Tendencje w~dystrybucji domen pozwalają na rozróżnienie kalevalicznych tricksterów jako \emph{większych} i~\emph{mniejszych}, a~ich cechy dodatkowe umożliwiają wyprowadzenie pojęcia \emph{nowego trickstera kalevalicznego}.\par \endgroup
\vspace{0.75cm}
\begin{center}
\textbf{Słowa kluczowe}\\
\vspace*{0.25\baselineskip}
głupota, językoznawstwo, Kalevala, kognitywizm,\\
prototyp, stereotyp, szaleństwo, trickster \\
\vspace{1cm}
\textbf{Dziedzina pracy (według programu Socrates-Erasmus)}\\
\vspace*{0.25\baselineskip}
09.0~Języki obce i~filologie \\
\vspace{1cm}
\textbf{Klasyfikacja tematyczna}\\
\vspace*{0.25\baselineskip}
151100~Metody stosowane w~językoznawstwie  \\
\vspace{1cm}
\textbf{The Madmen's Mill:}\\
\vspace*{0.25\baselineskip}
A prototype and a stereotype of a trickster-figure in the Finnish national epic\\
\end{center}
\vspace{0.25cm}
%%%%%%%%%%%%%%%%%%%%%
\begin{center}
\textbf{Summary}\\
\end{center}
\vspace*{-0.5\baselineskip}
\indent
%{\selectlanguage{english}
The approach of cognitive linguistics gives us a~possibility of carrying out research on concepts and phenomena, which are often ambiguous and fuzzy. One of them is \emph{trick\-ster} --- the divine madman, whose presence in The Kalevala is clearly visible, al\-though this subject has not yet been studied in the field of linguistics. A~thorough analysis of trickster domains, preceded by a~detailed presentation of cognitive science and complicated nature of the fool, proves that The Kalevala is full of trickster-figures, which are characterized mainly by magic, stupidity and a~link to a~deity. \begingroup \leftskip=0cm plus 0.5fil \rightskip=0cm plus -0.5fil \parfillskip=0cm plus 1fil The tendencies in the distribution of domains allow us to mark the kalevalic tricksters as \emph{major} and \emph{minor} ones and to create a~concept of a~\emph{new kalevalic trickster} by using their additional characteristics. \par \endgroup
%%%%%%%%%%%%%%%%%%%%
\vspace*{\baselineskip}

\newpage
\newpage\null\thispagestyle{empty}\newpage
\end{titlepage}
}
\makeatother

%\widowpenalty=300
%\clubpenalty=300

%\fancyhf{}

\hyphenation{tso-hatzidis}

\begin{document}
\tikzstyle{every node}=[draw=black,thick,rounded corners,anchor=west]
\tikzstyle{selected}=[draw=red,fill=red!30]
\tikzstyle{optional}=[dashed,fill=gray!50]
\tikzset{>=latex}

\maketitle

\pagestyle{fancy}







\setcounter{page}{5}
\tableofcontents
\listoffigures
\begingroup
\let\clearpage\relax
\listoftables
\endgroup

\cleardoublepage
\chapter*{Wstęp}
\addcontentsline{toc}{chapter}{Wstęp}
\begin{spacing}{1.3}

Niniejsza praca magisterska stawia za cel wykazanie obecności postaci trickstera w~fińskim eposie narodowym --- \emph{Kalevali}. Wybór tematu pracy został podyktowany zainteresowaniami autora dotyczącymi fińskiej literatury, wierzeń ludów Północy, jak również antropologicznego aspektu boskiego szaleńca, odnajdywanego we wszystkich mitologiach świata.

Trickster, zwany również ,,boskim szelmą'', ,,przecherą'' czy po prostu ,,spryciarzem'', jest zarówno postacią mityczną stanowiącą ucieleśnienie figlarnej wyobraźni człowieka, jak i~emanacją bohatera kulturowego o~statusie zdegradowanego boga lub świętego grzesznika. To właśnie boży szaleniec przeszkadza w~stworzeniu idealnego świata, przez co kreuje kulturę i~moralność. Początki badań tricksterskich sięgają końca XIX~wieku, a~na szerszą skalę przybliżył tę postać czytelnikowi amerykańskiemu Paul Radin w~swojej książce \emph{Trickster. Studium mitologii Indian północnoamerykańskich}. W~Polsce tematem zajmuje się m.in. Jerzy Wasilewski. Co się tyczy fińskiego eposu narodowego --- pełen jest on odniesień do tricksteryzmu, chociaż fińska literatura na temat boskiego prestidigitatora jest mocno okrojona.

Część teoretyczna dotycząca kognitywizmu czerpie zarówno z~tekstów badaczy polskich, m.in. Jerzego Bartmińskiego, Renaty Grzegorczykowej, Henryka Kardeli i~Anny Wierzbickiej, jak również z~prac językoznawców zagranicznych pokroju George’a Lakoffa i~Ludwiga Wittgensteina. Bibliografia dotycząca boskiego przechery została natomiast oparta na pracach Paula Radina i~Carla Gustava Junga, a~znaczącym źródłem informacji stała się również praca \emph{Szaleństwo w religiach świata} Jacka Sieradzana, omawiająca problematykę bożej głupoty w~różnych systemach religijnych. Kanwą do badań praktycznych stała się \emph{Kalevala} w~tłumaczeniu Jerzego Litwiniuka. Praca została podzielona na cztery części główne oraz aneks, w~którym przedstawiono tricksterów wymienionych w~niniejszej pracy.
\newpage
\vspace*{1.5\baselineskip}
Pierwsza część pracy wprowadza czytelnika w~ogólny świat nauk kognitywnych --- omówiono w~nim główne założenia nauk o~poznaniu, przedstawiono ich składowe, bazę oraz podział kognitywistyki na symboliczną i~subsymboliczną. Kolejne rozdziały skupiają się już stricte na gramatyce kognitywnej autorstwa Ronalda Langackera, do której swój wielki wkład wnieśli również opisani później badacze: Eleanor Rosch (teoria prototypów), George Lakoff (modele kognitywne) i~Ludwig Wittgenstein (teoria podobieństwa rodzinnego). Omówione zostały również stereotypy i~schematy, ich rozgraniczenie względem prototypów, jak również teoria zbiorów rozmytych Zadeha.

Druga część pracy poświęcona jest zagadnieniu trickstera. Rozdziały od piątego do siódmego przedstawiają mistyfikatora pod względem jego wymiaru antropologicznego, szamańskiego oraz groteskowej seksualności (Wakdjunkaga), bożego posłannictwa i~ambiwalencji (\`{E}\d{s}\`{u}-Elegbara, Hermes) oraz głupoty i~braku jej świadomości (Loki). Rozdział ósmy jest przedstawieniem Jahwe jako bóstwa kryminalnego, noszącego w~sobie cechy tricksterskie, na co zwraca uwagę Alf Hiltebeitel.

Trzecia część pracy stanowi syntezę przedstawionej wcześniej gramatyki kognitywnej z~postacią trickstera. Wnikliwa analiza występowania w~treści \emph{Kalevali} ośmiu domen tricksterskich oraz ich przyporządkowanie poszczególnym bohaterom --- Väinämöinenowi, Joukahainenowi, Ilmarinenowi, czarownicy Louhi, Lemminkäinenowi, Kullervo, Ukko, bogini Ilmatar oraz Marjatcie --- udowadniają, iż postacie eposu są odbiciami boskiego szaleńca. Ze względu na widoczne tendencje w~ilości cech tricksterskich oraz na zróżnicowane znaczenie bohaterów eposu, tricksterzy kalevaliczni podzieleni są na tricksterów \emph{większych} i~\emph{mniejszych}. Analiza dowodzi, iż stereotypowy fiński trickster jest magiem (zaklinaczem) niegrzeszącym mądrością, jednocześnie posiadającym koneksje z~bóstwem lub bóstwami.

Czwarta część pracy stanowi wyprowadzenie pojęcia \emph{nowego trickstera kalevalicznego} w~oparciu o~cechy, które nie przynależą stricte do charakterystyki boskiego głupca, jednakże występują regularnie u~tricksterskich bohaterów \emph{Kalevali}. Zalicza się do nich: postawę antyprometejską, postawę antytricksterską (w~której niemądre uczynki prowadzą do katastrofy lub zwycięstwa zła), pychę oraz ostateczną porażkę bohatera, za przykład której może posłużyć odpłynięcie Väinämöinena z~Kalevali.

\end{spacing}
%\newpage
%\newpage\null\thispagestyle{empty}\newpage

\part{GRAMATYKA KOGNITYWIZMU}
\clearpage{\pagestyle{empty}\cleardoublepage}

\chapter{Nauk kognitywnych zarys ogólny}
\begin{spacing}{1.3}

Mianem kognitywistyki (ang. \textit{cognitive science}, fiń. \textit{kognitiotiede}) określa się interdyscyplinarną naukę o~umyśle. Stanowiąca zunifikowany zbiór zajmujących się systemami poznawczymi nauk szczegółowych kognitywistyka --- zwana również \textit{nauką o~procesach poznawczych} --- jest znana m.in. z~tego, że ,,ma problemy z~samodefinicją'' (Wacewicz~2006). Przedstawiana w~przeszłości jako siatka przecinających się dyscyplin, do których zalicza się psychologię, filozofię, naukę o~ludzkim mózgu, sztuczną inteligencję oraz językoznawstwo, w~dzisiejszym świecie jest o~wiele większą kategorią, do której --- podług słów Terry’ego Winograda --- ,,zalicza się każdego, kto zajmuje się zjawiskami związanymi z~umysłem'' (za: Wacewicz~2006), dlatego też wzbogaca się ją dodatkowo o~antropologię, biologię oraz pedagogikę.

Dzięki dwóm rewolucjom wiedzy o~poznaniu --- których skutkiem było powstanie psychologii empirycznej, a~następnie odwrócenie się od behawioryzmu i~skierowanie w~stronę informatyki --- nauka ta nie posiada jednorodnej tożsamości, tak samo jak nie posiada właściwej tylko jej samej i~ujednoliconej metody badawczej. Dyscypliny weń wchodzące uzupełniają się wzajemnie (lingwistyka współdziała z~logiką, filozofia z~psychologią itd.) oraz posługują wieloma typami reprezentacji mentalnych, takimi jak myśli, pojęcia, modele poznawcze czy schematy wyobrażeniowe. Owe reprezentacje jednoczą ogół nauk poznawczych, działając jako znaki zastępujące elementy rzeczywistości. Dzięki zróżnicowanym dziedzinom (głównie logice, informatyce, filozofii oraz biologii) tworzy się obliczeniowe modele procesów poznawczych, które są następnie testowane poprzez psychologię i~neuronauki.\\
\vfill
\pagebreak
\vspace*{1.5\baselineskip}
\begin{figure}[ht]
\textsc{
\renewcommand*{\figurename}{Rys.}
\begin{center}
\begin{tikzpicture}[%
  grow via three points={one child at (0.5,-0.7) and
  two children at (0.5,-0.7) and (0.5,-1.4)},
  edge from parent path={(\tikzparentnode.south) |- (\tikzchildnode.west)}]
  \node {\textbf{\LARGE kognitywistyka}}
    child [missing] {}	
    child { node {\textbf{antropologia}}}
    child { node {\textbf{filozofia poznania}}
      child { node {filozofia języka}}
      child { node {filozofia percepcji}}
      child { node {filozofia umysłu}}
    }
    child [missing] {}				
    child [missing] {}				
    child [missing] {}
    child [missing] {}
    child { node {{\shortstack[l]{\textbf{informatyka} \\ \textbf{matematyka} \\ \textbf{sztuczna inteligencja}}}}}
    child [missing] {}
    child { node {{\textbf{językoznawstwo kognitywne}}}
      child { node {filozofia}}
      child { node {logika}}
      child { node {psychologia}}
    }
    child [missing] {}				
    %child [missing] {}				
    %child [missing] {}		
    child { node {{\shortstack[l]{\textbf{lingwistyka} \\ \textbf{logika}}}}}
    child [missing] {}
    child { node {{\textbf{neuronauki}}}
      child { node {biochemia}}
      child { node {neurobiologia}}
      child { node {neurofizjologia}}
      child { node {neurologia}}
      child { node {neuropsychologia}}
      %child [missing] {}
    }
    child [missing] {}				
    child [missing] {}	
    child [missing] {}
    child [missing] {}			
    child [missing] {}
    child [missing] {}
    child { node {{\textbf{psychologia}}}
      child { node {psychologia kognitywna (poznania)}}
      child { node {psychologia rozwoju dziecka}}
      child [missing] {}
    }
    child [missing] {}				
    child [missing] {}		
    child [missing] {}		
    child [missing] {};
   \end{tikzpicture}
\end{center}
}
\caption{Dyscypliny składowe kognitywistyki}
\end{figure}

\pagebreak

\begin{figure}[ht]
\textsc{
\begin{center}
\begin{tikzpicture}[%
  grow via three points={one child at (0.5,-0.7) and
  two children at (0.5,-0.7) and (0.5,-1.4)},
  edge from parent path={(\tikzparentnode.south) |- (\tikzchildnode.west)}]
%  edge from parent path={([xshift=2ex]\tikzparentnode.south west) |- (\tikzchildnode.west)},
%  growth parent anchor=west]
  \node {{\Large \textbf{baza kognitywna}}}
    child [missing] {}	
    child { node {{\textbf{ai}}}
      child { node {filozofia}}
      child { node {informatyka}}
      child { node {logika}}
      child { node {matematyka}}
    }
    child [missing] {}				
    child [missing] {}				
    child [missing] {}		
    child [missing] {}
    %child [missing] {}			
    child { node {{\textbf{neuronauki}}}
      child { node {biologia}}
      child { node {informatyka}}
      child { node {psychologia}}
    }
    child [missing] {}				
    child [missing] {}				
    child [missing] {}	
    %child [missing] {}	
    child { node {{\textbf{językoznawstwo kognitywne}}}
      child { node {filozofia}}
      child { node {logika}}
      child { node {psychologia}}
    }
    child [missing] {}				
    child [missing] {}				
    child [missing] {};
\end{tikzpicture}
\end{center}
}
\caption{Baza kognitywizmu}
\end{figure}
\vspace*{0.5\baselineskip}
Głównym celem nauk kognitywnych jest spojrzenie na ogólne zagadnienie poznania, m.in. poprzez badanie ludzkich możliwości empirycznych, samej struktury procesu myślenia oraz sposobu powstawania i~funkcjonowania umysłu (Wacewicz~2006) oraz ,,stworzenie zunifikowanego programu badawczego pozwalającego identyfikować różnorodne formy aktywności poznawczej i~wyjaśnić rolę, jaką pełnią one w~funkcjonowaniu umysłu jako całości'' (Klawiter~2006). Punktem wyjścia kognitywistyki jest ,,analiza sposobu postrzegania świata, próba zrozumienia, co dzieje się w~umyśle, gdy wykonujemy najprostsze czynności umysłowe'' (Duch~1998), a~jej drugim celem --- zbudowanie podobnego ludzkiemu symulatora sztucznego umysłu, posiadającego inteligencję i~zdolność myślenia. Jest to bezpośrednio związane z~faktem, iż poza człowiekiem kognitywizm uwzględnia w swojej ścieżce badawczej komputery jako systemy sztuczne (zaowocowało to m.in. stworzeniem jednego z~podstawowych problemów tej dziedziny, czyli zagadnienia maszyn myślących), czego podwaliny położył już David Hume w~swojej próbie stworzenia mechaniki umysłu na wzór newtonowskiej mechaniki klasycznej.\\

\pagebreak

\vspace*{1\baselineskip}
\begin{figure}[h]
\textsc{
\begin{center}
\begin{tikzpicture}[%
  grow via three points={one child at (0.5,-0.7) and
  two children at (0.5,-0.7) and (0.5,-1.4)},
  edge from parent path={(\tikzparentnode.south) |- (\tikzchildnode.west)}]
 % edge from parent path={([xshift=2ex]\tikzparentnode.south west) |- (\tikzchildnode.west)},
 % growth parent anchor=west]
  \node{{\shortstack[c]{\Large \textbf{kognitywistyka}\\ \Large \textbf{symboliczna}}}}
    child [missing] {}	
    child { node {{inteligencja kognitywna}}}
    child { node {{kognitywne podejmowanie decyzji}}}
    child { node {{modelowanie abstrakcyjnych funkcji myślowych}}}
    child { node {{rozwijanie inżynierii socjokognitywistycznej}}}
    child { node {{samoświadomość i definicja świadomości}}}
    child [missing] {}				
    child [missing] {}				
    child [missing] {}		
    child [missing] {}
    child [missing] {};
\end{tikzpicture}
\end{center}
}
\caption{Zadania kognitywistyki symbolicznej}
\end{figure}
\vspace{1cm}
\begin{figure}[ht]
\textsc{
\begin{center}
\begin{tikzpicture}[%
  grow via three points={one child at (0.5,-0.7) and
  two children at (0.5,-0.7) and (0.5,-1.4)},
  edge from parent path={(\tikzparentnode.south) |- (\tikzchildnode.west)}]
 % edge from parent path={([xshift=2ex]\tikzparentnode.south west) |- (\tikzchildnode.west)},
 % growth parent anchor=west]
  \node{{\shortstack[c]{\Large \textbf{kognitywistyka}\\ \Large \textbf{subsymboliczna}}}}
    child [missing] {}	
    child { node {{\shortstack[c]{łączenie wzorców nauk o~poznawaniu, \\psychologii, neurobiologii i~innych}}}}
    child [missing] {}	
    child { node {{\shortstack[c]{wyjaśnianie procesów myślowych opartych \\na własnościach sieci neuronowych}}}}
    child [missing] {}
    child [missing] {}	
    child [missing] {};
\end{tikzpicture}
\end{center}
}
\caption{Zadania kognitywistyki subsymbolicznej}
\end{figure}
\vspace{1.5cm}
Dzięki badaniom kognitywistycznym rozwija się poznanie procesów takich jak widzenie, zestawianie obiektów w~klasy oraz rozumienie wypowiedzi. Dodatkowo podejście poznawcze pomaga zrozumieć mechanizmy działania emocji, istotę świadomości i~roli podmiotowości w~empirii.\\

\end{spacing}

\chapter{Gramatyka kognitywna stricte}
\begin{spacing}{1.3}

Powstała na przełomie lat siedemdziesiątych i~osiemdziesiątych XX~wieku gramatyka kognitywna należy do młodszych teorii lingwistycznych oraz stanowi swoistą przeciwwagę dla formalnych teorii w~językoznawstwie, zwłaszcza dla tzw. \emph{rozszerzonej teorii standardowej} (ang. \textit{extended standard theory}, fiń. \textit{laajennettu standarditeoria}) Noama Chomsky’ego.

Twórcą gramatyki kognitywnej jest amerykański językoznawca Ronald Langacker, który w~swoim czasie zajmował się również rozwojem jej poprzedniczki --- gramatyki przestrzeni (ang. \textit{space grammar}, fiń. \textit{tilakielioppi}) pod koniec lat siedemdziesiątych. Datowanie początków kognitywizmu jest umowne, ponieważ już trzydzieści lat przed Langackerem elementy tejże analizy były dobrze widoczne w~niemieckiej teorii gestaltów (\textit{Gestalt Theorie}), a~później również w~pracach lingwistów uprawiających semantykę generatywną (m.in. George'a Lakoffa) oraz w~tekstach Ludwiga Wittgensteina (por.~1972), dotyczących teorii podobieństwa rodzinnego. Warto również zwrócić uwagę na pracę z~dziedziny psychologii Eleanor Rosch (1977:~1--72), mówiącą o~teorii prototypów.

W~latach dziewięćdziesiątych kognitywizm uległ zróżnicowaniu --- mamy do czynienia z~pracami o~charakterze formalistycznym (Jackendoff~1983;~1987), z~tekstami Langackera traktującymi o~kognitywizmie w~sposób określany przez Henryka Kardelę jako ,,umiarkowany formalizm'', z~pracami analitycznymi (por.~Turner~1987) oraz lingwi\-styczno-filozoficznymi (por.~Johnson~1987).

\newpage
%\vspace*{0.25\baselineskip}

\section{Encyklopedyczne pojęcie znaczenia}

Gramatyka kognitywna wprowadza w~językoznawstwie brak wyraźnego rozgraniczenia pomiędzy poziomami, za pomocą których opisywany jest język. Odrzuca m.in. zmodularyzowany opis języka według EST (por.~Chomsky~1982;~1986a;~1986b), pozostawiając trzy struktury językowe: fonologiczne, semantyczne oraz symboliczne (łączące dwie poprzednie) (Kardela 1992:~10). Zwracająca uwagę na semantyczne motywy zabiegów językowych gramatyka kognitywna stanowi model użycia języka (por. Langacker 1988a:~127--161), a~nie zbiór ściśle określonych zasad generatywnych, które nie mogą zostać naruszone. O~konieczności odrzucenia ścisłych podziałów i~modularności pisze Tomasz Krzeszowski (1988:~15):
\end{spacing}
%\newpage
\begin{adjustwidth}{1,5cm}{1,5cm}
\small{
,,Ponieważ symbolizm języka obejmuje wszystkie jego warstwy, od fonetyki przez leksykę, przeto językoznawstwo kognitywne nie uznaje za celowe, by postulować istnienie odrębnych komponentów języka, w~czym celowało językoznawstwo generatywne. Nie istnieją, na przykład wyraźne granice między składnią a~leksyką z~jednej strony czy składnią a~morfologią z~drugiej. [...] Również fałszywy jest opis języka oparty na wszelkich innych tradycyjnych podziałach i~rozróżnieniach, które zawsze i~nieuchronnie prowadzą do zaniedbania lub wręcz pomijania w~opisie przypadków pośrednich. Językoznawstwo kognitywne odrzuca zatem rozróżnienie między synchronią i~diachronią [...], między kompetencją a~wykonaniem [...], między gramatyką a~leksykonem, między morfologią a~składnią, między składnią a~semantyką, między semantyką a~pragmatyką [...], między denotacją a~konotacją [...].''} \\
\end{adjustwidth}
%\vspace{0.5cm}
\begin{spacing}{1.3}
Wprowadzając natomiast do kognitywistyki konieczność stosowania encyklopedycznej koncepcji znaczenia mówi Langacker (2005:~17):
\end{spacing}
\begin{adjustwidth}{1,5cm}{1,5cm}
\small{
,,Nie istnieje wyraźna granica między wiedzą językową i~nie-językową. Język wykorzystuje ogólną wiedzę o~świecie i~przywołuje tę wiedzę w~pewien sposób. Języka nie da się jednak wyraźnie oddzielić od tej wiedzy ani też nie istnieją wyraźne wyznaczniki tego, jakie elementy wiedzy mogą być przywoływane przez język.''}\\
\end{adjustwidth}
\begin{spacing}{1.3}

Encyklopedyczne pojęcie znaczenia jest w~pełni widoczne poprzez sformułowane przez tego badacza tezy (Langacker~1988b:~49--50): \\

\newpage
%\noindent

\begin{description}
  \item[TEZA A:]
  ,,Struktura semantyczna jest redukowalna do struktury kognitywnej (czyli mentalnego doświadczenia)'' --- struktura semantyczna jest ,,zanurzona'' w~wiedzy o~świecie, a~znaczenia ,,zależne są w~sposób oczywisty od tej wiedzy, czyli od szeroko pojętego kontekstu, tak lingwistycznego, jak i~»pozajęzykowego«'' (por. Langacker~1991:~109).
  \item[TEZA B:] ,,Często używane wyrażenie przywołuje sieć powiązanych ze sobą znaczeń (sensów)'' --- wyrażenia języka cechuje polisemia oraz stanowią one węzeł dostępu do danej sieci znaczeń. Elementy kategorii, które odchodzą od prototypu dają podwaliny tzw. \textit{kategoriom radialnym} (ang. \textit{radial categories}, fiń. \textit{säteittäiset kategoriat}), razem z~elementami prototypowymi (centralnymi) oraz peryferyjnymi (por. Tabakowska 2001:~54).
  \item[TEZA C:] ,,Struktury semantyczne określane są względem domen kognitywnych'' --- dana kategoria jest tym bardziej prototypowa, im więcej modeli z~tzw. \textit{wiązki modeli kognitywnych} (ang. \textit{cluster models}, fiń. \textit{klusterimallit, ryväsmallit}) jest przywoływanych dla zdefiniowania znaczenia danego słowa. Wyidealizowanymi modelami kognitywnymi (ang. \textit{idealized cognitive models}, fiń. \textit{idealisoidut kognitiiviset mallit}) --- na przykładzie modelu matki (m.in. jako ‘najbliższego przodka płci żeńskiej’ w~modelu genealogicznym, ‘żony ojca’ w~modelu małżeńskim czy też ‘dorosłej kobiety, która wychowuje i~jest żywicielką dziecka’ w~myśl modelu żywieniowego) --- zajmuje się m.in. Lakoff (1987: 74).
  \item[TEZA D:] ,,Struktura semantyczna uzyskuje swój status w~wyniku nałożenia profilu (desygnatu) na bazę'' --- wyrażenie zyskuje swoje znaczenie poprzez \emph{profilowanie}, nazywane również \emph{perspektywizacją} (Taylor 2001:~131), stanowiące wyodrębnianie określonych struktur z~tzw. \textit{bazy kognitywnej}. Dobrym przykładem jest proces definiowania ,,przeciwprostokątnej'', który nie ma prawa bytu bez odwołania się do pojęcia figury ,,trójkąta''.
  \item[TEZA E:] ,,Struktury semantyczne zawierają \emph{skonwencjonalizowane obrazowanie}'' \linebreak (ang. \emph{conventional imagery}, fiń. \textit{konventionaalinen mielikuvaisuus}) --- według Langackera (1988b:~63--64) scena (sytuacja) jest konstruowana na wiele różnych sposobów, przy czym zwraca się uwagę na aspekty takie jak: poziom uszczegółowienia, założenia dotyczące tła i~oczekiwań, skala i~zakres predykacji, perspektywa oraz subiektywizacja i~sekundarna aktywizacja. Pierwszy aspekt dotyczy precyzji definicji, widoczny np. w~ciągu \emph{zwierzę} $\textrm{\ding{221}}$ \emph{ssak} $\textrm{\ding{221}}$ \emph{gryzoń} $\textrm{\ding{221}}$ \emph{chomik} $\textrm{\ding{221}}$ \emph{chomik syryjski}, w~którym wyrażenie ostatnie jest najdokładniejsze.
\end{description}
\newpage

%\noindent
%TEZA A: ,,Struktura semantyczna jest redukowalna do struktury kognitywnej (czyli \linebreak mentalnego doświadczenia)'' --- struktura semantyczna jest ,,zanurzona'' w~wiedzy \linebreak o~świecie, a~znaczenia ,,zależne są w~sposób oczywisty od tej wiedzy, czyli od szeroko \linebreak pojętego kontekstu, tak lingwistycznego, jak i~«pozajęzykowego»'' (por. Langacker \linebreak 1991:~109).\\

%\noindent
%TEZA B: ,,Często używane wyrażenie przywołuje sieć powiązanych ze sobą znaczeń (sensów)'' --- wyrażenia języka cechuje polisemia oraz stanowią one węzeł dostępu do danej sieci znaczeń. Elementy kategorii, które odchodzą od prototypu dają podwaliny tzw. \textit{kategoriom radialnym} (ang. \textit{radial categories}), razem z~elementami prototypowymi \linebreak (centralnymi) oraz peryferyjnymi (por. Tabakowska 2001:~54).\\

%\newpage
%\noindent
%TEZA C: ,,Struktury semantyczne określane są względem domen kognitywnych'' --- \linebreak dana kategoria jest tym bardziej prototypowa, im więcej modeli z~tzw. \textit{wiązki modeli kognitywnych} (ang. \textit{cluster models}) jest przywoływanych dla zdefiniowania znaczenia danego słowa. Wyidealizowanymi modelami kognitywnymi (ang. \textit{Idealized Cognitive Models}, \textit{ICMs}) --- na przykładzie modelu matki (m.in. jako ‘najbliższego przodka płci żeńskiej’ w~modelu genealogicznym, ‘żony ojca’ w~modelu małżeńskim czy też ‘dorosłej kobiety, która wychowuje i~jest żywicielką dziecka’ w~myśl modelu żywieniowego) --- zajmuje się m.in. Lakoff (1987: 74).\\

%\noindent
%TEZA D: ,,Struktura semantyczna uzyskuje swój status w~wyniku nałożenia profilu (desygnatu) na bazę'' --- wyrażenie zyskuje swoje znaczenie poprzez \emph{profilowanie}, nazywane również \emph{perspektywizacją} (Taylor 2001:~131), stanowiące wyodrębnianie określonych struktur z~tzw. \textit{bazy kognitywnej}. Dobrym przykładem jest proces definiowania ,,przeciwprostokątnej'', który nie ma prawa bytu bez odwołania się do pojęcia figury ,,trójkąta''.\\

%\noindent
%TEZA E: ,,Struktury semantyczne zawierają \emph{skonwencjonalizowane obrazowanie} (ang. \linebreak \emph{conventional imagery})'' --- według Langackera (1988b:~63--64) scena (sytuacja) jest konstruowana na wiele różnych sposobów, przy czym zwraca się uwagę na aspekty takie jak: poziom uszczegółowienia, założenia dotyczące tła i~oczekiwań, skala i~zakres predykacji, perspektywa oraz subiektywizacja i~sekundarna aktywizacja. Pierwszy aspekt dotyczy precyzji definicji, widoczny np. w~ciągu \emph{zwierzę} $\textrm{\ding{221}}$ \emph{ssak} $\textrm{\ding{221}}$ \emph{gryzoń} $\textrm{\ding{221}}$ \emph{chomik} $\textrm{\ding{221}}$ \emph{chomik syryjski}, w którym wyrażenie ostatnie jest najdokładniejsze.\\

%\newpage
%%%%%%%%%%
%\vspace*{0.5\baselineskip}

Widoczny w~gramatyce kognitywnej jest również brak rozdziału semantyki od pragmatyki --- według Langackera podział wiedzy na językową i~pozajęzykową jest sztuczny i~pociąga za sobą niewłaściwe podejście naukowe --- podczas gdy zdaniem Kardeli (1992:~10) ,,jedyna właściwa koncepcja semantyki językoznawczej to taka, \linebreak która unika fałszywych dychotomii''. Idąc śladem tej myśli napotykamy problem właściwego konstruowania definicji. Dalej za Kardelą (tamże):
\end{spacing}
\begin{adjustwidth}{1,5cm}{1,5cm}
\small{
,,Czy w~definicji »kota« nie powinna np. znaleźć się informacja: »Z~reguły koty towarzyszą czarownicom podczas odbywającego się sabatu czarownic«? Odpowiedź brzmi: i~tak i~nie, w~zależności od okoliczności użycia słowa »kot«. Jeśli będziemy mieli do czynienia z~praktyką etnolingwistyczną, jest oczywiste, że w~opisie czarownicy informacja ta będzie miała kapitalne znaczenie, zupełnie nieistotna natomiast okaże się ona dla weterynarza (o~ile oczywiście nie wierzy on w~istnienie czarownic, w~związek, jaki łączy koty z~czarownicami i~nie daje ujścia swoim przekonaniom w~praktyce weterynaryjnej). Znaczy to, że o~tym, co składa się na informację centralną, a~co na peryferyjną, decyduje konkretne użycie języka, konkretny kontekst --- czas i~miejsce wypowiedzi, przekonania i~wierzenia mówiącego, sposób wartościowania, itd.''}\\
\end{adjustwidth}
\begin{spacing}{1.3}
%\noindent
Konkretne użycie języka ma zatem dla kognitywizmu znaczenie kluczowe, a~zmieniający się punkt widzenia rodem z~gestaltyzmu jest podstawową zasadą gramatyki kognitywnej, w~której ,,kompozycjonalny sposób derywowania znaczeń'' został odrzucony (Kardela 1992:~10).

Podejście kognitywne poświęca również więcej uwagi metaforyzacji oraz idiomatyczności. W~obu przypadkach do poprawnego zrozumienia wyrażeń konieczna jest znajomość całych struktur, a~nie tylko poszczególnych wyrazów (por. Lewicki 1982: 33--47). Za Langackerem (2003:~30):
\end{spacing}
\begin{adjustwidth}{1,5cm}{1,5cm}
\small{
,,Duże znaczenie przywiązuje się [...] do przypadków rzeczywistego użycia systemu językowego oraz do posiadanej przez użytkownika wiedzy na temat tego użycia; to gramatyka jest odpowiedzialna za posiadaną przez użytkowników języka wiedzę dotyczącą pełnego zakresu językowych konwencji [...]. Mamy tu do czynienia z~nieredukcjonistycznym podejściem do struktury języka.''\\
}
\end{adjustwidth}
%\vspace{0.5cm}
\begin{spacing}{1.3}
Metafory to według kognitywistów podstawowe i~kreatywne składniki języka, które generatywiści zepchnęli na margines jako te, do których nie można zastosować określonych reguł.
\newpage
\section{Domeny kognitywne}
Domeny kognitywne oznaczają kognitywne lub konceptualizacyjne struktury, za pomocą których ujmowana jest ludzka wiedza o~świecie. Łączą się one hierarchicznie, \linebreak a~następnie określają daną strukturę semantyczną. Domeny są bezpośrednio związane \linebreak ze wspomnianą wcześniej encyklopedyczną koncepcją znaczenia. Za Langackerem \linebreak (2005:~17):
\end{spacing}
\begin{adjustwidth}{1,5cm}{1,5cm}
\small{
,,Dane słowo reprezentuje konwencjonalny sposób dostępu do zbioru domen. Pewne domeny są istotne dla określenia wartości semantycznej sło\-wa, inne są peryferyjne lub bardziej »elastyczne« i~mogą być przywoływane w~zależności od kontekstu.''\\
}
\end{adjustwidth}
\begin{spacing}{1.3}
Domeny kognitywne dzielimy na podstawowe oraz złożone. Domeny podstawo\-we obejmują kryteria proste, np. wymiar, kolor lub czas. Domeny złożone cechuje natomiast hierarchia oraz subdomeny. Dobrym przykładem takiego zjawiska jest utworzenie domeny ,,kulinarnej''  (która znajduje zastosowanie podczas definiowania słów takich jak \emph{nóż} lub \emph{łyżka}), do której należą w~następnej kolejności subdomeny ,,zastawy stołowej'', ,,ułożenia sztućców na stole'' lub ,,drinków''. Zbiór domen kognitywnych dla danego pojęcia jest dobrze widoczny w~analizie pojęcia \emph{szklanki} według Langackera (2005:~16):\\

\begin{adjustwidth}{1,5cm}{1,5cm}
\begin{labeling}{ukierunkowanieeee}

\item [\textbf{przestrzeń}] domena podstawowa
\item [\textbf{ukierunkowanie}] wydłużenie wertykalne, zamknięcie u~podstawy
\item [\textbf{kształt}] cylinder zamknięty z jednej strony
\item [\textbf{materiał}] zazwyczaj szkło
\item [\textbf{rozmiar}] łatwo da się trzymać w~ręce
\item [\textbf{funkcja}] naczynie na płyn, rola w~procesie picia
\item [\textbf{inne}] metoda wyrobu, cena, łatwość tłuczenia

\end{labeling}
\end{adjustwidth}

%\begin{itemize}
%\item[\Lefttorque] \textsc{przestrzeń --- (domena podstawowa)}
%\item[\ding{51}] \textsc{przestrzeń --- [domena podstawowa]};
%\item[\ding{51}] \textsc{ukierunkowanie} --- wydłużenie wertykalne, zamknięcie u~podstawy;
%\item[\ding{51}] \textsc{kształt} --- cylinder zamknięty z jednej strony;
%\item[\ding{51}] \textsc{materiał} --- zazwyczaj szkło;
%\item[\ding{51}] \textsc{rozmiar} --- łatwo da się trzymać w~ręce;
%\item[\ding{51}] \textsc{funkcja nr 1} --- naczynie na płyn;
%\item[\ding{51}] \textsc{funkcja nr 2} --- rola w~procesie picia;
%\item[\ding{51}] \textsc{inne} --- cena, łatwość tłuczenia, metoda wyrobu.
%\item[ ]
%\end{itemize}
%\begin{spacing}{1.3}

\newpage
\section{Teoria podobieństwa rodzinnego}

Metaforyzacja niesie za sobą dodatkowo rozszerzenie granic kategorii, co jest bezpośrednio związane z~podstawowym problemem dotyczącym ich samej organizacji. Dwie możliwe drogi tejże organizacji to zastosowanie systemu warunków koniecznych oraz wystarczających w~przypadku gramatyk strukturalistycznych i~transforma\-cyjno-generatywnych, lub zasady podobieństwa rodzinnego według Ludwiga Wittgensteina (por.~1953) w~przypadku gramatyki kognitywnej, która to zasada jest --- obok teorii prototypów Eleanor Rosch --- jej punktem oparcia w~opisie kategoryzacji.

Leżąca u~podstaw gramatyki kognitywnej zasada podobieństwa rodzinnego wyróżnia ,,lepsze'' (prototypowe) oraz ,,gorsze'' (mniej prototypowe) egzemplarze w~danej kategorii (Kardela 1992:~11). Pokrewieństwo rodzinne jest w~ostatnich czasach uznawane za wytyczną nowoczesnej myśli lingwistycznej, co komentuje Anna Wierzbicka (1992:~33):
\end{spacing}
\begin{adjustwidth}{1,5cm}{1,5cm}
\small{
,,Wittgensteinowskie pojęcie »podobieństwa rodzinnego« odegrało kolosalną rolę w~rozwoju tzw. »semantyki prototypów« i~przybrało status niekwestionowanego dogmatu w~literaturze poświęconej znaczeniu [...]. Głosi się więc ostatnio hasła »antydefinicyjne« [...]. Wytworzył się wręcz klimat, w~którym ktoś, kto próbuje zdefiniować cokolwiek, może zostać uznany za człowieka staroświeckiego, nie idącego z~duchem czasu [...]. Jeśli więc semantyk ma iść z~duchem czasu, nie może mówić o~definicjach, lecz o~»pokrewieństwie rodzinnym«, prototypach i~rozmazaniu myśli ludzkiej [...]. Głosi się zatem coraz powszechniej, iż słów nie da się zdefiniować, ponieważ znaczenie zakodowane w~języku naturalnym jest w~swej istocie rozmyte.''\\
}
\end{adjustwidth}
\begin{spacing}{1.3}
%\noindent
Wittgenstein (1953: §~66--67) opisuje swoją teorię następująco:
\end{spacing}

\begin{adjustwidth}{1,5cm}{1,5cm}
\small{
,,[...] the result of this examination is: we see a~complicated network of similarities overlapping and criss-crossing: sometimes overall similarities, sometimes similarities of detail. [...] I can think of no better expression to characterize these similarities than »family resemblances«; for the various resemblances between members of a~family: build, features, colour of eyes, gait, temperament, etc. [...] overlap and criss-cross in the same way. [...] And the strength of the thread does not reside in the fact that some one fibre runs through its whole length, but in the overlapping of many fibres. But if someone wished to say: »There is something common to all these constructions --- namely the disjunction of all their common properties« --- I should reply: Now you are only playing with words. One might as well say: »Something runs through the whole thread --- namely the continuous overlapping of those fibres«.''\\
}
\end{adjustwidth}


\newpage

\chapter{Prototyp kognitywny}
\begin{spacing}{1.3}
Lata siedemdziesiąte XX~wieku przyniosły językoznawstwu prawdziwą rewolucję --- oto po raz pierwszy arystotelesowski model warunków koniecznych oraz warunków wystarczających w~celu stworzenia właściwej reprezentacji świata miano zastąpić kategoryzacją gradualną (Golka 2009:~32). Samo pojęcie \emph{prototypu} mogło zaistnieć poprzez odejście od opisu związków między wyrażeniami a~zjawiskami na rzecz opisywania samego ,,rozumienia wyrażeń poprzez społeczność mówiącą, a~więc w~ostateczności na opis pewnego typowego stanu mentalnego mówiących wiązanego z~używanymi wyrażeniami'' (Grzegorczykowa 1998:~109). W~ten sposób semantyka językowa zyskała w~swoim opisie kognitywny wydźwięk psychologiczny. Za Lindą Coleman i~Paulem Kayem (1981:~43, za:~Kleiber 1990:~59):
\end{spacing}
\begin{adjustwidth}{1,5cm}{1,5cm}
\small{
,,Bronimy idei, że wiele wyrazów ma jako znaczenie nie zbiór warunków koniecznych i~wystarczających, które jakaś rzecz lub zjawisko musi spełniać, by należeć do kategorii nazwanej daną nazwą, ale raczej znaczeniem jest przedmiot lub proces psychiczny, który nazwaliśmy prototypem.''
}
\end{adjustwidth}
\begin{spacing}{1.3}

\section{Prototypy Eleanor Rosch}

Eleanor Rosch, w~pracach której pojawiła się teoria prototypu, reprezentuje podejście probabilistyczne stwierdzając, iż istnieją obiekty mniej lub bardziej typowe dla określonej kategorii. Każda kategoria posiada zatem egzemplarze centralne oraz peryferyczne, z~czego te pierwsze pełnią rolę jej typowych reprezentantów, przy czym przejście między nimi pozostaje płynne (Golka 2009:~32), a~same jej granice --- rozmyte (Grzegorczykowa 1998:~109).

\newpage

Prototyp pełni zadanie ukazania zjawiska typowego, podczas gdy przypadki marginalne usuwa się z~pola widzenia. Im bardziej typowy jest pewien egzemplarz, tym większą posiada on wartość w~przynależności do swojej kategorii.

Jednym ze sposobów interpretacji pojęcia \emph{prototypu} jest nawiązanie do jego etymologii (gr. \textgreek{πρώτος} [pr\={o}tos] ‘pierwszy’, ‘wczesny’) i~ukazanie go jako pierwowzoru, który jest następnie powielany i~rozwijany. Według kognitywizmu opiera się prototyp jednakże na odrzuceniu aspektu genetycznego na rzecz oznaczenia go mianem ,,wzorca'' lub ,,wzorowego przykładu czegoś'' (Grzegorczykowa 1998:~110). W~pracach semantycznych można dostrzec trzy główne nurty interpretacji prototypu:\\
\begin{enumerate}[1)]
  \item jako najlepszego przykładu kategorii i przedstawiciela jej zakresu,
  \item jako zbioru cech typowych dla najlepszych przedstawicieli danej kategorii,
  \item jako przeciwstawione peryferyjności centrum znaczeniowe wyrazu.\\
\end{enumerate}
%\noindent

Pierwsza interpretacja jest bardzo dobrze widoczna w~pracach Cecila Browna \linebreak (por. 1990:~17--47) i~Anny Wierzbickiej (por. 1990:~347--367). Wiąże się ona bezpośrednio ze środowiskiem mówiącego, które wyznacza najlepszego przedstawiciela heterogenicznej kategorii posiadającej zróżnicowane cechy (kategoria ta nie ułatwia oczywiście badaczowi znalezienia jednego wspólnego dlań okazu) --- dlatego też prototypowym pojazdem będzie \emph{samochód} lub \emph{pociąg} (miast \emph{bryczki} czy \emph{dorożki}), a~ptakiem --- zależnie od społeczności --- polski \emph{wróbel} lub amerykański \emph{drozd} (Grzegorczykowa 1998: 111--112).

Interpretacja druga to zbiór cech charakterystycznych związanych z~danym pojęciem, przez co zbliża ona prototyp do semantycznego rozumienia wyrazów na drodze potoczności. Używając słowa \emph{matka} mówiący ma na myśli ‘kobietę, która urodziła i~wychowuje dziecko’, nie myśląc o~możliwości rozgraniczenia aspektu genetycznego od wychowawczego (Grzegorczykowa 1998:~112). W~niektórych kręgach językoznawców interpretacja ta znosi konieczność stosowania pojęcia prototypu, na rzecz  bezpośredniego badania znaczenia wyrażeń (por. Wierzbicka 1989:~731--769; 1990:~347--367).

Trzecia interpretacja prototypu jest powiązana ze swoją poprzedniczką. Prototyp \linebreak jest według niej tożsamy ze znaczeniowym centrum i~przez to przeciwstawiony cechom peryferyjnym. W~myśl tej interpretacji jest prototyp znaczeniem ,,podstawowym'' (Grzegorczykowa 1998:~112).

\newpage
\chapter{Stereotyp wobec prototypu}

Socjologiczne pojęcie stereotypu, oznaczające ,,schematyczne, jednoznaczne, społecznie utrwalone wyobrażenia i~przekonania wiązane przez członków pewnej społeczności mówiącej z~określonymi zjawiskami, przedmiotami, a~przede wszystkim osobami'' (por. Bartmiński 1985:~25--53; Bartmiński~\&~Panasiuk 1993:~363--387) zostało przeniesione na obszar lingwistyki przez postacie takie jak Lidia Jordanskaja oraz Igor Mielczuk (1988:~9--34) podczas badań nad konotacjami leksemów semantycznych. Są one w~tym przypadku sądami i~cechami, jakie ludzie kojarzą z~danymi obiektami, zjawiskami oraz osobami, przy czym ci~ostatni --- w~kwestii narodowości --- są najbardziej narażeni na powstawanie stereotypów (Pisarkowa 1976:~5--26).

Stereotyp w~semantyce kognitywnej jest pojęciem równoległym w~stosunku do omówionego wcześniej prototypu. Podczas gdy ten drugi jest najlepszym przedstawicielem danej kategorii, stereotyp jest zespołem cech, które go określają, i~któremu odpowiadają (Putnam 1975:~249).

%\newpage
\section{Grzegorczykowej rozgraniczenie prototypu i~stereotypu}
\enlargethispage{-5\baselineskip}
Prototyp powinien pozostać punktem centralnym danej kategorii lub stanowić zbiór cech charakterystycznych owemu centrum (obiektom, ich klasom, sytuacjom) (Grzegorczykowa 1998:~113--114). Dla przykładu prototypowa \emph{matka} to ‘rodzicielka oraz wychowawczyni’, a~prototypowa \emph{łódka} to ‘taka, która może pływać’ (zatem \emph{wrak} nie jest uznawany za łódkę stricte) (por. Wierzbicka 1990:~347--367). Stereotyp pozostaje natomiast ustabilizowanym społecznie zbiorem konotacji, jakie wiążemy z~danymi zjawiskami lub pojęciami. Owe konotacje są utrwalone poprzez związki frazeologiczne, derywaty oraz połączenia wyrazowe (np. \emph{matkować komuś} lub \emph{być jak matka} w~przypadku konotacji związanych ze stereotypową troskliwą rodzicielką). Podczas gdy prototypowy pies to ‘zwierzę szczekające o~typowych cechach, które najlepiej oddają rysunki dzieci (cztery łapy, ogon zakręcony do góry)’, jego stereotyp będzie związany z~cechami, jakie przypisuje mu dana część społeczeństwa, np. wiernopoddańczość (por. Mosiołek-Kłosińska 1992: 301--304).

%\vspace*{0.5\baselineskip}
\section{Langackera rozgraniczenie prototypu i~schematu}

Według Langackera (1987:~371) prototyp jest ,,typowym przykładem kategorii, \linebreak inne zaś elementy włączane są do niego na podstawie podobieństwa, jakie postrzega się między nimi a~prototypem; istnieją stopnie przynależności oparte na stopniach podobieństwa''. Poza tym pojęciem wyróżnia on również pojęcie schematu jako abstrakcyjnej charakterystyki, która jest w~pełni zgodna ze wszystkimi elementami definiowanej kategorii (tamże), przez co nie istnieje w~schemacie pojęcie stopniowalności, za to stanowi ,,zintegrowaną strukturę ucieleśniającą wspólnotę swoich elementów --- pojęć bardziej szczegółowych, będących kontrastującymi ze sobą uszczegółowieniami schematu''.

Każdy z~langackerowskich schematów jest jednostką składającą się z~dwóch biegunów: semantycznego oraz fonologicznego, przez co jest konkretyzowany poprzez jednostkę, która w~danej sytuacji jest ,,mniej schematyczna''. Część znaczeń słów jest powiązana ze sobą relacjami uszczegóławiającymi, a~część --- relacjami poszerzania znaczeń (Kardela 2006:~180), a~wszystkie razem tworzą dokładny model sieciowy. Langacker (1999:~12) podkreśla, że istnieje ścisły związek pomiędzy kategoryzacją z~użyciem prototypu a~kategoryzacją poprzez schemat:
%\newpage
\end{spacing}
\begin{adjustwidth}{1,5cm}{1,5cm}
\small{
,,Akceptując ogólną zasadność modelu radialnego opartego na rozszerzeniu [...] muszę jednocześnie podkreślić rolę schematyzacji i~relacji realizacji [...], choćby dlatego, że odnoszą się one do absolutnie kluczowych zjawisk językowych: wyciągania uogólnień i~sądów dotyczących prawidłowej budowy wyrażeń (konwencjonalności). [...] Oba typy kategoryzacji są w~każdym razie niezwykle blisko ze sobą spokrewnione (i~w~praktyce niekiedy trudno je od siebie odróżnić). Oba wiążą się z~operacją porównania standardu [...] z~obiektem docelowym [...].''\\
}
\end{adjustwidth}
\begin{spacing}{1.3}

\newpage
\vspace*{4.5\baselineskip}
\begin{figure}[ht]
\begin{center}
{\large \sc
\begin{forest} for tree={align=center,edge=->,}
[\textbf{ring}:\\circular entity,l*=2
  [\textbf{ring}:\\circular mark,l*=2]
  [\textbf{ring}:\\circular object,name=2,l*=2
     [\textbf{ring}:\\circular piece\\of jewellery,l*=3
        [\textbf{ring}:\\circular piece\\of jewellery\\worn around finger,name=4,l*=2]
        [,phantom]
        [,phantom]
        [\textbf{ring}:\\circular piece\\of jewellery\\worn around finger,name=5,l*=2]
     ]
%     [,phantom]
%     [,phantom]
  ]
  [\textbf{ring}:\\circular arena,edge=dotted,name=3,l*=2]
]
\draw[->,dotted] (2) to[out=south,in=south] (3);
\draw[->,dotted] (4) to[out=south east,in=south] (5);
\draw[->,dotted] (5) to[out=south west,in=south] (4);
\end{forest}
}
\end{center}
\caption{Model sieciowy dla angielskiego rzeczownika \emph{ring}}% ‘pierścień', `ring (bokserski)’}
\end{figure}

\newpage

\vspace*{0.5\baselineskip}
\section{Zadeha teoria zbiorów rozmytych}

W~roku 1965 Lotfi Zadeh po raz pierwszy opisał teorię zbiorów rozmytych, której założeniem miało być formalizowanie wiedzy niepewnej lub aproksymatywnej, w~celu jej późniejszego przetworzenia przez systemy wykorzystujące sztuczną inteligencję. Przyjęcie pojęcia \emph{zbioru (podzbioru) rozmytego} pozwala na zastosowanie stopniowalnej przynależności elementu do danej klasy, a~sama logika rozmyta dopuszcza możliwość, że zdanie jest częściowo prawdziwe lub częściowo fałszywe (Golka 2009:~33). Kontrastowość tej teorii w~stosunku do logiki dwuwartościowej jest analogiczna do różnic między modelem warunków koniecznych i~wystarczających a~teorią kognitywnego prototypu.

Pierwsza próba formalizacji teorii prototypu miała miejsce poprzez połączenie jej z~teorią zbiorów rozmytych, gdyż obie te teorie ukazywały wyraźne ową nieostrość kategoryzacyjnych aspektów. Daniel Osherson i~Edward Smith (1981:~59--72) stwierdzili jednak, że unifikacja obu teorii prowadzi do nieprawidłowości na poziomie pojęć złożonych, m.in. w~kwestii pojęć koniunktywnych (składających się z~rzeczownika i~z~przymiotnika) oraz pojęć pustych i~uniwersalnych. Pierwszy przypadek dotyczy sytuacji, w~której przedmiot jest bardziej prototypowy dla pojęcia złożonego niż prostego, z~których to pojęć się składa. Mamy wtedy do czynienia z~paradoksem, ponieważ jego przynależność do pojęcia prostego będzie większa niż do pojęcia złożonego --- opisano to wyraźnie na przykładzie \emph{jabłka w~paski} (ang. \emph{stripped apple}), które intuicyjnie należy do kategorii ,,jabłko w~paski'', ale poprzez definicję iloczynów zbiorów bliżej jest mu do kategorii ,,jabłko'' (Golka 2009:~35). Drugi przypadek natomiast dotyczy pojęć pustych takich jak ,,jabłko, które nie jest jabłkiem''.

Unifikacja teorii prototypów i~teorii zbiorów rozmytych nie jest zatem tak prosta, jak semantyczna kompozycjonalność, w~myśl której pojęcia proste składają się na definicje pojęć złożonych, jednakże ,,trudności opisane przez Oshersona i~Smitha nie dyskredytują teorii zbiorów rozmytych jako bazy dla teorii prototypu. [...] Jakikolwiek opis pojęć o~budowie prototypowej jest niemożliwy bez odwołania się do zbiorów rozmytych'' (Zadeh 1982:~291--297). Twórca teorii zbiorów rozmytych zaznacza również, iż owe problemy wynikają z~błędnego rozumienia prototypu przez w/w badaczy.

\end{spacing}

\part{BOŻY SZALENIEC}
\clearpage{\pagestyle{empty}\cleardoublepage}
\chapter{Trickstera zarys ogólny}

\begin{spacing}{1.3}
Mianem \emph{trickstera} określa się w~religioznawstwie archetyp boskiego szaleńca, występującego naprzeciw odgórnym zasadom, autorytetom, normom i~porządkowi świata. Jego rola opiera się na przewrotnej zabawie rzeczywistością, przeskakiwaniu pomiędzy prawdą, fikcją, pragnieniami i~wszelkimi tematami tabu, oraz dbaniu o~cykliczność procesów kreacji i~destrukcji. Ta wieloznaczna postać obecna jest w~mitologiach wielu ludów, przy czym jej emanacją może być bóg (bogini), półbóg, śmiertelnik, duch lub niekiedy nawet rzecz (we współczesnej literaturze science fiction, np. Wirus --- por. Ziemiański 2012:~306), odgrywająca rolę kłamcy i/lub szaleńca, będącego często kluczową, pozytywną (lub przynajmniej ambiwalentną) postacią danej historii. Według Carla Gustava Junga (1954:~203) tricksterem może być:
\end{spacing}

\begin{adjustwidth}{1,5cm}{1,5cm}
\small{
,,[...] zarazem Bóg, człowiek, jak i~zwierzę. Jest zarówno podczłowiekiem jak i~nadczłowiekiem, bestią i~istotą boską, której główną i~najbardziej alarmującą cechą jest nieświadomość.''\\
}
\end{adjustwidth}

\begin{spacing}{1.3}
Troista natura trickstera znajduje również swoje potwierdzenie w~pracach innych badaczy, m.in. Lawrence'a Sullivana (1987:~45) oraz Beatrice Otto (2001:~40). Wypowiada się o~niej również Paul Radin (1956: 169) podkreślając ambiwalencję tej postaci:
\end{spacing}

\begin{adjustwidth}{1,5cm}{1,5cm}
\small{
,,[...] on pozostaje wszystkim w~każdym człowieku: bogiem, zwierzęciem, człowiekiem, bohaterem, bufonem, tym, który był, zanim zaistniały dobro i~zło, tym, który afirmuje i~neguje, niszczycielem i~twórcą.''\\
}
\end{adjustwidth}
\pagebreak
\vspace*{1\baselineskip}
\begin{spacing}{1.3}
Podług słów Roberta Peltona (1980:~6) i~Maca Linscotta Rickettsa (1987:~50) terminu \emph{trickster} użył po raz pierwszy w~roku 1868 Daniel Garrison Brinton w~swoim dziele \emph{The Myths of the New World: A~Treatise on the Symbolism and Mythology of the Red Race of America}, aczkolwiek data jego pierwszego wydania jako data narodzin tego określenia jest dyskusyjna (por. Sieradzan 2005:~565). Kolejną wzmiankę --- mówiącą o~,,elementach trickstera'' --- możemy znaleźć w~trzecim wydaniu dzieła Brintona z~1896 roku na stronie 194 (tamże).

Słowo \emph{trickster} jako \emph{terminus technicus} jest nieprzetłumaczalne, chociaż w~języku polskim --- poza angielskim określeniem --- spotyka się również pisownię `trikster' (por. Stokłosa 2012; Wasilewski). W~języku rosyjskim występuje czasami pod zrusycyzowaną postacią \emph{triukacz} (ros. {\otherlanguage{russian}{\large трюкач}}) (Wasilewski), chociaż zwykle używa się terminu oryginalnego (ros. \foreignlanguage{russian}{\large трикстер}). %трикстер

%(ros. \foreignlanguage{russian}{трюкач})

Ze względu na charakter tej postaci oraz jej skomplikowaną naturę jest \emph{trickster} określany jako ‘przechera’ lub ‘oszust’, które to nazwy mają ścisły związek z~jego aktywnością w~zakresie robienia sztuczek (ang. \emph{tricks}). W~zależności od sytuacji jest on również ‘zwodzicielem’, ‘deceptorem’ (fr. \emph{decepteur}), ‘boskim szelmą’ (niem. \emph{der göttliche Schelm}), ‘boskim urwisem’ (fr. \emph{le fripon divin}) lub ‘boskim oszustem’ (nid. \emph{de goddelijke bedrieger}). W~języku polskim można również spotkać się z~określeniami takimi jak ‘szachraj’, ‘szalbierz’, ‘spryciarz’, `prestidigitator', ‘macher’, ‘krętacz’ i~‘mistyfikator’, oraz z~innymi --- związanymi z~wszelkiego rodzaju oszustwami, magią i~szaleństwem. W~języku fińskim najbliższe tricksterowi określenia to m.in.: \emph{huijari} (‘aferzysta’, ‘oszust’, ‘szuler’), \emph{huiputtaja} (‘oszust’, ‘szalbierz’, ‘nabieracz’), \emph{petkuttaja} (‘oszust’, ‘szachraj’, ‘nabieracz’), \emph{temppuilija} (‘sztukmistrz’, ‘kawalarz’) oraz \emph{veijari} (‘szelma’, ‘krętacz’, ‘szachraj’, ‘szuja’, ‘łajdak’, ‘zbytnik’, ‘psotnik’) (Kudzinowski 1988), które to słowo oznacza również ‘sowizdrzała’ (por. \emph{veijariromaani} ‘powieść sowizdrzalska’) --- stare określenie prześmiewczego żartownisia o~niepokornej postawie wobec autorytetów i~oficjalnych norm społeczno-kulturowych.
\end{spacing}

\section{Idioty wymiar antropologiczno-szamański}
\begin{spacing}{1.3}
Głupota trickstera jest pierwszą cechą rzucającą się w~oczy czytelnikowi. Już pierwsi Indianie Winnebago opowiadali historie o~Wakdjunkadze, który zachowywał się tak, ,,jakby zerwał z~siebie uprząż rozumu'' (Wasilewski):
\end{spacing}
\pagebreak
\vspace*{0.5\baselineskip}
\begin{adjustwidth}{1,5cm}{1,5cm}
\small{
,,[...] nie potrafi zapanować nad uzbrojoną w~nóż ręką, która zaczyna walczyć z~drugą. Z~groteskową bezmyślnością --- a~może złośliwie --- uśmierca czwórkę dzieci, powierzonych jego opiece przez napotkanego myśliwego. W~panicznym strachu przed jego zemstą ucieka tak szybko, że błyskawicznie dociera na kraj świata, tam skacze do oceanu, pływa chaotycznie pytając ryb o~drogę i~żaląc się, że spotkało to właśnie jego, »wielkiego wojownika«, po czym okazuje się, że cały czas taplał się tuż przy brzegu. Jest głodny, a~że nie wychodzi mu złowienie ryby garnkiem, więc zadowala się ugotowaniem »zupy« na wodzie, przez którą ona przepływała --- zjada ją zadowolony, mlaszcząc ze smakiem. Rozmawia z~pniakiem, biorąc go za człowieka.''\\
}
\end{adjustwidth}

\begin{spacing}{1.3}
Nieokrzesany i~niedojrzały, pada trickster również ofiarą swoich własnych figli (por. Leeming 1990, za: Stokłosa 2012:~121). Z~drugiej strony, występuje w~kulturach łowieckich często pod postaciami sprytnych zwierząt (m.in. Kojota i~Kruka, które razem z~Zającem i~Pająkiem określi później Wasilewski mianem tricksterskiej \emph{Wielkiej Czwórki}) --- co daje podstawy Radinowi (por. 2010) do nazwania go ,,mądrym głupcem''. Uważa się, iż trickster jest wytworem całej ludzkości oraz stanowi odpowiedź na pytania zadawane od zarania dziejów, jak również swoiste zwierciadło wewnętrznych i~zewnętrznych problemów człowieka (Radin 2010:~22). To właśnie nieuznający żadnych wartości wyższych trickster sprawia, iż mogą one zaistnieć.

Kulturowa rola trickstera jest o~tyle ważna, iż pomimo tego, że małpuje on boga (Jung 1954: 195--204) (łac. \emph{simia dei}), to jest mu bardzo blisko do jednoczesnego bycia zbawcą. Greckiego Prometeusza współbogowie uważali za szaleńca i~przestępcę, ponieważ wykradł on ich święty ogień i~przekazał go maluczkim ludziom, którzy z~kolei obwołali go bohaterem i~dobroczyńcą (a~następnie świętym męczennikiem). W~tej postaci kozioł ofiarny współistnieje ze zbawicielem (por. Sieradzan 2005:~97). Dodatkowo jest trickster również uniwersalnym i~transpersonalnym ucieleśnieniem ,,figlarności wyobraźni'' (por. Fairclough 2005:~17), który stworzył świat wbrew woli głównego bóstwa (lub niezgodnie z~jego poleceniami), tudzież przeszkodził tejże doskonałej kreacji.

W~postaci tricksterów widzi świat magików, zaklinaczy, zdegradowanych bogów i~bohaterów kulturowych, określając ich czasami ,,świętymi grzesznikami'' (por. Sieradzan 2005:~95--97), co można odnieść do magii (a~dalej --- do szamanizmu), tępionej w~średniowieczu przez różne kościoły. Znamienne jest również, iż greckie \textgreek{γόης} [g\={o}es] --- synonim \emph{trickstera} --- oznacza zarówno ‘oszusta’, jak i~‘szamana’ (chociaż tłumaczy się to słowo również jako ‘trickstera’ stricte --- por. Bolton 1962:~35). W~tym kontekście przytacza go Radin (2010: 119) na przykładzie \emph{Cyklu Zająca} Winnebagów:
\end{spacing}
\pagebreak
\begin{adjustwidth}{1,5cm}{1,5cm}
\small{
„Płakał nad ludźmi. Było to w czasie, gdy rozmyślał nad stworzeniem Rytuału szamańskiego. A oto historia powstania Rytuału szamańskiego. Wszystkie ptaki latające, które były złe albo dręczyły ludzi, Zając pozabijał. Dlatego powiada się, że wypchnął je wyżej w powietrze. Wszystkie złe duchy, które mieszkały na tej ziemi, też pozabijał, żeby już dłużej się po niej nie pałętały. O to właśnie chodzi, kiedy się mówi, że głębiej wdeptał te, których grzbiety wystawały z ziemi.”\\
}
\end{adjustwidth}

\section{Groteskowa seksualność}
\begin{spacing}{1.3}

Trickstera cechuje dodatkowo nadzwyczajny apetyt, wielce spotęgowana cielesność oraz pozbawiona jakichkolwiek zahamowań seksualność (por. Stokłosa 2012:~121--122), ocierająca się bardzo często o~groteskę. Boskim szaleńcem, będącym orędownikiem i~personifikacją cielesnej strony życia (por. Kerényi 2010:~210), włada chuć wespół z~infantylizmem i~głupotą. Nie potrafi on pojąć norm społecznych i~etycznych, jest amoralny i~nieodpowiedzialny (por. Leeming, za: Stokłosa 2012:~121--122), a~przez swoje twórcze kłamstwa zaspokaja on zarówno swoje pożądanie seksualne, jak i~głód. W~mitach tricksterskich mamy do czynienia z~przełamywanym seksualnego tabu, co ciekawie przytacza Wasilewski:
\end{spacing}
\begin{adjustwidth}{1,5cm}{1,5cm}
\small{
,,Dawno, dawno temu w~pewnej indiańskiej osadzie mieszkał wódz, który właśnie miał wyruszyć na wojenną ścieżkę. Zwołał wszystkich wojowników na ucztę, w~trakcie której gdzieś nagle zniknął. Jakież było zaskoczenie szukających, gdy znaleźli go spółkującego z~kobietą. Rozeszli się, przerażeni, że oto wódz łamie najsurowsze tabu, jakie panuje w~okresie zbrojnej wyprawy.''\\
}
\end{adjustwidth}

\begin{spacing}{1.3}
Postać trickstera rozwija się wraz z~rozwojem samej historii. Nabywa on świadomość własnego ciała poprzez groteskowe przeobrażenia cielesne, oraz uczy się seksualności, folgując jednocześnie swojej głupocie w~pełnym tego słowa znaczeniu (Wasilewski):
\end{spacing}

\begin{adjustwidth}{1,5cm}{1,5cm}
\small{
,,Wakdjunkaga kładzie się spać, a~obudziwszy się widzi wysoko nad sobą coś, co bierze za dumnie powiewający wodzowski sztandar. Dopiero po chwili orientuje się, że jest to derka, którą był przykryty, uniesiona przez jego własny członek w~erekcji. Prosi go, by zmiękł, po czym [...] chowa do pudła, które zarzuca sobie na plecy. Dostrzega kobiety wodza, kąpiące się na drugim brzegu jeziora, wydobywa członek z~pudła i~posyła go w~kierunku córki wodza. Za pierwszym razem członek płynie po powierzchni, ale to spłoszyłoby kobiety, więc Trikster przywiązuje doń duży kamień, a~członek porusza się po dnie. To z~kolei za głęboko, bohater znajduje więc mniejszy kamień, ale takie rozwiązanie wywołuje fale na powierzchni. Dopiero za czwartym razem udaje się Triksterowi posłać penis wprost ku córce wodza i~to tak skutecznie, że przerażone dziewczęta nie mogą go wyciągnąć --- [...] a~Trikster zaśmiewa się z~kobiet, że niepotrzebnie psują mu przyjemność.''\\
}
\end{adjustwidth}

\chapter{Boski szaleniec ponad podziałami}
\begin{spacing}{1.3}
Trickster jest postacią zarówno konstruującą, jak i~znoszącą wszelkiego rodzaju granice, zarówno w~aspekcie fizycznym, jak i~społecznym oraz religijnym. To on --- jako boski posłaniec --- przenosi bogów na niebiosa, pozostając jednocześnie na ziemi, by tworzyć świat, dobra i~kulturę. Bywa określany również jako bóg, który ,,chętnie płata figle innym bogom, a czasami także ludziom i~zwierzętom'' (Stokłosa 2012: \linebreak 121--122), oraz jako tzw. \emph{bóg progu}, gdyż jest ,,wszędzie tam, gdzie przebiega granica --- po to, by ją przekroczyć, zlikwidować, przenieść czy też po prostu wyznaczyć'' (Hyde 2008:~8). Granice te dotyczą kategorii dobra i~zła, świętości i~świeckości, problematyki czystości, płci, wieku oraz życia i~śmierci (Hyde 1998:~7). Ze względu na ogromny zakres swojej specyficznej działalności jest trickster według Radina ,,reprezentantem wszelkich możliwości'' (za: Blaeser 1996:~139), a~Jarold Ramsey określa go po prostu jako tego, który ,,nie jest ani tym, ani tamtym, lecz i~tym, i~tamtym'' (tamże), co wskazuje na jego wszechmoc oraz ambiwalentny charakter, w~którym nieoznaczoność współistnieje wespół z~potencjałem twórczym. Według Williama Doty’ego i~Williama Hynesa (1993: 34--42) każdy trickster posiada kilka z~sześciu charakterystycznych cech:

\begin{enumerate}[1)]
  \item jest zasadniczo niejednoznaczną i~nietypową postacią,
  \item jest shapeshifterem (zmiennokształtnym),
  \item jest posłańcem i~naśladowcą bogów,
  \item odwraca/przekręca daną sytuację,
  \item łączy świętość ze sprośnością,
  \item jest oszustem i~zwodzicielem.
\end{enumerate}
\end{spacing}

\pagebreak
\vspace*{1\baselineskip}
\section{\`{E}\d{s}\`{u}-Elegbara}
\begin{spacing}{1.3}
Jorubańskie bóstwo \`{E}\d{s}\`{u}-Elegbara jest przykładem boga-trickstera działającego ponad granicami, będącego jednocześnie posłańcem losu i~dostarczycielem boskich informacji, zarówno dla innych bogów, jak i~dla ludzi. Do jego cech typowych zalicza się kapryśność (cecha typowo tricksterska, gdyż postać ta jest zdana na łaskę swoich namiętności, nie znając wartości moralnych), zwinność oraz szybkość ciała i~umysłu, która to cecha uczyniła z~niego również językoznawcę i~inteligentnego metafizyka (por. Davis 2000: 48--49).
\end{spacing}

\section{Hermes}
\begin{spacing}{1.3}
Trafnym przykładem trickstera jako istoty ponad granicami, podobnej \`{E}\d{s}\`{u}-Elegba\-rze, jest postać greckiego Hermesa. Ten nieprzypisany jednoznacznie ani do Olimpu, ani do świata ziemskiego posłaniec zwycięża poprzez swój pacyfizm oraz poprawne stosunki z~wszelkimi istotami --- poczynając od mieszkańców sfer niebieskich, na duszach w~Hadesie kończąc. Znamienne jest również to, iż jego syn --- Hermafrodyt --- posiadał zarówno męskie, jak i~żeńskie narządy płciowe (Mendecki 2000: 135). Tricksteryzm Hermesa w~praktyce jest wyraźnie widoczny, chociaż Kerényi (2010: 124--125) krytykuje takie podejście za to, iż podobnego jemu Wakdjunkagi nigdy nie uznano za bóstwo równe olimpijskim, nawet według standardów indiańskich.
\end{spacing}

\vspace*{1\baselineskip}
\chapter{Troisty głupiec prapoczątku}
\begin{spacing}{1.3}
Postać boskiego głupca jest powiązana z~prapoczątkiem świata, a~odkrywane na całym globie mity tricksterskie należą do najstarszych wytworów ludzkości (Radin 2010: 21). Kształtujący się, niedoskonały i~chaotyczny świat czasów Genesis jest widoczny w~ogólnej postawie trickstera, który jest postacią nierozwiniętą, bez stałej formy, o~nieokreślonych proporcjach (Radin 2010: 22). Razem z~rozwojem człowieka wykształconego i~pojawieniem się pojęcia dobra i~zła ulega pierwotny trickster przedefiniowaniu, jak również stanowi on ideologiczną \emph{axis mundi} ludzkiej cywilizacji (por. Diamond 2010: 9--11). Świat został stworzony ,,głupio'' i~niedoskonale właśnie za sprawą trickstera, który paradoksalnie umożliwił mu prawidłowe zaistnienie w~obecnej, akceptowalnej postaci (Wasilewski).
\end{spacing}

\pagebreak
\vspace*{1\baselineskip}
\section{Loki}
\begin{spacing}{1.3}

Trickster to prawdopodobnie najstarsza postać we wszystkich ziemskich mitologiach (por. Radin 1956: 164). Dobrym przykładem owego pradawnego szalbierza jest nordycki jötunn Loki, olbrzym zaliczony w~poczet bogów, do którego można odnieść słowa Peltona (1980: 269): ,,Trickster nie jest ani Bogiem, ani człowiekiem, ani zwierzęciem; on jest nimi wszystkimi''. Będący symbolem ognia i~oszustwa Loki wywodzi się od olbrzyma imieniem Fárbauti, a~jego potomstwo jest pod wszelkimi względami mocno zróżnicowane: z~olbrzymką Angrboðą ma on troje dzieci, z~czego jedno jest wilkiem (Fenrir), jedno wężem (Jörmungandr), a~jedno boginią umarłych (Hel). Ze swoją żoną-boginią Sygin ma on z~kolei synów Váliego i~Narfiego (Nariego), co nie przeszkodziło mu urodzić jeszcze po drodze pod postacią klaczy ośmionogiego konia Sleipnira (Welch 2011: 47), jak i~szeregu pomniejszych potworów pod postacią kobiety. Świadczy to o~złożoności postaci Lokiego, który stając się ojcem ludzi, zwierząt i~bóstw, jest istotą ponad wszelkimi granicami, zarówno społecznymi, jak i~nawet biologicznymi.

Chociaż widoczne jest w~tej postaci podobieństwo do Hermesa --- wszakże przez całe życie Loki nie trzymał z~żadną ze stron (bóstwami, olbrzymami bądź ludźmi) --- to w~przeciwieństwie do greckiego posłańca nie charakteryzuje się on ani pacyfizmem, ani tym bardziej nawet poprawnymi stosunkami w~kwestiach dyplomatycznych. Pomimo tego, iż łączą go z~bogami dobre stosunki (dzięki niemu zdobyli oni wszakże swoje artefakty), to nie przeszkadza mu to w~spóźnieniu się na ich wielką biesiadę i~--- miast uszanowania gospodarzy --- wytykania im ich niechlubnych uczynków z~przeszłości (por. \emph{Kłótnia Lokiego}, Edda Poetycka 1986: 108--124). Tricksterskie szaleństwo Lokiego nie ma praktycznie żadnych granic --- poczynając od ścięcia pięknych włosów bogini Sif, a~kończąc na morderstwie boga Baldura, za co zostaje przez jego kompanów w~iście prometejski sposób przywiązany do skały, na której cierpi katusze aż do dnia Ragnaröku. Wtedy to ma on się oswobodzić, popłynąć na swoim statku do Asgårdu by --- ostatecznie opowiadając się po stronie olbrzymów --- poprowadzić swoją armię na wojnę, w~której spłonie cały świat.
\end{spacing}

\chapter{Ambiwalentny kryminalista imieniem Jahwe}
\begin{spacing}{1.3}
Pojęcia ,,kryminalnego bóstwa'' użył w~odniesieniu do religii hinduistycznej historyk religii Alf Hiltebeitel (1989:~1), chociaż ambiwalencję bogów wskazał już siedem lat wcześniej Bede Griffiths (1982:~94) w~stosunku do Śiwy, Kali i~Jahwe. Hiltebeitel określił tym mianem boga łamiącego ,,święte przepisy i~granice, którymi kierują się w~życiu inni bogowie i~ludzie'', przez co staje się on podobny demonom. Bóstwem kryminalnym jest według tej koncepcji Jahwe, którego tricksterską naturę dostrzega się m.in. na przykładzie ambiwalentnego podejścia do ogólnego pojęcia dobra i~zła w~Biblii, jak również personifikacji tychże pojęć w~osobie samego bóstwa (Diamond 2010: 15). Jung (1958b: 709--710) określa boga chrześcijańskiego jako niemoralną istotę, na której --- pomimo faktu, iż jest ,,strażnikiem prawa i~porządku'' --- nie można polegać, gdyż jest ona mściwa, wybuchowa i~niesprawiedliwa. Podobne zdanie ma również Sieradzan (2005: 196), według którego jest bóg chrześcijański paradoksalnie ,,przyczyną istniejącego w~świecie chaosu'':
\end{spacing}

\begin{adjustwidth}{1,5cm}{1,5cm}
\small{
,,Jahwe [...] pierwotnie zawierał [...] w~sobie również aspekt demoniczny i~często nakazywał łamanie praw, które sam ustanowił, jeśli korzystał na tym albo on sam, albo jego »lud wybrany«.''\\
}
\end{adjustwidth}

\begin{spacing}{1.3}
Dualizm tej postaci objawia się również w~manipulowaniu ludźmi oraz inspirowaniu ich do czynów zarówno pozytywnych, jak i~tych o~wiele mniej chwalebnych. Ma to z~kolei związek z~nakazywaną przez niego podwójną moralnością, według której lud izraelicki może kierować się etycznymi zasadami dziesięciorga przykazań, a~w~stosunku do ludów barbarzyńskich stosować iście odmienną jurysdykcję --- idealnym przykładem obydwu tych mechanizmów jest posłużenie się zmanipulowanym Mojżeszem w~celu zesłania niszczycielskich plag na Egipt i~wymordowania pierworodnych.

Uznanie Jahwe za bóstwo kryminalne jest ściśle związane z~etyczną ambiwalencją początków świata (w~tym zaś przypadku okresu przed powstaniem ksiąg profetycznych), kiedy boska natura cechowała się amoralnością, zawierając w~sobie pierwiastek dobra i~zła (por. Diamond 2010: 19--20). Tricksterski dualizm boga Izraelitów jest widoczny w~Starym Testamencie (Pwt~32:~39; Iz~45:~7; Lm~3:~38; Am~3:~6):
\end{spacing}

\begin{adjustwidth}{1,5cm}{1,5cm}
\small{
,,[...] Ja zabijam i~Ja sam ożywiam, Ja ranię i~Ja sam uzdrawiam, że nikt z~mojej ręki nie uwalnia. [...]
Ja tworzę światło i~stwarzam ciemności, sprawiam pomyślność i~stwarzam niedolę. Ja, Pan, czynię to wszystko. [...]
Czyż nie pochodzi z~ust Najwyższego i~niedola, i~szczęście? [...]
Czyż zdarza się w~mieście nieszczęście, by Pan tego nie sprawił?''\\
}
\end{adjustwidth}

\begin{spacing}{1.3}
%\noindent
Stwórca świata jest zatem demonicznym, plemiennym szaleńcem, na co zwrócił uwagę Jung (1954:~196), pisząc o~jego ,,bezsensownych orgiach zniszczenia''. Przykładem jego prymitywizmu jest np. utrata zaufania do Mojżesza, która automatycznie wiązała się z~próbą zabicia go (por. Wj~4:~24). Liczne przykłady destrukcyjnej natury boga-trickstera zostały przytoczone w~biblijnej Księdze Ezechiela (Ez~21:~33--37):
\end{spacing}

\begin{adjustwidth}{1,5cm}{1,5cm}
\small{
,,W~miejscu, gdzie zostałeś stworzony, i~w~kraju, gdzie się zrodziłeś, będę cię sądził. Wyleję na ciebie mój gniew, rozniecę przeciwko tobie ogień mojej zapalczywości i~wydam cię w~ręce ludzi dzikich, sprawców zniszczenia. Staniesz się strawą dla ognia, krew twoja będzie płynąć w~środku kraju i~nie pozostawisz po sobie wspomnienia --- ponieważ Ja, Pan, tak powiedziałem.''\\
}
\end{adjustwidth}

\begin{spacing}{1.3}
Jahwe jest również tajemniczym graczem, który ukrywa się przed ludźmi (łac. \emph{deus absconditus}), przez co wystawia ich na różnego rodzaju próby, np. w~drodze do obiecanej ziemi Kanaanu. Jest jednocześnie postacią ,,z~koszmarnego snu, tragedii greckiej (gdzie każde wyjście jest złe)'', można go również interpretować jako ,,wytwór kapryśnej nieświadomości albo freudowskiego {\sc ID}'' (Sieradzan 2005:~197). Jednakże chyba najciekawszym przykładem jego ambiwalencji, tricksteryzmu i~szaleństwa jest fakt, iż --- za Sieradzanem (tamże) --- ,,»duch Jahwe« jest traktowany jako byt odrębny od samego Boga: Jahwe prowadzi z~nim dialog jak z~drugim człowiekiem'' (por. 2~Krn~18: 19--23). W~okresie późniejszym Biblia przeszła swego rodzaju cenzurę, podczas której postać niegodziwego Jahwe podzielono na dwa byty: dobry bóg stał się ,,Panem'', podczas gdy jego ,,zła, mroczna strona, która stymulowała jego zachowania trickstera, stała się »Szatanem«'' (Sieradzan 2005:~197).

Znamiennym przykładem i~dowodem na tricksteryzm boga Izraelitów jest oczywista analogia do pierwotnego trickstera północnoamerykańskiego --- Wakdjunkagi --- o~którym opowieści samym Indianom Winnebago sprawiały kłopoty interpretacyjne, jak również ,,mieli [oni] nawet zasadniczą trudność ze zdecydowaniem, czy była to postać dobra czy zła. Wyjaśniając to białemu badaczowi, jedni porównywali go do chrześcijańskiego Diabła, podczas gdy inni do Jezusa'' (Wasilewski).

\end{spacing}


\addtocontents{toc}{\protect\vspace*{60pt}}
\addtocontents{toc}{\protect\vspace*{3.70\baselineskip}}
\clearpage{\pagestyle{empty}\cleardoublepage}
\part{SYNTEZA \& ANALIZA}
\clearpage{\pagestyle{empty}\cleardoublepage}

\chapter{Metoda na szalbierzy w eposie}% w eposie}
\begin{spacing}{1.3}
Niniejsza część pracy ukazuje tricksteryzm postaci występujących w~fińskim eposie narodowym --- \emph{Kalevali}. Ukazanie wzorowego przykładu kalevalicznego przechery następuje poprzez analizę bohaterów eposu pod względem domen kognitywnych przypisanych do ogólnego pojęcia ,,trickstera'', nałożonych następnie na bazę kognitywną, którą w~tym przypadku jest treść eposu. Owe domeny tricksterskie zostały stworzone na podstawie opisanych wcześniej cech kłamców, błaznów i~szaleńców w~myśl zasady tricksterskich punktów kluczowych Doty’ego i~Hynesa, odpowiednio rozszerzonych. Należą do nich:\\


\begin{adjustwidth}{\parindent}{0cm}
\begin{labeling}{deus ex machinaaaa}

\item [\textbf{więź z bóstwem}] naśladownictwo, posłannictwo
\item [\textbf{ambiwalencja}] egzystencja na granicy lub ponad granicami
\item [\textbf{shapeshifting}] zmiennokształtność
\item [\textbf{magia rytualna}] śpiewy, zaklinania
\item [\textbf{sacrofanum}] seksualność połączona ze świętością
\item [\textbf{deus ex machina}] nagła zmiana sytuacji i/lub osiągnięcie danego celu
\item [\textbf{kłamstwo}] oszustwa i zwodzenia
\item [\textbf{idiotyzm}] głupota oraz jej nieświadomość\\

\end{labeling}
\end{adjustwidth}


%\begin{enumerate}
%  \item {\sc więź z bóstwem} --- naśladownictwo, posłannictwo,
%  \item {\sc ambiwalencja} --- egzystencja na granicy lub ponad granicami,
%  \item {\sc shapeshifting} --- zmiennokształtność,
%  \item {\sc magia + zaklęcia} --- rytuały, śpiewy, zaklinania,
%  \item {\sc sacrofanum} --- seksualność i~łączenie jej ze świętością,
%  \item {\sc deus ex machina} --- nagła zmiana sytuacji i/lub osiągnięcie danego celu,
%  \item {\sc kłamstwo} --- oszustwa i zwodzenia,
%  \item {\sc idiotyzm} --- głupota oraz jej nieświadomość.
%\end{enumerate}
%\noindent

Postacie, do których przypisana zostanie większa liczba powyższych domen będą stanowiły lepsze (bardziej wzorcowe) odbicie kategorii trickstera na płaszczyźnie kalevalicznej, podczas gdy reszta z~bohaterów będzie jego ,,słabszymi'' emanacjami.
\end{spacing}

\chapter{Tricksterzy kalevaliczni}
\begin{spacing}{1.3}
Bohaterowie \emph{Kalevali} noszą w~sobie cały wachlarz cech klasyfikujących ich jako tricksterów. Czytając zredagowane przez Eliasa Lönnrota dzieło napotykamy na zaklinaczy, szamanów i~czarownice parających się szeroko pojętą magią, którzy stosunkowo często wpadają w~dość groteskowe perypetie. Na kartach eposu wiedza tajemna miesza się nierzadko z~czystą głupotą oraz brakiem rozsądku, przeplatanymi nagłymi rozwiązaniami akcji w myśl greckiego \emph{deus ex machina}.

\section{Väinämöinen}
\vspace*{0.5\baselineskip}
\subsection*{Więź z bóstwem}

Väinämöinen jest dostojnym, wiekowym zaklinaczem, wymienionym już na samym początku eposu. Głównego bohatera \emph{Kalevali} wiele łączy z~bóstwami, gdyż jego matką jest bogini Ilmatar, władczyni powietrza (I:~107--110):
\end{spacing}
{\centering
\small
\begin{longtable}{c c}
[...]                         & [...] \\
Sam się zrodził Väinämöinen,  & yksin syntyi Väinämöinen, \\
Piewca odwieczny pojawił      & ilmestyi ikirunoja, \\
Z panieńskiej prarodzicielki, & kapehesta kantajasta, \\
Dziewy powietrza Ilmatar.     & Ilmattaresta emosta. \\
\end{longtable}
}
%\vspace*{\baselineskip}
\begin{spacing}{1.3}
W~kwestii mówienia o~sobie samym jest Väinämöinen podobny do Jahwe --- podczas gdy bóg Izraelitów przemawia w~2~Księdze Kronik do swojego ducha jako osobnego bytu, tak kalevaliczny zaklinacz nie zawsze wyraża się o~sobie bezpośrednio --- ale tak, jakby miał na myśli osobę trzecią (XL:~223--226):
\end{spacing}

\newpage
{\centering
\small
\begin{longtable}{c c}
,,Dałoby się z rybiej kości  & ,,Näistäpä toki tulisi \\
Zrobić muzyckie narzędzie,   & kalanluinen kanteloinen \\
Jeśliby się ktoś pokusił     & kun oisi osoajata, \\
I do grania je przymusił.''  & soiton luisen laatijata.'' \\
\end{longtable}
}

\begin{spacing}{1.3}
%\noindent
Jest to dość osobliwe, zwłaszcza że wszyscy doskonale zdają sobie sprawę, kogo ma on w~danej chwili na myśli. Kiedy indziej zaś mówi o~sobie jako o~tym, który dalej przekazuje bożego ducha, co już zakrawa na tricksterską bezczelność (XLV:~329--336):
\end{spacing}
{\centering
\small
\begin{longtable}{c c}
Nic nie sprawiłbym od siebie  & ,,En minä mitänä voine, \\
Bez pomocy Stwórcy w niebie,  & kun ei Luojani luvanne. \\
Użycz, Stwórco, swej pomocy,  & Avun Luoja antakohon, \\
Wspomożenia użycz, Boże,      & avun tuokohon Jumala \\
Tam, gdzie oczy me zobaczą,   & minun silmin nähtyäni, \\
Moje ręce namacają,           & käsin päällä käytyäni, \\
Moje usta wypowiedzą,         & suin sulin puheltuani, \\
Tam, gdzie tchnę swojego ducha! & hengin henkäeltyäni! \\
\end{longtable}
}
%\vspace{\baselineskip}
\begin{spacing}{1.3}
%\noindent
Warto zwrócić uwagę na fakt, iż na samym początku wypowiada się Väinämöinen dość pokornie, lecz zaraz jego pokora zmienia się w~pychę i~poczucie swoistego awansu z~pozycji ziemskiego znachora do pozycji zbawcy-uzdrowiciela, w~którego działaniach dopiero pojawić się ma boskie tchnienie. Väinämöinen pragnie wielkiej chwały (i~zbytnio się z~tym nie kryje), tak samo jak ma w~zwyczaju bezwstydnie wysługiwać się bóstwem w~wielu sprawach (posiłkując się dodatkowo fałszywą skromnością), co mu się dość często udaje (XVIII:~29--40):
\end{spacing}
{\centering
\small
\begin{longtable}{c c}
,,Znijdź teraz do czółna, Boże, & ,,Tule nyt purtehen, Jumala, \\
Do mej łodzi, miłosierny, & aluksehen, armollinen, \\
Przydaj mocy mnie, słabemu, & väeksi vähän urohon, \\
Męstwa człekowi skromnemu, & miehen pienen miehuueksi \\
Na tych wodach niezmierzonych, & noilla väljillä vesillä, \\
Rozfalowanych przestworzach. & lake'illa lainehilla! \\
Wietrze, kołysz moje czółno, & ,,Tuuittele, tuuli, purtta, \\
Falo, nieś na sobie statek, & aalto, laivoa ajele \\
Bym nie wiosłował rękoma, & ilman sormin soutamatta, \\
Nie mącił powierzchni wody & ve’en kieron rikkomatta, \\
Na rozległej głębi morza, & väljille meren selille, \\
Na otwartej jego toni!'' & ulapalle aukealle!'' \\
\end{longtable}
}
%\vspace{\baselineskip}
\begin{spacing}{1.3}
Widać to również bardzo wyraźnie, gdy wyrusza Väinämöinen jako znachor na wojnę z~chorobami bogini bólu Loviatar --- po przygotowaniu sauny prosi on Ukko o~tchnięcie boskiej mocy w~parę, jak i~o~usunięcie pary złej (wszakże bogini podobnie w~saunie powiła swoje potomstwo). Widać wyraźnie szelmowską pewność siebie Väinämöinena co do pomocy ze strony boga (XLV:~229--236):
\end{spacing}
\newpage
{\centering
\small
\begin{longtable}{c c}
Nikt nie poźre nas niewinnych, & ,,Ei nyt meitä syyttä syöä, \\
Nie ruszy nas, uzdrowionych,& eikä tauitta tapeta, \\
Skoro Stwórca nie zezwoli,& ei luvatta suuren Luojan, \\
Nie umorzy Bóg na niebie,& ilman surmatta Jumalan. \\
Kto by pożarł nas niewinnych,& Kenpä meitä syyttä söisi, \\
Ten by zjadł swe własne słowa,& suuhunsa omat sanansa, \\
Zło przywołał na swą głowę,& päähänsä pahat panonsa, \\
Zamysł na samego siebie!&ajatukset itsehensä! \\
\end{longtable}
}
%\vspace{\baselineskip}
\begin{spacing}{1.3}
Jednakże nie zawsze kalevaliczny czarownik jest herosem o~zbyt wybujałym ego --- w~wielu sytuacjach podkreśla on rolę Ukko jako najwyższego dobroczyńcy, którego zalecenia wykonuje na chwałę zarówno jego, jak i~ludzi, którym pomaga w~duchu prometeizmu (II:~294--308):
\end{spacing}
{\centering
\small
\begin{longtable}{c c}
[...] & [...]\\
Rozsypywać jął nasiona & siementä sirottamahan. \\
Wypowiadając te słowa: & Itse tuon sanoiksi virkki: \\
,,Rzucam ziarna, przesiewając & ,,Minä kylvän kyyhättelen \\
Pomiędzy palcami Stwórcy, & Luojan sormien lomitse, \\
Z prawicy Wszechmogącego, & käen kautta kaikkivallan \\
By wyrosły na tej roli, & tälle maalle kasvavalle, \\
Wybujały na polanie. & ahollen ylenevälle. \\
Starko żyjąca pod ziemią, & ,,Akka manteren-alainen, \\
Babko pól, włodarko roli! & mannun eukko, maan emäntä! \\
Daj się wykluć z ziemi ziołom, & Pane nyt turve tunkemahan, \\
I roli, by je wspomogła; & maa väkevä vääntämähän! \\
Niechaj nigdy na tym świecie & Eip' on maa väkeä puutu \\
Mocy w roli nie ubywa, & sinä ilmoisna ikänä, \\
Póki łask nie skąpią dawcy, & kun lie armo antajista, \\
Córy stworzenia obietnic. & lupa luonnon tyttäristä. \\
\end{longtable}
}
%\vspace{\baselineskip}

\begin{spacing}{1.3}
Wiele jest w~\emph{Kalevali} fragmentów, w~których Väinämöinen urasta do rangi zbawcy lub co najmniej Prometeusza --- obsiewa on krynicę Kalevy nasionami, a~jego pieśni sprawiają światu radość, zwłaszcza wtedy, gdy gra na swoim kantele. Mówi o~nim wszakże matka Joukahainena do swojego syna (VI:~117--124):
\end{spacing}
{\centering
\small
\begin{longtable}{c c}
Jeślibyś zastrzelił Väinö, & ,,Ampuisitko Väinämöisen, \\
Co z Kalevy się wywodzi, & kaataisit kalevalaisen, \\
Radość uszłaby ze świata, & ilo ilmalta katoisi, \\
Pieśń zniknęłaby z tej ziemi; & laulu maalta lankeaisi. \\
{[...]} & {[...]}\\
\end{longtable}
}
%\vspace{\baselineskip}
\begin{spacing}{1.3}
%\noindent
Väinämöinen ukazuje się jako heros i~dobroczyńca również wtedy, gdy wyrusza razem ze swoimi kompanami do Pohjoli po Sampo --- magiczny młynek, który ma zmienić jego ojczyznę w~krainę mlekiem i~miodem płynącą.
\end{spacing}
\vspace*{0.5\baselineskip}
\subsection*{Ambiwalencja}
\begin{spacing}{1.3}
W~kwestii graniczności należy w~przypadku Väinämöinena wymienić dwa istotne aspekty: jest on postacią łączącą w~sobie mądrość i~głupotę, jak również jest on tym, który może egzystować na granicy światów. W~pierwszym przypadku widoczne jest to w~sytuacjach groteskowych, takich jak grabienie morza w~poszukiwaniu utraconego instrumentu (XLIV:~55--66):
\end{spacing}
{\centering
\small
\begin{longtable}{c c}
Stary mądry Väinämöinen & Vaka vanha Väinämöinen \\
U steru się usadowił, & itse istuihe perähän; \\
Głębie morza jął przegarniać, & läksi merta luutimahan, \\
Fale grabiami zamiatać; & lainetta lakaisemahan. \\
Wszystkie grążele zgromadził, & Luopi lumpehet kokohon, \\
Wygrabił na brzegu śmiecie, & haravoipi rannan raiskat, \\
{[...]} & {[...]} \\
Lecz nie znalazł, nie odzyskał & eipä saanut, ei tavannut \\
Narzędzia ze szczęk szczupaczych, & hauinluista soittoansa, \\
{[...]} & {[...]} \\
\end{longtable}
}
\begin{spacing}{1.3}
%\noindent
Drugi przypadek natomiast to ,,stypa po niedźwiedziu'' --- rytualny obrzęd odbywający się po polowaniu, wspominający życie zwierzęcia, opisujący jego śmierć, oraz mający pomóc mu w~przejściu do świata ,,na wysokościach'', a~następnie w~zmartwychwstaniu (por. Haavio 1979:~21--47). W~fińskim eposie obrzędowi przewodzi Väinämöinen, pełniący rolę greckiego Charona jako przewoźnika dusz. Mediator światów kłamie m.in. na temat niezdejmowania z~niedźwiedzia futra, jednocześnie to czyniąc --- a~mając na celu jego ochronę przed profanacją. Tricksterskie zamierzenie, zrealizowane poprzez kłamstwo i~(pozorny) brak logiki zostaje osiągnięte (XLVI:~295--304):
\end{spacing}
{\centering
\small
\begin{longtable}{c c}
[...] & [...]\\
Nie zepsujemy ci futra, & Ei tuhota turkkiasi, \\
Sierści twej nie oglądamy, & karvojasi ei katsota \\
By ją wdział na siebie nędznik, & herjojen hetalehiksi, \\
Sponiewierał ją kaleka.'' & vaivaisien vaattehiksi.'' \\
Po czym stary Väinämöinen & Siitä vanha Väinämöinen \\
Zdjął z niedźwiedzia jego futro, & otatti otsolta turkin, \\
Zaniósł je na piętro spichrza, & pani aitan parven päähän; \\
Do kotła wsadził mięsiwo, & lihat liitti kattilahan, \\
Do naczynia złoconego, & kuparihin kullattuhun, \\
Do saganka z dnem miedzianym. & vaskipohjahan patahan. \\
\end{longtable}
}

\begin{spacing}{1.3}
%\noindent
Przekraczalność granic w~przypadku stypy po niedźwiedziu można w~\emph{Kalevali} jeszcze interpretować w~kwestii samego statusu zwierzęcia --- skoro osiągnięcie jego pięknego, dorosłego stanu umożliwiły m.in. leśne włodarki, to oznacza, że nie tylko przewodzi Väinämöinen pogrzebowi niedźwiedzia, ale również, iż uśmiercił on półboga. W~taki oto sposób trickster, będący w~bliskości w~stosunku do bogów, ludzi i~zwierząt sprawia niedźwiedziowi pożegnanie godnego samego króla, jednocześnie zaprzeczając wypadkom w~iście groteskowy sposób, w~którym również jest metoda (mająca przynieść powodzenie w~przyszłych polowaniach) (por. XLVI:~107--116):
\end{spacing}
{\centering
\small
\begin{longtable}{c c}
,,Mój misiaczku, mój jedyny, & ,,Otsoseni, ainoiseni,\\
Miodołapy mój pięknisiu! & mesikämmen, kaunoiseni!\\
Nie gniewaj się bez powodu, & Elä suutu suottakana!\\
To nie ja ciebie zabiłem, & En minä sinua kaannut:\\
Jenoś ty sam spadł z pałąka, & itse vierit vempeleltä,\\
Sturlał się z czubka choiny, & hairahit havun selältä\\
Drzewa portki ci podarły, & puhki puiset kaatiosi,\\
Stargały kurtę z igliwia & halki haljakan havuisen.\\
Gryzące pluchy jesienne & Sykysyiset säät lipeät,\\
I posępne chmurne dzionki. & päivät pilviset pimeät.\\
\end{longtable}
}

\begin{spacing}{1.3}
%\noindent
Pozostając w~temacie przejścia między światami, należy przypomnieć, że ów zaklinacz udaje się również na własne życzenie do krainy umarłych --- Tuoneli --- w~celu odnalezienia potrzebnego mu zaklęcia do budowy łodzi, co powoduje jego napięte stosunki z~boginią Tuonetar (por.~XVI:~131--412).

Ciekawą sprawą dotyczącą graniczności w~przypadku Väinämöinena jest jego tendencja do wywyższania się --- gdy wielki i~mocarny pieśniarz lamentuje po straconym instrumencie (por. XLII:~496--502), kowal Ilmarinen również zaczyna rwać sobie włosy z~głowy i~narzekać na swój los. Dopiero wtedy Väinämöinen zmienia diametralnie swoje nastawienie do sytuacji --- tutaj mamy do czynienia z~tricksterskim brakiem jednoznacznej moralności, gdyż w~cwany sposób ukazuje on swoją wyższość nad kompanem, zmuszając go wpierw (niebezpośrednio) do współpłaczu. Nawet jeśli chciano przez to pokazać jedynie umiejętność wychodzenia z~kłopotów i~dawania sobie rady w~trudnych chwilach (co jest dość częstym motywem \emph{Kalevali}), rezultat jest jak najbardziej tricksterski.

Ma w~sobie Väinämöinen również cząstkę zła. Mówi o~niej chociażby Antero Vipunen, po wnętrznościach którego urządził sobie kalevaliczny szaman iście groteskową wycieczkę w~poszukiwaniu odpowiedniego zaklęcia na zatamowanie krwotoku z~rany. Vipunen nazywa zaklinacza ,,diabelskim psem'', ,,kundlem z~Manali'', ,,znajduchem'', ,,zakałą ziemi'', ,,nasłanym złem'' oraz innymi epitetami (które później, po zakończeniu perypetii, szybko odwołuje). Dla przykładu (XVII:~239--244, 589--594):
\end{spacing}

{\centering
\small
\begin{longtable}{c c}
{[...]} & {[...]}\\
Czyś mnie stamtąd zauroczył, & ,,Sieltäkö sinäki puutuit,\\
Z owych stron, plago, naskoczył, & sieltä, taikea, tapahuit\\
By mi serce niepokoić, & syämehen syyttömähän,\\
Bezeceństwa w brzuchu broić, & vatsahan viattomahan,\\
Żeby kąsać mnie złośliwie, & syömähän, kaluamahan,\\
Rozdzierać i żreć łapczywie? & puremahan, louhtamahan?\\
{[...]} & {[...]}\\
,,Piłem i jadałem wiele, & ,,Mont' olen syönyt, monta juonut,\\
Moc tysięcy pochłonąłem, & tuhonnut tuhatlukuja;\\
Jeszczem nie jadł nic takiego & moint' en vielä konsa syönyt\\
Jak ten stary Väinämöinen, & kuin söin vanhan Väinämöisen!\\
Dobrze, że trafiłeś do mnie, & Hyvin laait tultuasi,\\
Jeszcze lepiej, że powracasz.'' & teet paremmin, kun paloat.''
\end{longtable}
}
\begin{spacing}{1.3}
%\noindent
Zło  Väinämöinena widoczne jest również w~ostatnim runie, kiedy to przychodzi mu wydać sąd nad dzieckiem Marjatty --- kieruje się on wtedy potrzebą pozbycia się konkurencji (przepowiednia mówi wszakże, że bobas zostanie wielkim królem), co widać w~słowach (L:~447--454):
\end{spacing}

{\centering
\small
\begin{longtable}{c c}
Stary mądry Väinämöinen & Vaka vanha Väinämöinen,\\
Takie wydał orzeczenie: & tuop' on tuossa tuomitsevi:\\
,,Skoro z bagna wydobyty & ,,Kun lie poika suolta saatu,\\
I od jagody pochodzi, & maalta marjasta si'ennyt,\\
Niech chłopaka w ziemię wsadzą, & poika maahan pantakohon,\\
Przy kępie jagód zostawią, & marjamättähän sivulle,\\
Lub na bagno wyprowadzą, & tahi suolle vietäköhön,\\
Palicą w głąb bagna wtłoczą!'' & puulla päähän lyötäköhön!''\\
\end{longtable}
}
\vspace*{0.5\baselineskip}
\subsection*{Shapeshifting}
\begin{spacing}{1.3}
W \emph{Kalevali} dwukrotnie mamy do czynienia z~Väinämöinenem zmieniającym swoją postać. Za pierwszym razem wpada on w~niemałe tarapaty, kiedy pani świata umarłych nie chce go wypuścić z~Tuoneli, a~z~jej rozkazu syn Tuoni zastawia na zaklinacza żelazne sieci. Väinämöinen przez chwilę biadoli nad swoim losem i~nieszczęściem, jakie go spotkało, a~następnie przeobraża się i~bezproblemowo ratuje swoje życie, wracając z~powrotem do świata żywych (XVI:~369--376):
\end{spacing}
{\centering
\small
\begin{longtable}{c c}
Rychło przybrał inną postać, & Pian muuksi muuttelihe,\\
Raptownie się przeobraził, & ruton toiseksi rupesi:\\
W wodę zapadł poczerniały, & meni mustana merehen,\\
Jak wydra między turzyce, & saarvana sara'ikkohon;\\
Zwijał się jak gad żelazny, & matoi rautaisna matona,\\
Czołgał jak wąż jadowity & kulki kyisnä käärmehenä\\
W poprzek rzeki Tuoneli, & poikki Tuonelan joesta,\\
Między sieciami Tuoni. & läpi Tuonen verkkoloista.\\
\end{longtable}
}
\begin{spacing}{1.3}
Drugie przeobrażenie Väinämöinena ma miejsce, gdy podróżuje on po wnętrznościach Antero Vipunena. Podczas gdy incydent w~Tuoneli doprowadził do zmiany jego właściwości fizycznych, w~kiszkach odwiecznego pieśniarza zmienia się Väinämöinen w~kowala (XVII:~131--140):
\end{spacing}
{\centering
\small
\begin{longtable}{c c}
Tedy stary Väinämöinen & Silloin vanha Väinämöinen\\
Już w kowala się przedzierzgnął, & löihen itsensä sepoksi, \\
Wziął się do kucia żelaza, & rakentihe rautioksi;\\
Uczynił z koszuli kuźnię, & painoi paitansa pajaksi,\\
Miechy z rękawów porobił, & hiat paian palkehiksi,\\
Popielnik zasię z kożucha, & turkkinsa tuhottimeksi,\\
Dymnik sprawił z nogawicy, & housut hormiksi rakenti,\\
Wylot dymnika z pończochy, & sukat hormin suulliseksi,\\
Kowadłem było kolano, & polvensa alasimeksi,\\
Młotkiem zasię własny łokieć. & vasaraksi kyynäspäänsä.\\
\end{longtable}
}
\begin{spacing}{1.3}
Dowodzi to jego nadnaturalnych zdolności oraz pomysłowości, dzięki którym osiągnął swój cel, którym było zdobycie potrzebnych mu informacji. Jest to jedyny przypadek w~\emph{Kalevali}, kiedy Väinämöinen przejmuje czasowo funkcje innego bohatera --- funkcja kowala jest wszakże zarezerwowana dla Ilmarinena (który paradoksalnie wyposażył Väinämöinena przed tą podróżą w~niezbędny kowalski osprzęt).
\end{spacing}
\vspace*{0.5\baselineskip}
\subsection*{Magia}
\begin{spacing}{1.3}
Związek Väinämöinena z~magią, szamanizmem i~zaklęciami jest widoczny już na samym początku eposu --- wymieniony jest on jako ten, który ,,dzierży nuty i~zaklęcia''. Według jego własnych słów, jest on ,,na wieszcza urodzony'' oraz ,,przeznaczony na pieśniarza'' (XXI:~359--360), a~zwrotki i~pieśni mu się nie kończą (XXI:~369--374):
\end{spacing}
{\centering
\small
\begin{longtable}{c c}
Śpiewał stary Väinämöinen & Lauloi vanha Väinämöinen,\\
Tak jak potrafił najlepiej & sekä lauloi jotta taitoi:\\
I słów mu nie brakowało, & ei sanat sanoihin puutu,\\
Więcej zwrotek miał w zapasie & virret veisaten vähene;\\
Niźli kamieni po skałach, & ennen kalliot kiviä,\\
Na leśnych jeziorach kwiatów. & umpilammit lumpehia.\\
\end{longtable}
}
\begin{spacing}{1.3}
%\noindent
Gdy gra on na kantele, słuchają go zarówno ludzie, jak i~zwierzęta, a~świat przepełnia radość (XLI:~18--30):
\end{spacing}
{\centering
\small
\begin{longtable}{c c}
{[...]} & {[...]}\\
Jął wygrywać znakomicie & alkoi soittoa somasti\\
Na brzękadle z rybich żeber, & hauinruotaista romua,\\
Gęślach ze szczupaczej ości; & kalanluista kanteletta.\\
Wygięte podnosił palce, & Sormet nousi notkeasti,\\
Lekko wznosząc kciuk do góry. & peukalo ylös keveni.\\
Szło wesele za weselem, & Jo kävi ilo ilolle,\\
Radość za radością biegła, & riemu riemulle remahti,\\
Po nucie zdążała nuta & tuntui soitto soitannalle,\\
I pieśń z pieśni wynikała; & laulu laululle tehosi.\\
Podźwiękiwał ząb szczupaka, & Helähteli hauin hammas,\\
Ogon ryby się unosił, & kalan pursto purkaeli,\\
Zawodziły źrebców grzywy, & ulvosi upehen jouhet,\\
Dźwięczało włosie rumaka. & jouhet ratsun raikkahuivat.\\
\end{longtable}
}
\newpage
\begin{spacing}{1.3}
%\noindent
Zaklinania Väinämöinena są potężne i~nierzadko mogą kontrolować zachowanie natury (III:~295--300):
\end{spacing}
{\centering
\small
\begin{longtable}{c c}
Śpiewał stary Väinämöinen: & Lauloi vanha Väinämöinen:\\
Mżyła toń i drżała ziemia, & järvet läikkyi, maa järisi,\\
Trzęsły się spiżowe góry, & vuoret vaskiset vapisi,\\
Twarde głazy rozpadały, & paaet vahvat paukahteli,\\
Skały pękały na poły, & kalliot kaheksi lenti,\\
Calizny na brzegu morza. & kivet rannoilla rakoili.\\
\end{longtable}
}
\begin{spacing}{1.3}
%\noindent
Umożliwiają mu również panowanie nad prawami fizyki i~przechodzenie ,,niemożliwych'' prób, jakie stawiają przed nim inni bohaterowie, zwykle przy okazji jego groteskowych miłosnych podbojów (VIII:~100--104):
\end{spacing}
{\centering
\small
\begin{longtable}{c c}
{[...]} & {[...]}\\
W mig niteczkę przepołowił & jouhen halki halkaisevi\\
Nożem, nie mającym ostrza, & veitsellä kärettömällä,\\
Niedostrzegalnym sposobem, & aivan tutkaimettomalla;\\
Zawiązał jajko na węzeł, & munan solmuhun vetävi\\
Co się wyczuć nie zdał zgoła; & solmun tuntumattomaksi.\\
{[...]} & {[...]}\\
\end{longtable}
}
\begin{spacing}{1.3}
Kalevaliczny szamanizm w~osobie Väinämöinena jest widoczny we fragmentach związanych z~wszelkimi dolegliwościami, którym należy zaradzić, począwszy od skaleczenia, a~skończywszy na torsjach, transie i~konwulsyjnej walce z~żywiołem bólu. W~pierwszym przypadku widzimy, jak ważną rolę odgrywają w~fińskim eposie zaklęcia, będące kluczem do zrozumienia praprzyczyny określonego zjawiska (czego brak w~przypadku Väinämöinena skończył się zalaniem świata krwią z~jego kolana) (VIII:~177--186):
\end{spacing}
{\centering
\small
\begin{longtable}{c c}
Jął tedy zamawiać ranę    & Loihe siitä loitsimahan,\\
Pomocne dobierać słowa  & sai itse sanelemahan.\\
Wedle przyczyn wszystkich rzeczy, & Luki synnyt syitä myöten,\\
Przyrody wszego stworzenia, & luottehet lomia myöten,\\
Wszakże nie pamiętał wielu & mutt' ei muista muutamia\\
Słów skutecznych na żelazo, & rauan suuria sanoja,\\
Coby krzepko ryglowały, & joista salpa saataisihin,\\
Zamek bezpieczny dawały & luja lukko tuotaisihin\\
Przeciw rozdarciu żelazem, & noille rauan ratkomille,\\
Skaleczeniu siną stalą. & suu sinervän silpomille.\\
\end{longtable}
}
\begin{spacing}{1.3}
%\noindent
W~drugim zaś przypadku mamy do czynienia z~klasycznym szamańskim transem, podczas którego wojuje Väinämöinen z~bólem za pomocą umysłu (jego ciało pozostaje bezwładne, narażone na wszelkiego rodzaju pomiatania --- por. Haavio~1979: 341) (IX:~517--528):
\end{spacing}
\newpage
{\centering
\small
\begin{longtable}{c c}
Ledwie przyłożono maści, & Kun oli voie päälle pantu,\\
Potężnych użyto leków, & nuot on katsehet vakaiset,\\
Jęło trząść Väinämöinena, & murti se puolipyörryksihin,\\
Bez czucia nim pomiatało: & Väinämöisen väännyksihin:\\
Tłukło tędy i owędy, & lyökse sinne, lyökse tänne,\\
Spokoju zaznać nie dało. & vaan ei löytänyt lepoa.\\
Tak starzec z bólem wojował, & Niin ukko kipuja kiisti,\\
Dolegliwość swą wypędzał & työnti tuosta tuskapäitä\\
Do wnętrza Góry Boleści, & keskelle Kipumäkeä,\\
Na wierzchołek Kopca Bólu, & Kipuvuoren kukkulalle\\
By kamienie tam nękała, & kiviä kivistämähän,\\
Żeby tam dręczyła skały,& paasia pakottamahan.\\
{[...]} & {[...]}\\
\end{longtable}
}
\begin{spacing}{1.3}
Używa czarów Väinämöinen również do spełniania swoich pilnych potrzeb, np. wyśpiewania łodzi gotowej do podróży, a~także załogi w~pełnym oporządzeniu (XVI: 101--112; XXXIX: 275--293):
\end{spacing}
{\centering
\small
\begin{longtable}{c c}
Tedy stary Väinämöinen,		               &  Siitä vanha Väinämöinen\\
Wieszcz wiedzący od wiek wieka,	&  tietäjä iän-ikuinen,\\
Łódź zgotował swoją wiedzą,		&  teki tieolla venettä,\\
Czółno czynił czarowaniem,		&  laati purtta laulamalla\\
Klecąc je z dębowych desek,		&  yhen tammen taittumista,\\
Z bierwion grubo narąbanych.		&  puun murskan murenemista.\\
Pieśń zaśpiewał – dno ułożył,		&  Lauloi virren: pohjan puutti,\\
Drugą pieśnią burty spoił,		&  lauloi toisen: liitti laian;\\
Jakoż i zaśpiewał trzecią,			&  lauloi kohta kolmannenki\\
Gdy dla wioseł dulki ciosał,		&  hankoja hakatessansa,\\
Kiedy końce wręg zakładał,		&  päitä kaaren päätellessä,\\
Kiedy zapuszczał je w pazy.		&  liitellessänsä limiä.\\
{[...]} 					& {[...]}\\
Tedy stary Väinämöinen			& Siitä vanha Väinämöinen\\
Mamrotać począł zaklęcia: 		& lauloa hyrähtelevi.\\
Na jednej burcie wyśpiewał		& Lauloi ensin laitapuolen\\
Zalotników z gładką głową, 		& sukapäitä sulhosia,\\
Z gładką głową, krzepką dłonią, 		& sukapäitä, piipioja,\\
Na nogach się buty mienią; 		& saapasjalkoja jaloja.\\
Na drugiej burcie wyśpiewał		& Lauloi toisen laitapuolen\\
Panny w cynowej ozdobie, 		& tinapäitä tyttäriä,\\
Cynopierśne, w pasach z miedzą, 		& tinapäitä, vaskivöitä,\\
Ze złotymi pierścionkami. 		& kultasormia somia.\\
Do tego wyśpiewał Väinö		& Lauloi vielä Väinämöinen\\
Pełno na ławach załogi, 			& teljot täytehen väkeä,\\
Wszystko byli ludzie starzy, 		& ne on vanhoa väkeä,\\
Siedzący przez całe życie, 		&  iän kaiken istunutta,\\
Miejsca mieli nie za wiele, 		& kuss' oli vähän sijoa\\
Skoro wpierw zasiedli młodzi. 		& nuorukaisilta esinnä.\\
Sam usadził się na rufie, 			&  Itse istuvi perähän,\\
Koło steru brzozowego, 			& kokan koivuisen kuvulle,\\
By pokierować komiegą, 		& lasketteli laivoansa.\\
{[...]} 					& {[...]}\\
\end{longtable}
}
\newpage
\begin{spacing}{1.3}
%\noindent
Pilną potrzebą jest również otwarcie skrytki, w~której czarownica Louhi zamknęła młynek Sampo. Fakt, iż Ilmarinen musiał dodatkowo natłuszczać masłem zasuwy (żeby nie zgrzytnęły przy otwieraniu) rzuca cień na jakość czarów Väinämöinena w~dziedzinie usypiania swoich przeciwników (tutaj --- obrońców Sampo) grą na kantele oraz profilaktycznym zaszywaniem powiek (XLII:~101--112):
\end{spacing}
{\centering
\small
\begin{longtable}{c c}
Stary mądry Väinämöinen		& Tuossa vanha Väinämöinen\\
Zaczął mamrotać zaklęcia		& lauloa hyrähtelevi\\
U wrót wierchu miedzianego, 		& vaaran vaskisen ovilla,\\
W bramie zamku kamiennego; 		& kivilinnan liepehillä:\\
Ugną się zamku zasuwy, 			& jopa liikkui linnan portit,\\
Zadrżą żelazne zawiasy. 			& järkkyi rautaiset saranat.\\
Zasię kowal Ilmarinen, 			& Itse seppo Ilmarinen,\\
Który towarzyszył tutaj, 			& tuopa tuossa toisna miesnä.\\
Namaścił masłem zasuwy, 		& Voilla voiti lukkoloita,\\
Tłuszczem pomazał zawiasy, 		& saranoita rasvasilla,\\
Żeby nie zgrzytnęły wrota, 		& jottei ukset ulvahuisi\\
Nie zajęczały zawiasy; 			& eikä naukuisi saranat.\\
{[...]}					& {[...]}\\
\end{longtable}
}
\begin{spacing}{1.3}
Kolejnym celem väinämöinenowego czarowania jest wpędzanie kompanów w~tarapaty, czego najlepszym przykładem jest kowal Ilmarinen --- najpierw zwabiony podstępem do obejrzenia wspaniałego drzewa, na czubku którego ma świecić księżyc i~gwiazdy, a~potem błyskawicznie poniesiony wiatrem na północ do Pohjoli (z~czego druga część incydentu ma na celu również zapobiegnięcie skompromitowania Väinämöinena przez mówiące drzewo) (X:~31--42, 159--169):
\end{spacing}
{\centering
\small
\begin{longtable}{c c}
Po czym stary Väinämöinen		    &  Siitä vanha Väinämöinen\\
Wyśpiewywał i zaklinał, 			&  laulelevi, taitelevi:\\
Wyczarował świerk strzelisty, 		&  lauloi kuusen kukkalatvan,\\
Cały w świeczkach, w igłach złotych,&  kukkalatvan, kultalehvän;\\
Co wierzchołek wzniósł pod niebo, 	&  latvan työnti taivahalle,\\
Na wylot przebił obłoki, 			&  puhki pilvien kohotti,\\
Rozpostarł w przestwór konary, 		&  lehvät ilmoille levitti,\\
Rozgałęził je w niebiosa. 		    &  halki taivahan hajotti.\\
{[...]}				             	&  {[...]}\\
Wyśpiewał miesiąc, by świecił		&  lauloi kuun kumottamahan\\
W świerku ze złotym wierzchołkiem, 	&  kultalatva-kuusosehen,\\
A z nim Wielką Niedźwiedzicę. 		&  lauloi oksillen otavan.\\
{[...]}					        &  {[...]}\\
Wówczas stary Väinämöinen		&  Silloin vanha Väinämöinen\\
Uciekł się do czarowania, 		&  lauloa hyrähtelevi:\\
Wyśpiewał trąbę powietrzną, 	&  lauloi tuulen tuppurihin,\\
Pobudził wicher zajadły, 		&  ilman raivohon rakenti;\\
Rzecze takimi słowami, 			&  sanovi sanalla tuolla,\\
Takie wypowie zaklęcie: 		&  lausui tuolla lausehella:\\
,,Weź go na swe czółno, wietrze,&  ,,Ota, tuuli, purtehesi,\\
Na łódź, Podmuchu Północny, 	&  ahava, venosehesi\\
Porwij do swojej Pohjoli, 		&  vieä vieretelläksesi\\
Ziemi w mroku pogrążonej.'' 	&  pimeähän Pohjolahan!''\\
\end{longtable}
}
\begin{spacing}{1.3}
Zaklinanie widzimy również przed omawianą już wcześniej stypą po niedźwiedziu. Ma ono na celu ochronę Väinämöinena podczas polowania. Po polowaniu natomiast inkantuje Väinämöinen kolejne zaklęcie --- tym razem o~charakterze dziękczynnym, skierowane do Stwórcy (XLVI:~101--104):
\end{spacing}
{\centering
\small
\begin{longtable}{c c}
,,Pochwalony bądź, Jumala,	    & ,,Ole kiitetty, Jumala,\\
Wielbię cię, Stwórco jedyny, 	& ylistetty, Luoja yksin,\\
Żeś misia mi dał w udziale,	    & kun annoit otson osaksi,\\
Złoto puszczy jako zdobycz!'' 	& salon kullan saalihiksi!''\\
\end{longtable}
}

\vspace*{0.5\baselineskip}
\subsection*{Sacrofanum}
\begin{spacing}{1.3}
Połączenia świętości ze sprośnością w~\emph{Kalevali} nie należy szukać w~dosłownym tych słów znaczeniu. Sferę \emph{sacrum}, czyli eliadowski ,,czas święty'' stanowią obrzędy szamańskie, zaklinania oraz teksty związane z~bliską obecnością bóstwa przy danym bohaterze. Sferę \emph{profanum} oraz ,,czas świecki'' stanowią zatem wszelkie odniesienia do cielesności i~tricksterskiej chuci, która w~przypadku Väinämöinena jest widoczna w~incydencie z~Aino --- wyzwany przez Joukahainena na pojedynek zaklinacz akceptuje ofertę ,,najobfitszego okupu'' (III:~345--352), którym okazuje się być siostra pokonanego. Väinämöinen bardzo się cieszy z~takiego obrotu sprawy, czego nie można powiedzieć o~samej Aino.

,,Sprośności'' Väinämöinena należy doszukiwać się w~jego bezpośrednim stosunku do kobiet, czego przykładem jest jego pierwsze spotkanie z~Aino (IV:~9--20):
\end{spacing}
{\centering
\small
\begin{longtable}{c c}
Już do domu powracała, & Jo astui kohin kotia,\\
Olszniakiem śpieszyła żwawo, & lepikköä leuhautti.\\
Aż tu stary Väinämöinen & Tuli vanha Väinämöinen;\\
Zobaczył w gaju dziewczynę, & näki neitosen lehossa,\\
Smukłą kibić śród murawy & hienohelman heinikössä.\\
I takimi rzecze słowy: & Sanan virkkoi, noin nimesi:\\
,,Nie dla innych, młoda panno, & ,,Eläpä muille, neiti nuori,\\
Jeno dla mnie, młoda panno, & kuin minulle, neiti nuori,\\
Nosić masz perły na szyi, & kanna kaulanhelmilöitä,\\
Krzyżyk na piersi zawieszać, & rinnanristiä rakenna,\\
Warkoczem owijać głowę, & pane päätä palmikolle,\\
Jedwabiem włosy oplatać!'' & sio silkillä hivusta!''\\
\end{longtable}
}
\begin{spacing}{1.3}
%\noindent
W~kolejnych runach sytuacje damsko-męskie (np. spotkanie z~Siewierzanką) będą wyglądały podobnie (VIII:~34--50):
\end{spacing}
{\centering
\small
\begin{longtable}{c c}
{[...]} & {[...]}\\
Tak się ozwie i przemówi: & itse lausui, noin nimesi:\\
,,Zejdź, dziewczyno, do mych sani, & ,,Tule, neiti, korjahani,\\
Do wasąga wsiądź mojego.'' & laskeite rekoseheni!''\\ \newpage
{[...]} & {[...]}\\
,,Na to zda się w saniach panna, & ,,Siksi neittä korjahani,\\
W moim wasągu niewiasta: & tyttöä rekoseheni:\\
By piekła słodkie kołacze, & mesileivän leipojaksi,\\
Żeby piwo mi warzyła, & oluen osoajaksi,\\
Ośpiewała każdy stołek & joka lautsan laulajaksi,\\
I okna opromieniała & ikkunan iloitsijaksi\\
W okolicach mej Väinöli, & noilla Väinölän tiloilla,\\
Na podwórcach Kalevali.'' & Kalevalan kartanoilla.''\\
\end{longtable}
}
%\begin{spacing}{1.3}
\vspace*{0.5\baselineskip}
\subsection*{Deus ex machina}
\begin{spacing}{1.3}
Väinämöinen jako postać tricksterska ma w~wielu sytuacjach niebywałe szczęście, dzięki któremu udaje mu się uniknąć poważnych kłopotów. Po raz pierwszy mamy do czynienia z~nagłym rozwiązaniem akcji w~runie siódmym, kiedy to dryfującego po morzu zaklinacza ratuje orzeł (na podobieństwo uwięzionego w~Orthanku tolkienowskiego Gandalfa) (VII:~90--116):
\end{spacing}
{\centering
\small
\begin{longtable}{c c}
,,Nie martw się bynajmniej o to, & ,,Ellös olko milläskänä!\\
Jeno siadaj mi na grzbiecie, & Seisotaite selkähäni,\\
Wdrap się na skrzydła szerokie, & nouse kynkkäluun nenille!\\
Ponad morzem cię przeniosę, & Mie sinun merestä kannan,\\
{[...]} & {[...]}\\
Uniósł wzwyż Väinämöinena, & kantoi vanhan Väinämöisen,\\
Wraz z wichurą go unosił, & viepi tuulen tietä myöten,\\
Zimnego wiatru koleją & ahavan ratoa myöten\\
W odległe kraje północy, & Pohjan pitkähän perähän,\\
Do Sarioli mgłą spowitej; & summahan Sariolahan.\\
Tam zrzucił Väinämöinena, & Siihen heitti Väinämöisen,\\
A sam uniósł się w powietrze. & itse ilmahan kohosi.\\
\end{longtable}
}
\begin{spacing}{1.3}
Innym razem zwabia Väinämöinen Ilmarinena do Pohjoli pod pretekstem podziwiania wspaniałego drzewa, lecz gdy postęp wychodzi na jaw, a~kowal orientuje się, że jest jedynie kartą przetargową, wielki piewca improwizuje, by ratować swój plan od klęski (X:~159--169):
\end{spacing}
{\centering
\small
\begin{longtable}{c c}
Wówczas stary Väinämöinen & Silloin vanha Väinämöinen \\
Uciekł się do czarowania, & lauloa hyrähtelevi: \\
Wyśpiewał trąbę powietrzną, & lauloi tuulen tuppurihin, \\
Pobudził wicher zajadły, & ilman raivohon rakenti; \\
Rzecze takimi słowami, & sanovi sanalla tuolla, \\
Takie wypowie zaklęcie: & lausui tuolla lausehella: \\
,,Weź go na swe czółno, wietrze, & ,,Ota, tuuli, purtehesi, \\
Na łódź, Podmuchu Północny, & ahava, venosehesi \\
Porwij do swojej Pohjoli, & vieä vieretelläksesi \\
Ziemi w mroku pogrążonej.'' & pimeähän Pohjolahan!'' \\
\end{longtable}
}
\newpage
\begin{spacing}{1.3}
Kolejnym przykładem rozwiązującego akcję boga z~maszyny jest przemiana Väinämöinena (XVI:~369--376), która umożliwia mu ucieczkę ze świata bogini Tuonetar --- w~którą to ucieczkę jeszcze chwilę wcześniej sam bardzo wątpi, tak jak wątpił w~wydostanie się z~kiszek Antero Vipunena (XVII:~115--118):
\end{spacing}
{\centering
\small
\begin{longtable}{c c}
,,Chyba przyjdzie na mnie zguba, 		&  ,,Jo taisi tuhoni tulla,\\
Zła godzina mnie pogłębi		&  hätäpäivä hämmenteä\\
W tych diabelskich zakamarkach, 	&  tämän hiien hinkalossa,\\
W owej zagródce Kostuchy.'' 		&  tämän kalman karsinassa.''\\
\end{longtable}
}
\begin{spacing}{1.3}
\emph{Deus ex machina} w~wykonaniu pieśniarza widać również po uzgodnieniu z~Joukahainenem wysokości okupu po przegranym pojedynku na wiedzę (III:~475--476):
\end{spacing}
{\centering
\small
\begin{longtable}{c c}
{[...]} & {[...]}\\
Wspak obrócił święte słowa, 		&  pyörti pois pyhät sanansa,\\
Odwołał wszystkie zaklęcia. 		&  perin laski lausehensa.\\
\end{longtable}
}
\begin{spacing}{1.3}
Ostatnim przykładem tricksterskiego wyjścia z~opresji jest ucieczka przed łodzią plemienia Pohjoli po kradzieży Sampo. Zastanawia się Väinämöinen długo, jak rozwiązać problem, po czym wrzuca do wody hubkę i~krzesiwo, które na skutek zaklęcia zamieniają się w~wielką skałę pośrodku morza, o~którą rozbija się łódź wojaków Louhi (XLIII: 107--134).
\end{spacing}

\vspace*{0.5\baselineskip}
\subsection*{Kłamstwo}
\begin{spacing}{1.3}

Väinämöinen jest notorycznym kłamcą. Oszukuje zarówno swoich kompanów (m.in. wysłanego podstępem do Pohjoli Ilmarinena), jak również same bóstwa. Co ciekawe, uparcie trwa w~swoich zmyślonych wersjach wydarzeń i~dopiero po wielu dowodach na kłamstwo decyduje się wyjawić prawdę, nieświadomy nierozsądności swojego zachowania. Idealnym przykładem Väinämöinena-kłamcy jest wymiana zdań między nim a~Tuonetar, do królestwa której udaje się on w~poszukiwaniu zaklęć (XVI:~159--262):
\end{spacing}
{\centering
\small
\begin{longtable}{c c}
Stary mądry Väinämöinen		&  Vaka vanha Väinämöinen \\
Wołać jął i pohukiwać		&  jo huhuta huikahutti \\
Poza rzekę Tuoneli, 		&  tuossa Tuonelan joessa, \\
Ponad niziną Manali: 		&  Manalan alantehessa: \\
,,Przywiedź łódź, córo Tuoni, 		&  ,,Tuo venettä, Tuonen tytti, \\
Sprowadź prom, Manali dziecię, 		&  lauttoa, Manalan lapsi, \\
Niech na drugi brzeg przejadę, 		&  yli salmen saa'akseni, \\
Przedostanę się przez rzekę!'' 		&  joen poikki päästäkseni!'' \\
Niepozorna córka Tuoni, 		&  Lyhykäinen Tuonen tytti, \\
Niska niewiastka Manali			&  matala Manalan neiti, \\
Prała chusty pobrukane, 		&  tuo oli poukkujen pesijä, \\
Biła bieliznę kijanką			&  räpähien räimyttäjä \\
W czarnej rzece Tuoneli, 		&  Tuonen mustassa joessa, \\
W podziemnych wodach Manali; 		&  Manalan alusve'essä. \\
Tak ozwała się do niego, 			&  Sanan virkkoi, noin nimesi, \\
Te wypowiedziała słowa: 		&  itse lausui ja pakisi: \\
,,Łódkę tutaj ci sprowadzę, 		&  ,,Vene täältä tuotanehe, \\
Jeśli podasz mi powody, 			&  kuni syy sanottanehe, \\
Z jakich jesteś tu w Manali		&  mi sinun Manalle saattoi \\
Nie uśmiercony zarazą, 			&  ilman tauin tappamatta, \\
Ani też zwyczajnym zgonem, 		&  ottamatta oivan surman, \\
Ni żadnym innym sposobem.'' 		&  muun surman musertamatta.'' \\
{[...]} & {[...]}\\
,,Tuoni mnie sprowadziła, 		&  ,,Tuoni minun tänne tuotti, \\
Mana ściągnęła ze świata.'' 	&  Mana mailtani veteli.'' \\
{[...]} & {[...]}\\
,,Tom trafiła na bajarza! 			&  ,,Jopa keksin kielastajan! \\
Gdyby cię przywiodła Tuoni, 		&  Kunp' on Tuoni tänne toisi, \\
Mana zabrała ze świata, 			&  Mana mailta siirteleisi, \\
Tuoni dałaby na drogę, 			&  Tuoni toisi tullessansa, \\
Na podróż Pani Podziemia		&  Manalainen matkassansa	 \\
Kaptur Tuoni na ramiona, 		&  Tuonen hattu hartioilla, \\
Do rąk Many rękawice, 			&  Manan kintahat käessä. \\
Prawdę mów, Väinämöinenie, 		&  Sano totta, Väinämöinen: \\
Co cię tutaj sprowadziło?'' 		&  mi sinun Manalle saattoi?'' \\
{[...]} & {[...]}\\
,,Żelazo mnie tu przywiodło, 		&  ,,Rauta mun Manalle saattoi, \\
Stal strąciła do Tuoni.'' 			&  teräs tempoi Tuonelahan.'' \\
{[...]} & {[...]}\\
,,I po tym poznać bajarza! 		&  ,,Tuosta tunnen kielastajan! \\
Gdyby tknęło cię żelazo, 			&  Kun rauta Manalle saisi, \\
Stal strąciła do Tuoni, 			&  teräs toisi Tuonelahan, \\
Odzież krwią by ociekała, 		&  verin vaattehet valuisi, \\
Krwią broczyłyby twe rany, 		&  hurmehen hurahteleisi. \\
Prawdę mów, Väinämöinenie, 		&  Sano totta, Väinämöinen, \\
Drugim razem prawdę powiedz!'' 	&  sano totta toinen kerta!'' \\
{[...]} & {[...]}\\
,,Woda zniosła mnie do Many, 		&  ,,Vesi sai minun Manalle, \\
Przypływ zniósł do Tuoneli.'' 		&  aalto toi on Tuonelahan.'' \\
{[...]} & {[...]}\\
,,Przejrzałam na wylot kłamcę! 		&  ,,Ymmärrän valehtelijan! \\
Gdyby woda cię przyniosła, 		&  Jos vesi Manalle saisi, \\
Przypływ zniósł do Tuoneli, 		&  aalto toisi Tuonelahan, \\
Ubiór wodą by ociekał, 			&  vesin vaattehet valuisi,	 \\
Wilgoć zeń by się sączyła; 		&  helmasi herahteleisi. \\
Powiedz prawdę bez ogródek, 		&  Sano tarkkoja tosia: \\
Co się tutaj sprowadziło? 		&  mi sinun Manalle saattoi?'' \\
Tedy stary Väinämöinen			&  Tuossa vanha Väinämöinen \\
Zmyślił jeszcze raz opowieść: 		&  vielä kerran kielastavi: \\
,,Ogień mnie tu przyprowadził, 		&  ,,Tuli toi mun Tuonelahan, \\
Do Many przywiodło płomię.'' 		&  valkea Manalle saattoi.'' \\
{[...]} & {[...]}\\
,,Odgaduję bałamuta! 			&  ,,Arvoan valehtelijan! \\
Gdyby ogień cię sprowadził, 		&  Jos tuli Manalle toisi, \\
Białe płomię do Tuoneli, 			&  valkeainen Tuonelahan, \\
W strąkach miałbyś swe kędziory, 	&  oisi kutrit kärventynnä, \\
Brzydko przypaloną brodę. 		&  partaki pahoin palanut. \\
{[...]} & {[...]}\\
I już dalszych kłamstw zaniechaj, 	&  valehia viimeisiä, \\
{[...]} & {[...]}\\
,,Jeżeli skłamałem nieco, 		&  ,,Jos vähän valehtelinki, \\
Wymyśliłem raz i drugi, 			&  kerran toisen kielastelin, \\
Prawdę jednakże wyjawię: 		&  toki ma sanon toetki. \\
Wiedzą łódź wybudowałem, 		&  Te'in tieolla venettä, \\
Czółno czyniłem czarami, 		&  laain purtta laulamalla. \\
Zaklinałem dzień i drugi, 		&  Lauloin päivän, lauloin toisen, \\
Dnia trzeciego na ostatek		&  niin päivällä kolmannella \\
Stuknęły sanie śpiewaka, 		&  rikkoihe reki runoilta, \\
Pękła płoza pieśniarzowi; 		&  jalas taittui lausehilta: \\
Wybrałem się więc po wiertło, 		&  läksin Tuonelta oroa, \\
Świder chciałbym mieć z Tuoni		&  Manalalta vääntiätä \\
Żeby naprawić saneczki, 			&  rekosen rakentoani, \\
Wasąg z pieśniami odnowić; 		&  laulukorjan laatiani. \\
Przywiedź mi tutaj łódeczkę, 		&  Tuopa nyt venoista tänne, \\
Prom do brzegu mi przyprowadź, 	&  laita mulle lauttoasi \\
Żebym przebyć mógł cieśninę, 		&  yli salmen saa'akseni, \\
Przeprawić się poprzez rzekę!'' 		&  joen poikki päästäkseni!'' \\
\end{longtable}
}
\begin{spacing}{1.3}
%\noindent
Ciekawe jest również to, iż Väinämöinen sam również wytyka kłamstwo innym (XLIII:~50--54):
\end{spacing}
{\centering
\small
\begin{longtable}{c c}
,,Już zełgałeś po próżnicy! 		&  ,,Jo vainen valehtelitki! \\
Bo to chmurka być nie może, 		&  Ei se pilvi ollekana, \\
Ani obłoku okrawek, 			&  pilven lonka lienekänä: \\
Jest to czółno pod żaglami; 		&  se on pursi purjehinen. \\
Spojrzyj baczniej po raz drugi!'' 		&  Katso toiste tarkemmasti!'' \\
\end{longtable}
}
\vspace*{0.5\baselineskip}
\subsection*{Idiotyzm}
\begin{spacing}{1.3}
Väinämöinen często postępuje bardzo niemądrze, zarówno w~kwestiach poważnych, jak i~w~drobnych perypetiach. Wielki zaklinacz, który pokonuje Joukahainena w~pojedynku na wiedzę o~początkach świata, i~który nie poznaje własnej matki (jednej z~pierwszych bogiń), ma pecha jeśli chodzi o~mądrość praktyczną, m.in. w~zakresie wędkarstwa --- łowi on rybę, po czym nie orientuje się, że złowił swoją utraconą Aino (V:~59--72):
\end{spacing}
{\centering
\small
\begin{longtable}{c c}
Przyglądał się jej, obracał, 		&Katselevi, kääntelevi. \\
Te wypowiadając słowa: 		&Sanan virkkoi, noin nimesi: \\
,,Czy ta ryba to jest ryba, 			&,,Onp' on tuo kala kalanen, \\
Nie widzi mi się to zgoła, 		&kun en tuota tunnekana! \\
Błyszcząca jest niby sieja, 		&Sileähk' on siikaseksi, \\
Niby troć połyskująca, 			&kuleahka kuujaseksi, \\
Jaśniejąca niczym szczupak, 		&haleahka haukiseksi, \\
Od ikrzaka wszak smuklejsza, 		&evätöin emäkalaksi; \\
Kształtna zda się jako człowiek, 		&ihala imehnoksiki, \\
Wstążkę wokół ma jak panna, 		&päärivatoin neitoseksi, \\ \pagebreak
Lecz nie pas, niczym syrena, 		&vyötöin veen on tyttöseksi, \\
Uszy ma nie jak gołąbka, 		&korvitoin kotikanaksi: \\
Rozmiar morskiego łososia, 		&luopuisin meriloheksi, \\
Okonia z wielkiej głębiny.'' 		&syvän aallon ahveneksi.'' \\
\end{longtable}
}
\begin{spacing}{1.3}
Innym razem zaś wybiera się z~Ilmarinenem żeby złowić tylko jedną rybę, ale wcześniej nakazuje stworzyć do tego celu specjalne sieci rybackie. Kiedy długie na pięćset sążni akcesoria nie przynoszą oczekiwanego efektu, ucieka się Väinämöinen do pomocy bóstwa, proponując wymianę odzieży (XLVIII:~123--134):
\end{spacing}
{\centering
\small
\begin{longtable}{c c}
,,Vellamo, włodarko wody, 		& ,,Vellamo, ve'en emäntä, \\
Pani wód z piersią w sitowiu, 		& ve'en eukko ruokorinta! \\
Przybądź zmienić swą koszulę, 		& Tules paian muuttelohon, \\
Pozamieniać się na stroje! 		& vaattehen vajehtelohon! \\
Ty masz trzcinową koszulę		& Sinull' on rytinen paita, \\
I pierzynę z morskiej piany, 		& merenvaahtivaippa päällä, \\
Którą córa wiatru tkała, 			& tuulen tyttären tekemä, \\
Aallotar, fal dziewoja: 			& Aallottaren antelema: \\
Ja ci lnianą oddam w zamian, 		& minä annan liinapaian, \\
Czyściuteńką ci nałożę, 			& panen aivan aivinaisen; \\
Córa Księżyca ją tkała, 			& se on Kuuttaren kutoma, \\
Córa Słoneczka uprzędła. 		& Päivättären kehreämä. \\
\end{longtable}
}
\begin{spacing}{1.3}
%\noindent
Potem jeszcze prosi on Ahto o~zagonienie wszystkich ryb w~morzu do sieci przy użyciu laski. W~końcu osiąga swój cel (XLVIII:~180--200), aczkolwiek ogrom wkładu własnego do złapania jednego szczupaka jest nie tylko groteskowy, co po prosty idiotyczny.

Zadaje również Väinämöinen ogólne pytania na temat funkcjonowania świata, które również nie grzeszą błyskotliwością (VII:~33--42; XVI:~125--130):
\end{spacing}
{\centering
\small
\begin{longtable}{c c}
Nic mi zgoła nie wiadomo, 		& ,,Enkä tuota tieäkänä, \\
Jak tu sobie w życiu radzić		& miten olla, kuin eleä \\
W owym nieprzyjaznym czasie, 		& tällä inhalla iällä, \\
U schyłku wieku starego: 		& katovalla kannikalla: \\
Czy budować dom na wietrze, 		& tuulehenko teen tupani, \\
Czy na falach izbę ciosać? 		& vetehenkö pirtin veistän? \\
Jak na wietrze dom zbuduję, 		& ,,Teen mä tuulehen tupani: \\
Nie znajdzie na nim oparcia, 		& ei ole tuulessa tukea; \\
Jak wyciosam chatę w wodzie, 		& veistän pirttini vetehen: \\
Woda ciesiołkę uniesie.'' 			& vesi viepi veistokseni.'' \\
{[...]}					& {[...]}\\
Zastanawiał się, zgadywał, 		& Arvelee, ajattelevi, \\
Skąd by tu wydostać słowa, 		& mistäpä sanoja saisi, \\
Skuteczne znaleźć zaklęcie: 		& loisi lempiluottehia; \\
Może by z jaskółczych łebków, 		& pääskyjenkö päälaelta, \\
W kluczu ciągnących łabędzi, 		& joutsenkarjan juonen päästä, \\
Lub z łopatek stada gęsi? 		& hanhilauman hartioilta? \\
\end{longtable}
}
\pagebreak
\begin{spacing}{1.3}
%\noindent
Kiedy indziej kaleczy się zaklinacz siekierą (toporem), odpowiednio to komentując (VIII:~169--176):
\end{spacing}
{\centering
\small
\begin{longtable}{c c}
,,Toporze zadartonosy, & ,,Oi sie kirves kikkanokka, \\
Równo ostrzona siekiero, & tasaterä tapparainen! \\
Myślałaś, że to budulec, & Luulitko puuta purrehesi, \\
Że to na czółenko sosna, & honkoa hotaisnehesi, \\
Choina do wytrzebienia, & petäjätä pannehesi, \\
Pień brzozowy do zrąbania, & koivua kohannehesi, \\
Żeś się w mięso ześliznęła. & kun sa lipsahit liha'an, \\
Że ugodziłaś mnie w żyłę!'' & solahutit suonilleni?'' \\
\end{longtable}
}
\begin{spacing}{1.3}
Zdarza się również Väinämöinenowi (wespół z~Sampsą) ściąć osikę, ktora ma potem do niego o~to pretensje (przecież łódź z~niej wykonana by przeciekała). Po osice pada zwichrowana sosna (zaklinacz pyta ją o~potencjalne bycie budulcem nie czekając na odpowiedź), a~dopiero za trzecim razem uzyskuje najpierw Väinämöinen opinię drzewa (tym razem dębu). Gdyby nie ten nagły przebłysk inteligencji, przez upartego dziada i~wyrobnika mogłoby paść na darmo pół lasu (por.~XVI).

Głupota praktyczna Väinämöinena objawia się podczas podróży do wnętrzności Antero Vipunena, kiedy to znachor dość późno orientuje się, iż może zostać zjedzony (XVII:~115--118), jak również gdy grabi on morze w~poszukaniu kantele. Podczas wykonywania pieśni śpiewa o~tym, że pomylił księżyc ze sprzączką, kołnierz koszuli ze słońcem i~czepiec na głowie --- ze statkiem (XXV:~643--648).

Również zaklęcia na stypie po niedźwiedziu nie są na pierwszy rzut oka zbyt bystre (XLVI:~107--116):
\end{spacing}
{\centering
\small
\begin{longtable}{c c}
,,Mój misiaczku, mój jedyny, & ,,Otsoseni, ainoiseni, \\
Miodołapy mój pięknisiu! & mesikämmen, kaunoiseni! \\
Nie gniewaj się bez powodu, & Elä suutu suottakana! \\
To nie ja ciebie zabiłem, & En minä sinua kaannut: \\
Jenoś ty sam spadł z pałąka, & itse vierit vempeleltä, \\
Sturlał się z czubka choiny, & hairahit havun selältä, \\
Drzewa portki ci podarły, & puhki puiset kaatiosi, \\
Stargały kurtę z igliwia & halki haljakan havuisen. \\
Gryzące pluchy jesienne & Sykysyiset säät lipeät, \\
I posępne chmurne dzionki. & päivät pilviset pimeät. \\
\end{longtable}
}
\begin{spacing}{1.3}
Ciekawy jest również fakt, iż często głupio postępujący Väinämöinen-trickster również sam wytyka innym niemądre postępowanie, co widać bardzo wyraźnie, kiedy Ilmarinen przychodzi do niego po poradę (okazuje się, że jego nowa wykuta ze złota żona nie jest tak wspaniała jak miała być). Zaklinacz komentuje (XLIX:~55--58):
\end{spacing}
{\centering
\small
\begin{longtable}{c c}
,,Kowalu Ilmarinenie, & ,,Ohoh seppo Ilmarinen! \\
Rzecz zrobiłeś niedorzeczną! & Jo nyt laait joutavia! \\
Złoto jak miesiąc nie świeci, & Ei kumota kulta kuuna, \\
Srebro nie grzeje jak słońce.'' & paista päivänä hopea.'' \\
\end{longtable}
}
\vspace*{\baselineskip}
\section{Joukahainen}
\vspace*{0.5\baselineskip}
\subsection*{Zaklinacz}
\begin{spacing}{1.3}
Joukahainen jest pochodzącym z~Laponii młody zaklinaczem, którego cechuje pycha i~zbytnia pewność siebie. Wymieniony na samym początku \emph{Kalevali}, podobnie do Väinämöinena dzierży nuty i~zaklęcia (I:~29--36):
\end{spacing}
{\centering
\small
\begin{longtable}{c c}
{[...]} & {[...]}\\
Owe nuty wydobyte & noita saamia sanoja, \\
I zaklęcia nanizane & virsiä virittämiä \\
{[...]} & {[...]}\\
Z Joukahainena cięciwy & Joukahaisen jousen tiestä, \\
Na gołoborzach Północy, & Pohjan peltojen periltä, \\
Wrzosowiskach Kalevali. & Kalevalan kankahilta. \\
\end{longtable}
}
\begin{spacing}{1.3}
%\noindent
O~byciu zaklinaczem (oraz o~chęci utrzymania pozycji mistrza w~tej dziedzinie) świadczą słowa (III:~21--60):
\end{spacing}
{\centering
\small
\begin{longtable}{c c}
Joukahainen, chłopak młody, & Olipa nuori Joukahainen, \\
Z Laponii dryblas żylasty, & laiha poika lappalainen. \\
Zabłądziwszy raz do wioski & Se kävi kylässä kerran; \\
Usłyszał przedziwne słowa, & kuuli kummia sanoja, \\
Co sposobniej brzmią w pamięci, & lauluja laeltavaksi, \\
Do zaklęć są przydatniejsze & parempia pantavaksi \\
{[...]}  &  {[...]}\\
Od tych, jakich sam używał, & kuin mitä itseki tiesi, \\
Pouczony przez rodzica. & oli oppinut isolta. \\
Rozeźliło go to srodze, & Tuo tuosta kovin pahastui, \\
Zazdrość go nie opuszczała, & kaiken aikansa kaehti \\
Gdy pojął, że Väinämöinen & Väinämöistä laulajaksi \\
Lepszym odeń zaklinaczem; & paremmaksi itseänsä. \\
{[...]}  &  {[...]}\\
Skoro zechcę się z nim zmierzyć, & Jos tahon tasalle panna, \\
Iść o lepsze jak mąż z mężem, & miesten verroille vetäitä, \\
Sam zaklęć swoich użyję, & itse laulan laulajani, \\
Zamawiania swe wypowiem, & sanelen sanelijani: \\
Najlepszego piewcę pieśni & laulan laulajan parahan \\
W najlichszego przeinaczę, & pahimmaksi laulajaksi, \\
{[...]}  &  {[...]}\\
\end{longtable}
}
\begin{spacing}{1.3}

Joukahainen mści się na Väinämöinenie za przegraną w~pojedynku poprzez przeszycie go strzałą, której zaklęcie kierujące potwierdza związek Lapończyka z~magią i~szamanizmem (pierwsza z~wypuszczonych strzał dociera do nieba, druga poprzez piaski do krainy śmierci, dopiero trzecia dosięga pieśniarza)  (VI:~156--162):
\end{spacing}
\newpage
{\centering
\small
\begin{longtable}{c c}
,,Teraz tnij, brzozowy grocie, & ,,Iske nyt, koivuinen sakara, \\
Przeszyj, sosnowa gałęzi, & petäjäinen selkä, lyö'ös, \\
Zadzwoń, napięta cięciwo! & jänne liina, lippaellos! \\
Jeśli obsunie się ręka, & Min käsi alentanehe, \\
Ty puść górą strzałę w locie, & sen nuoli ylentäköhön; \\
Jeżeli ręka zgóruje, & min käsi ylentänehe, \\
Niechaj lot obniży strzała!'' & sen nuoli alentakohon!'' \\
\end{longtable}
}
\vspace*{0.5\baselineskip}
\subsection*{Idiotyzm i kłamstwa}
\begin{spacing}{1.3}
Nie wykazuje się Joukahainen zbytnim rozsądkiem, przez co przegrywa pojedynek zaklinaczy. Jego głupotę wytyka mu Väinämöinen już podczas ich pierwszego spotkania (III:~102--104):
\end{spacing}
{\centering
\small
\begin{longtable}{c c}
,,Powiedz, z czyjegoś ty rodu, & ,,Kuit' olet sinä sukua, \\
Że tak głupio przesz do przodu & kun tulit tuhmasti etehen, \\
Nie pytający o drogę? &  vastahan varattomasti? \\
\end{longtable}
}
\begin{spacing}{1.3}

Fakt, iż Lapończyk wie, kogo spotkał, a~mimo wszystko pyta się zaklinacza o~tożsamość i~traktuje lekceważąco, wskazuje na jego bezbrzeżną głupotę (przez którą okrywa się on później hańbą i~będzie zmuszony oddać swoją siostrę w~ramach okupu). Wyzywa on butnie Väinämöinena na pojedynek (III:~111--114):
\end{spacing}
{\centering
\small
\begin{longtable}{c c}
,,Jestem młody Joukahainen, & ,,Mie olen nuori Joukahainen. \\
Ale ty daj znać o sobie, & Vaan sano oma sukusi: \\
Coś za jeden i skąd rodem, & kuit' olet sinä sukua, \\
Z czyjej gromady, hołyszu?'' & kuta, kurja, joukkioa?'' \\
{[...]} & {[...]} \\
,,A cóż znaczy młodość męża, & ,,Vähä on miehen nuoruuesta, \\
Albo też jego starszeństwo? & nuoruuesta, vanhuuesta! \\
Który z nas więcej ma wiedzy & Kumpi on tieolta parempi, \\
I w porę ją sobie wspomni, & muistannalta mahtavampi, \\
Do tego droga należy, & sep' on tiellä seisokahan, \\
A drugi niech się usunie. & toinen tieltä siirtykähän. \\
Możesz wszak, Väinämöinenie, & Lienet vanha Väinämöinen, \\
Wieszczu wiedzący od wieka & laulaja iän-ikuinen, \\
W śpiewaniu ze mną się zmierzyć, & ruvetkamme laulamahan, \\
W zaklinaniu pójść w zawody, & saakamme sanelemahan, \\
Mąż niechaj poucza męża, & mies on miestä oppimahan, \\
Niechaj jeden z nich zwycięża.'' & toinen toista voittamahan!'' \\
\end{longtable}
}
\begin{spacing}{1.3}

%\noindent
Na dodatek grozi, że zamieni przeciwnika w~świnię lub kiernoza, jeśli ten nie stanie do walki. W~obliczu poniesionej sromotnej klęski wygląda to śmiesznie i~świadczy o tym, iż jest Joukahainen głupcem o~zbyt wybujałym ego. Jest on również kłamcą, o~czym świadczą jego argumenty podczas pojedynku (III:~217--227):
\end{spacing}
\newpage
{\centering
\small
\begin{longtable}{c c}
Bo wspominam sobie czasy, & Muistanpa ajan mokoman, \\
Gdy zaorywałem morze, & kun olin merta kyntämässä,\\
Wytyczałem brzeg motyką, & meren kolkot kuokkimassa,\\
Dla ryb wydrążałem tonie, & kalahauat kaivamassa,\\
{[...]} & {[...]}\\
Byłem jako szósty z mężów, & ,,Viel' olin miesnä kuuentena, \\
Siódmy między mocarzami & seitsemäntenä urosna\\
Przy urządzaniu tej ziemi & tätä maata saataessa,\\
{[...]} & {[...]}\\
\end{longtable}
}


\vspace*{\baselineskip}
\section{Ilmarinen}
\vspace*{0.5\baselineskip}
\subsection*{Więź z bóstwem}
\begin{spacing}{1.3}
Ilmarinen, wielki kowal Kalevali, jest postacią bogobojną --- modli się zarówno do najwyższego Ukko (XVIII:~417--436), jak i~do pomniejszych bogiń, np. do pani mgieł i~majaków (XIX:~135--142):
\end{spacing}
{\centering
\small
\begin{longtable}{c c}
Wybierając się na łowy & Kävi siitä suistamahan; \\
Takie wypowie zaklęcie: & itse noin sanoiksi virkki: \\
,,Terhenetär, mgieł dziewojo, & ,,Ututyttö Terhenetär! \\
Porozsiewaj mgły przez sito, & Seulo seulalla utua, \\
Porozpylaj swe opary & terhenistä tepsuttele \\
W ostojach leśnej zwierzyny, & viljan vierimäsijoille, \\
By nie słyszała, jak idę, & jottei kuule kulkevaksi \\
Nie uciekała ode mnie!'' & eik' on eestäni pakene!'' \\
\end{longtable}
}
\begin{spacing}{1.3}

%\noindent
Fakt, iż wykuł on nieboskłon, skłania czytelnika do refleksji na temat jest ambiwalentnego znaczenia --- zarówno dla nieba, jak i~dla ziemi.

\vspace*{0.5\baselineskip}
\subsection*{Sacrofanum}

Jako wielbiciel kobiecego piękna, również Ilmarinen ma na swoim koncie porwanie niewiasty, a~w~przeciwieństwie do Väinämöinena posuwa się nawet o~krok dalej (XXXVIII:~113--124):

\end{spacing}
{\centering
\small
\begin{longtable}{c c}
{[...]} & {[...]} \\
Zaciął usta, schylił głowę, & murti suuta, väänti päätä, \\
Smoliste zagryzł kędziory. & murti mustoa haventa. \\
Porwał pannę w mig na ręce, & Saautti tytön samassa, \\
Mocarną dłonią zagarnął, & käärälti käpälihinsä, \\
Wyleciał jak wicher z chaty, & läksi tuiskuna tuvasta, \\
Rzucił się ku swoim saniom & riepsahti rekensä luoksi; \\ \newpage
I do sań usadził pannę, & työnnälti tytön rekehen, \\
Wepchnął do swego wasąga, & koksahutti korjahansa. \\
W jednej chwili w drogę ruszył, & Läksi kohta kulkemahan, \\
Już gotowy do podróży, & valmistui vaeltamahan, \\
Z ręką na lejcach rumaka, & käsi ohjassa orosen, \\
Z drugą na piersi dziewczyny. & toinen neien nännisillä. \\
\end{longtable}
%\end{center}
}
\vspace*{0.5\baselineskip}
\subsection*{Magia}
\begin{spacing}{1.3}
Kuźnia Ilmarinena jest domem dla magicznych zaklęć (I:~29--36). On sam, obecny przy stworzeniu świata, narodził się w~nadprzyrodzonych, magicznych  okolicznościach (IX:~100--116):
\end{spacing}
{\centering
\small
\begin{longtable}{c c}
Po bagnisku wilk przebiegał, & ,,Susi juoksi suota myöten, \\
Błądził niedźwiedź z wrzosowiska, & karhu kangasta samosi; \\
Bagno brnęło śladem wilka, & suo liikkui suen jälessä, \\
Piach pod łapami niedźwiedzia, & kangas karhun kämmenissä: \\
Tak przyszła ruda żelaza, & siihen nousi rautaruoste \\
A z nią hartowanie stali, & ja kasvoi teräskaranko \\
Za odciskiem wilczej łapy, & suen sorkkien sijoille, \\
I legowiskiem niedźwiedzia. & karhun kannan kaivamille. \\
I tak kowal Ilmarinen & ,,Syntyi seppo Ilmarinen, \\
Narodził się i wychował & sekä syntyi jotta kasvoi. \\
Na wzgórzu, gdzie tlił się węgiel, & Se syntyi sysimäellä, \\
W miejscu, gdzie go wypalano, & kasvoi hiilikankahalla \\
Ze spiżowym młotem w dłoni, & vaskinen vasara käessä, \\
Z obcęgami w drugiej ręce. & pihet pikkuiset piossa. \\
Urodził się kowal nocą, & ,,Yöllä syntyi Ilmarinen, \\
Za dnia wybudował kuźnię, & päivällä pajasen laati. \\
Szukał miejsca pod budowę, & Etsi paikkoa pajalle, \\
Gdzie by tu ustawić miechy, & levitystä lietsimille. \\
{[...]} & {[...]} \\
\end{longtable}
}
\begin{spacing}{1.3}
%\noindent
Jest również do pewnego stopnia Prometeuszem w~stosunku do ludzi, gdyż nakazał żelazu, by było posłuszne człowiekowi (IX:~177--192):
\end{spacing}
{\centering
\small
\begin{longtable}{c c}
Więc zaklęło się żelazo, & ,,Siinä vannoi rauta raukka, \\
Poprzysięgło uroczyście & vannoi vaikean valansa \\
Na kowadło, palenisko, & ahjolla, alasimella, \\
Na młot i żar pałający. & vasaroilla, valkkamilla; \\
Te wypowiedziało słowa, & sanovi sanalla tuolla, \\
Takie wyrzecze zaklęcie: & lausui tuolla lausehella: \\
,,Jestem do rąbania drzewa, & 'Onpa puuta purrakseni, \\
Do gryzienia serc kamienia, & kiven syäntä syöäkseni, \\
Lecz nie do rąbania brata, & etten veistä veikkoani, \\
Cięcia dzieci mojej matki; & lastua emoni lasta. \\
Lepiej byłoby mi zostać, & Parempi on ollakseni, \\
Zacniej, gdybym miało przeżyć & eleäkseni ehompi \\ \newpage
Jako towarzysz wędrowca, & kulkijalla kumppalina, \\
Sprzęt do ręki idącego, & käyvällä käsiasenna, \\
Niźli własny ród pożerać & kuin syöä omaa sukua, \\
Krzywdzić swoje własne plemię.'' & heimoani herjaella.'' \\
\end{longtable}
}
\vspace*{1.5\baselineskip}
\begin{spacing}{1.3}
Największym osiągnięciem panującego nad żywiołem ognia Ilmarinena jest wykucie cudownego młynka Sampo, swoistego rogu obfitości Północy, o~który wybucha później wielka wojna Väinöli z~Pohjolą (X:~403--416):
\end{spacing}
{\centering
\small
\begin{longtable}{c c}
Tedy kowal Ilmarinen & Se on seppo Ilmarinen \\
Na ostatek dnia trzeciego & päivän kolmannen perästä \\
Nachylił się, by popatrzeć, & kallistihe katsomahan \\
Co tam na dnie paleniska: & ahjonsa alaista puolta: \\
Ujrzał, jak powstaje Sampo. & näki Sammon syntyväksi, \\
W górę wieczko pstre wyrasta. & kirjokannen kasvavaksi. \\
Tedy kowal Ilmarinen, & Siitä seppo Ilmarinen, \\
Majster sławny z dawien dawna, & takoja iän-ikuinen, \\
Jął wykuwać jeszcze mocniej, & takoa taputtelevi, \\
Zawzięciej młotkiem łomotać. & lyöä lynnähyttelevi. \\
Kował Sampo umiejętnie: & Takoi Sammon taitavasti: \\
Żeby z jednej mełło mąkę, & laitahan on jauhomyllyn, \\
Sól zasię z sąsiedniej strony, & toisehen on suolamyllyn, \\
Z trzeciej wreszcie grosz gotowy. & rahamyllyn kolmantehen. \\
\end{longtable}
}
\begin{spacing}{1.3}
Ilmarinen dokonuje też wielu pomniejszych magicznych sztuk, gdy stara się o~rękę córki pani na Pohjoli (XIX:~91--94, 143--146):
\end{spacing}
{\centering
\small
\begin{longtable}{c c}
Tak zorał żmijowe pole, & Siitä kynti kyisen pellon, \\
Zbruździł ziemię pełną gadu, & vakoeli maan matoisen, \\
Pługiem wypłoszył gadziny, & nosti kyitä kynnökselle, \\
Węże oraniem rozwłóczył, & käärmehiä käännökselle. \\
{[...]} & {[...]} \\
Związał wilczy pysk kagańcem, & Sai sutosen suitsi-suuhun, \\
Niedźwiedzia w żelazne pęta, & karhun rautakahlehesen \\
Na wrzosowiskach Tuoni, & tuolta Tuonen kankahalta, \\
Wśród modrego matecznika, & sinisen salon sisästä. \\
{[...]} & {[...]} \\
\end{longtable}
}
\begin{spacing}{1.3}
\vspace*{0.5\baselineskip}
\subsection*{Deus ex machina}

Gdy do Ilmarinena przybywa Louhi (a~kowal odpowiada na pytanie czarownicy o~to, czym się aktualnie zajmuje) dochodzi do diametralnej zmiany całej sytuacji --- żądna zemsty czarownica zmienia się w~potulnego gołąbka, oddając jednocześnie światu przywłaszczone sobie wcześniej słońce i~księżyc (XLIX:~351--378).
\vspace*{\baselineskip}
\newpage
\subsection*{Oszustwa}

Ilmarinen dopuszcza się drobnych oszustw podczas przechodzenia prób, które mają udowodnić, że jest godnym mężem dla panny z~Pohjoli. Jego wybranka --- czarownica --- pomaga mu m.in. w~schwytaniu podziemnego wilka (XIX:~121--128):
\end{spacing}
{\centering
\small
\begin{longtable}{c c}
,,Kowalu Ilmarinenie, & ,,Ohoh seppo Ilmarinen, \\
Majstrze sławny z dawien dawna, & takoja iän-ikuinen! \\
Wykowaj stalowe wodze, & Teräksestä tehkös suitset, \\
Kantar przysposób z żelaza & päitset rauasta rakenna \\
Na jednym wodnym kamieniu, & yhellä vesikivellä, \\
W trzech spienionych wodospadach, & kolmen kosken kuohumilla! \\
Spętaj nim niedźwiedzia Tuoni, & Niillä tuonet Tuonen karhut, \\
W wodzach wiedź wilka z Manali.''  & suistanet suet Manalan.'' \\
\end{longtable}
}
\begin{spacing}{1.3}
\vspace*{0.5\baselineskip}
\subsection*{Idiotyzm}

Zdarza się Ilmarinenowi nie grzeszyć bystrością --- przecież Väinämöinenowi udaje się nawet skłonić go do wdrapania się na drzewo (po czym wiatr porywa kowala i~niesie w~kierunku Północy), jak również do wybrania się z~nim na groteskowy połów jednej ryby. W~runie trzydziestym dziewiątym wielki kowal ani myśli wziąć się do wiosłowania, ustępując miejsca młodzieży, starcom i~niewiastom, a~w~runie siódmym wspaniały twórca nieboskłonu zostaje podczas obróbki żelaza oszukany przez owada, co wygląda jeszcze bardziej idiotycznie (VII:~231--244):
\end{spacing}
{\centering
\small
\begin{longtable}{c c}
Szerszeń, kusego latawiec, & ,,Herhiläinen, Hiien lintu, \\
{[...]} & {[...]}\\
Polatując w mig nawrzucał & ,,Lenteä hyrähtelevi; \\
Diabelskich obrzydliwości, & viskoi Hiien hirmuloita, \\
Naznosił jadu od żmii, & kantoi käärmehen kähyjä, \\
Czarnego śluzu robaka, & maon mustia mujuja, \\
Kwasu od czerwonych mrówek, & kusiaisen kutkelmoita, \\
Utajonej żabiej złości & sammakon salavihoja \\
Do sporządzenia żelaza, & teräksenteko-mujuihin, \\
Do wyrobu modrej stali. & rauankarkaisu-vetehen. \\
\end{longtable}
}
\begin{spacing}{1.3}
Gdy zaś nad Kalevalą zalegają mroźne ciemności, kuje Ilmarinen złoty księżyc i~srebrne słońce, czego brak logiki zostaje mu natychmiast wytknięty przez Väinämöinena (XLIX:~49--58):
\end{spacing}
{\centering
\small
\begin{longtable}{c c}
,,Księżyc wykuwam ze złota, & ,,Kuuta kullaista kuvoan, \\
Ze srebra sklepuję słońce, & hope'ista aurinkoa \\
By je wznieść nad czubkiem nieba, & tuonne taivahan laelle, \\
Nad sześciorgiem pstrych pokrywek.'' & päälle kuuen kirjokannen.'' \\
Tedy stary Väinämöinen & Silloin vanha Väinämöinen \\
Takie słowa wypowiada: & itse tuon sanoiksi virkki: \\ %\newpage
\end{longtable}
}
{\centering
\small
\begin{longtable}{c c}
,,Kowalu Ilmarinenie, & ,,Ohoh seppo Ilmarinen! \\
Rzecz zrobiłeś niedorzeczną! & Jo nyt laait joutavia! \\
Złoto jak miesiąc nie świeci, & Ei kumota kulta kuuna, \\
Srebro nie grzeje jak słońce.'' & paista päivänä hopea.'' \\
\end{longtable}
}
\begin{spacing}{1.3}
%\noindent
Ilmarinen jest nieświadomy swojej głupoty i~dopiero, gdy stawia na swoim --- instalując swoje dzieła --- dostrzega, że napracował się daremnie (XLIX:~69--74):
\end{spacing}
{\centering
\small
\begin{longtable}{c c}
Wysoko umieścił księżyc, & Saipa kuun kohotetuksi, \\
Jak najwyżej słońce wsadził, & auringon asetetuksi, \\
Księżyc na wierzchołku jodły, & kuun on kuusen kukkuralle, \\
Słoneczko na szczycie sosny: & päivyen petäjän päähän: \\
Jednak nie zajaśniał miesiąc, & eipä kuu kumotakana \\
Ani słońce nie ogrzało. & eikä päivyt paistakana. \\
\end{longtable}
}
\begin{spacing}{1.3}
Boski kowal w~niezwykły sposób łączy kunszt z~potężną dawką kretynizmu, czego apogeum osiąga, gdy wykuwa sobie nową żonę ze złota (XXXVII:~39--60, 139--162):
\end{spacing}
{\centering
\small
\begin{longtable}{c c}
{[...]} & {[...]}\\
Nazbierał na morzu złota, & poimi kultia mereltä, \\
Srebra zgromadził na fali, & hope'ita lainehilta; \\
W sągach równo poukładał & keräsi kekosen puita, \\
Drzewa sań ponad trzydzieści, & kolmekymmentä rekoista; \\
W węglarni węgiel wyżarzył & puunsa poltti hiililöiksi, \\
I rozłożył w palenisku. & hiilet ahjohon ajeli. \\
Po czym wziął zebrane złoto, & Otti noita kultiansa, \\
Srebra pozbierał okruchy, & valitsi hope'itansa \\
Były jak jesienna jarka, & sykysyisen uuhen verran, \\
Spore jak zając zimowy, & verran talvisen jäniksen. \\
{[...]} & {[...]} \\
Zasię kowal Ilmarinen & Itse seppo Ilmarinen \\
Krzątał się przy palenisku, & ahjoa kohentelevi, \\
Starał się o posąg złoty, & pyyti kullaista kuvaista, \\
Srebrnolitą pannę młodą. & hope'ista morsianta. \\
{[...]} & {[...]} \\
Panna wyszła z paleniska, & Neiti ahjosta ajaikse, \\
Srebrnolica, złotowłosa, & kultaletti lietsehestä, \\
Ze złotymi warkoczami & pää hopea, kassa kulta, \\
I postacią urodziwą: & varsi kaikki kaunokainen. \\
Inni się jej przestraszyli, & Muut tuota pahoin pelästyi, \\
Lecz się nie zląkł Ilmarinen. & ei pelästy Ilmarinen. \\
Tedy kowal Ilmarinen & Siitä seppo Ilmarinen \\
Kuć jął posąg pozłocisty, & takoi kullaista kuvoa, \\
Nocą kuł bez odpocznienia, & takoi yön levähtämättä, \\
Za dnia bez chwili wytchnienia; & päivän pouahuttamatta. \\
Aż ukształtował panience & Jalat laati neitoselle, \\
Obydwie nogi i ręce, & jalat laati, käet kuvasi: \\
Jednak nie niosły jej nogi, & eipä jalka nousekana, \\
Ręce uścisnąć nie mogły. & käänny käet syleilemähän. \\
{[...]} & {[...]} \\
,,Piękna byłaby to panna, & ,,Oisi tuo sorea neito, \\
Gdyby mogła być z niej żona, & kun oisi sanallisena, \\
Duchem, mową obdarzona.'' & mielellisnä, kielellisnä.'' \\
\end{longtable}
}
\begin{spacing}{1.3}
%\noindent
Nie dość, że z~opóźnieniem orientuje się, że martwy posąg nie zastąpi prawdziwego człowieka, to dodatkowo pojawia się kolejny poważny problem (XXXVII:~175--196):
\end{spacing}
{\centering
\small
\begin{longtable}{c c}
Gdy wykąpał się do woli, & Kylpi seppo kyllitellen, \\
Wypluskał za wszystkie czasy, & valelihe vallotellen. \\
Ułożył się obok panny & Neien vierehen venähti \\
Na cienkiej jak mgła pościeli, & utuisehen uutimehen, \\
W łożu pod stalowym stropem, & teltahan teräksisehen, \\
W łożu żelazem zwieńczonym. & rankisehen rautaisehen. \\
I tam kowal Ilmarinen & Siinä seppo Ilmarinen \\
Już od razu pierwszej nocy & heti yönä ensimäisnä \\
Jął dopraszać się okrycia, & kyllä peitettä kysyvi, \\
Chustami się zabezpieczył, & vaippoja varustelevi, \\
Dwa, trzy futra wziął niedźwiedzie, & kahet, kolmet karhuntaljat, \\
Pięć, sześć zebrał chust wełnianych, & viiet, kuuet villavaipat, \\
Żeby usnąć przy małżonce, & maata kera puolisonsa, \\
Przy swym posągu ze złota. & tuon on kultaisen kuvansa. \\
Ten bok ciepły mu pozostał, & Se oli kylki kyllä lämmin, \\
Którym dotykał okrycia; & ku oli vasten vaippojansa; \\
Ten zaś, którym z panną młodą, & ku oli nuorta neittä vasten, \\
Z posągiem złotym się stykał, & vasten kullaista kuvoa, \\
Miał calutki wyziębiony, & se oli kylki kylmimässä, \\
Jak gołoledź wystudzony, & oli hyyksi hyytymässä, \\
Jak kra w morzu zamrożony & meren jääksi jäätymässä, \\
I niczym kamień stwardniały. & kiveksi kovoamassa. \\
\end{longtable}
}
\begin{spacing}{1.3}
%\noindent
Boski kowal wiezie swoją złotą narzeczoną do Väinämöinena (XXXVII: 206--210), który słusznie krytykuje dzieło rzemieślnika i przestrzega kolejne pokolenie przed tworzeniem podobnych złotych straszydeł.

Warto zwrócić uwagę na scenę, w~której Ilmarinen, widząc wodnego potwora Iku-Turso, jest śmiertelnie przerażony (XLII:~401--406):
\end{spacing}
{\centering
\small
\begin{longtable}{c c}
{[...]} & {[...]} \\
Krew mu nagle zbiegła z twarzy, & veret vieri kasvoiltansa, \\
Policzki zbladły rumiane; & puna poskilta putosi. \\
Głowę zasłonił wojłokiem, & Veti viltin päänsä päälle, \\
Po same uszy się schował, & yli korvien kohenti, \\
Nakrył twarz tak, jak potrafił, & peitti kasvot kaunihisti, \\
Jak mógł, zacisnął powieki. & silmänsä sitäi paremmin. \\
\end{longtable}
}
\begin{spacing}{1.3}
%\noindent
Zważywszy na fakt, że ów stwór okazuje się małym zwierzęciem pokroju ośmiornicy, którego Väinämöinen łapie jedną ręką i~wrzuca z~powrotem do wody, jest to kolejny dowód na tricksterską niemądrość Ilmarinena.
\end{spacing}
\newpage
\section{Louhi}
\vspace*{0.5\baselineskip}
\subsection*{Ambiwalencja}
\begin{spacing}{1.3}
Louhi wydaje się postacią negatywną --- wszakże to ona zamyka cudowny młynek Sampo w~kamiennej górze Pohjoli (XIII:~423--432), oszukuje Ilmarinena i~nie zgadza się na jego małżeństwo ze swoją córką. To ona pragnie zemścić się na całej Kalevali, a~o~pomoc prosi bezpośrednio samego Ukko (XLV:~15--22):
\end{spacing}
{\centering
\small
\begin{longtable}{c c}
,,Ukko, Boże najważniejszy! & ,,Oi Ukko, ylijumala! \\
Wygub naród Kalevali, & Kaataos Kalevan kansa \\
Wybij go żelaznym gradem, & rakehilla rautaisilla, \\
Igłami z ostrzem stalowym, & neuloilla teräsnenillä! \\
Albo zabij go zarazą, & Tahikka tauilla tapata, \\
Uśmierć plemię niegodziwe, & surmoa suku katala, \\
Mężów w rozległych zagrodach, & miehet pitkille pihoille, \\
Niewiasty na deskach obór!'' &  naiset läävän lattioille!'' \\
\end{longtable}
}
\begin{spacing}{1.3}
%\noindent
To Louhi przyczynia się do sprowadzenia Loviatar z~Tuoneli, żeby ta mogła wydać na świat choroby (XLV:~147--184):
\end{spacing}
{\centering
\small
\begin{longtable}{c c}
Tedy owa zła zakała, & Siinä tuo paha pahennus, \\
Śleporoda córa Tuoni, & Tuonen tyttö umpisilmä \\
Ulżyła swemu brzuchowi, & jopa vatsansa vajenti, \\
Powiła plemię złośliwe & latoi lapsensa vihaiset \\
Pod pstro wyszywaną płachtą, & alla vaipan vaskikirjan, \\
Pod płótnianką cieniusieńką. & alla uutimen utuisen. \\
Zrodziła dziewiątkę dziatwy & Teki poikoa yheksän \\
W ciągu jednej nocy letniej, & yhtenä kesäisnä yönä, \\
Za jednym zagrzaniem pary, & yhen löylyn lyötävillä, \\
Jednym sauny zgotowaniem, & yhen saunan saatavilla, \\
Jednym wysileniem brzucha & yhestä vatsan väestä, \\
Z żywota przeciążonego. & kohuntäyestä kovasta. \\
Ponazywała swych synów, & Nimitteli poikiansa, \\
Dała dziatwie swej przezwiska, & laaitteli lapsiansa, \\
Wedle zadań dla każdego, & kuin kuki tekemiänsä, \\
Co któremu przeznaczone: & itse ilmi luomiansa: \\
Który kłuł, przezwała kolką, & minkä pisti pistokseksi, \\
Który dusił, dychawicą, & kunka änkäsi ähyksi, \\
Gośćcem tego, który łamał, & minkä laati luuvaloksi, \\
Który ściskał --- przepukliną, & kunka riieksi risasi; \\
Tego, co ropiał --- czyrakiem, & minkä painoi paiseheksi, \\
Tego, co zasychał --- strupem, & kunka ruohutti ruveksi, \\
Tego, co pożerał --- rakiem, & minkä syöjäksi sysäsi, \\
Tego, co powalał --- dżumą. & kunka ruhtosi rutoksi. \\
Jeden został bez imienia, & Jäi yksi nimittämättä, \\
Synalek najpośledniejszy, & poika pahnan-pohjimmainen. \\
Zaraz jako guślarzowi & Senpä sitte käski tuonne, \\
Kazała mu iść nad wody, & työnti velhoiksi vesille, \\ \newpage
By zauroczył parowy, & noi'iksi noroperille, \\
By swą zawiść wszędzie wsączył. & katehiksi kaikin paikoin. \\
{[...]} & {[...]} \\
Podjudziła wściekłe stwory, & Ärrytti äkäiset luomat, \\
Pchnęła przeróżne choroby & tavattomat tauit työnti \\
Przeciw ludowi Väinöli, & vasten Väinölän väkeä, \\
Na zgubę rodu Kalevy. & surmaksi su'un Kalevan. \\
\end{longtable}
}
\begin{spacing}{1.3}
%\noindent
Warto zwrócić uwagę na to, iż w~celu ułatwienia porodu Louhi przygotowuje Loviatar saunę, czyli paradoksalnie stosuje tą samą metodę co Väinämöinen, który choroby ma później pokonać.

Ma jednakże Louhi i~dobre uczynki na kartach eposu. To właśnie ona przygotowuje dla (prawie) całego ludu Północy wspaniałe przyjęcie (XX:~609--614):
\end{spacing}
{\centering
\small
\begin{longtable}{c c}
Zaprosiła lud z Północy & Kutsui kaiken Pohjan kansan \\
I wszystko plemię Kalevy, & ja kaiken Kalevan kansan, \\
Nadto chudych komorników, & nuotki hoikat huonemiehet, \\
Kozaków w ciasnych kaftanach, & kaitakauhtanat kasakat. \\
Jeno junaka Ahtiego & Yks' on aino Ahti poika, \\
W zaprosinach pominęła. & --- senp' on heitti kutsumatta. \\
\end{longtable}
}
\begin{spacing}{1.3}
%\noindent
Czarownica chwali również swojego zięcia (XXI:~189--218):
\end{spacing}
{\centering
\small
\begin{longtable}{c c}
{[...]} & {[...]} \\
Przynieś podpałkę brzozową, & Tuopa tulta tuohen päässä, \\
Szczypkę smolnego łuczywa, & temmo tervaksen nenässä \\
Bym zięciowi się przyjrzała, & katsellakseni vävyä, \\
Popatrzeć mu mogła w oczy, & nähäkseni sulhon silmät, \\
Modre ma, czy może piwne, & sinisetkö vai punaiset \\
Albo niczym len białawe!'' & vaiko vaatevalkeuiset!'' \\
{[...]} & {[...]} \\
,,Widzę oczy mego zięcia, & ,,Jo näen vävyni silmät: \\
Ani modre, ani piwne, & ei siniset, ei punaiset \\
Ani jako len białawe, & eikä vaatevalkeuiset; \\
Białe są jak piana w morzu, & meren on vaahen valkeuiset, \\
Brunatne niby morszczyny, & meren ruo'on ruskeuiset, \\
Piękne jak morskie sitowie. & meren kaislan kauneuiset. \\
\end{longtable}
}
\begin{spacing}{1.3}
Najbardziej jednak tricksterskim fragmentem dotyczącym Louhi jest ten, w~którym posyła ona niedźwiedzia na zgubę Kalevali (XLVI:~17--20). Wielka czarownica pragnie w~akcie zemsty dokonać spustoszenia krainy, ale zamiast tego mieszkańcy łapią drapieżnika i~mają wkrótce potem wspaniałą stypę po niedźwiedziu. Dowodzi to faktu, iż trickster wykonuje paradoksalnie dobry uczynek, postępując źle.
\newpage
\subsection*{Shapeshifting}

W~runie czterdziestym trzecim Louhi zmienia diametralnie swoją postać, gdy wyrusza na wojnę z~Väinämöinenem (XLIII:~145--160, 173--180):
\end{spacing}
{\centering
\small
\begin{longtable}{c c}
,,Jaką by tu znaleźć radę, & ,,Mikäs neuvoksi tulevi, \\
U kogo szukać pomocy?'' & kukas pannahan etehen?'' \\
Wnet przybrała inną postać, & Jopa muiksi muutaltihe, \\
Przemienić się odważyła: & tohti toisiksi ruveta. \\
Pięć kos naprędce znalazła, & Otti viisi viikatetta, \\
Po sześciu motykach szczątki, & kuusi kuokan kuolioa: \\
Z nich skleciła sobie szpony, & nepä kynsiksi kyhäsi, \\
Jak pazury przytwierdziła; & kohenteli kouriksensa; \\
Połowę strzaskanej łodzi & puolen purtta särkynyttä: \\
Jak zbroję sobie nasadzi, & senpä allensa asetti; \\
Burty przypnie jako skrzydła, & laiat siiviksi sivalti, \\
Jako ptasi ogon stery, & peräpuikon purstoksensa; \\
I stu mężów pod skrzydłami, & sata miestä siiven alle, \\
Tysiąc na czubku ogona, & tuhat purston tutkaimehen, \\
Stu mężów z mieczami w ręku, & sata miestä miekallista, \\
Tysiąc z łukami mocarzy. & tuhat ampujaurosta. \\
{[...]} & {[...]} \\
Stary mądry Väinämöinen & Vaka vanha Väinämöinen \\
Odwrócił od słońca głowę, & käänti päätä päivän alta, \\
Na północo-wschód popatrzy, & luopi silmät luotehesen, \\
Rzuci okiem poza siebie; & katsoi taaksensa vähäisen: \\
Leci tam starka z Pohjoli, & jo tulevi Pohjan eukko, \\
Frunie potworne ptaszysko. & lintu kumma liitelevi, \\
Skrzydła ma jak u sokoła, & harte'ista kuin havukka, \\
Kadłub do gryfa podobny! & vaakalintu vartalolta! \\
\end{longtable}
}
\begin{spacing}{1.3}
Gdy w~końcu czarownica dopada zaklinacza, mamy do czynienia z~prawdziwym starciem tricksterów --- oto skrzydlata Louhi atakuje łódź Väinämöinena, wbija się w~nią pazurami, po czym sadowi na maszcie. Wydarzenie to odegrało wielką rolę w~fińskiej kulturze, a~samą scenę (nazwaną \emph{Obroną Sampo}) utrwalił na płótnie Akseli Gallén-Kallela.
\vspace*{0.5\baselineskip}
\subsection*{Magia}

Louhi jest czarownicą. Jej magiczne zdolności objawiają się m.in. podczas wielkiej biesiady, kiedy to uwarzone przez nią samą wcześniej piwo (XX:~425--434) zaczęło wygłaszać własną opinię na temat braku śpiewaków na przyjęciu (XX:~551--556):
\end{spacing}
\newpage
{\centering
\small
\begin{longtable}{c c}
,,Skoro nie macie śpiewaka, & ,,Kun et laita laulajata, \\
Pieśniarza jako się godzi, & laaullista laulajata, \\
Co by umiał kukać ładnie, & kunnollista kukkujata, \\
Co by pieśń zanucił składnie, & kaunista karehtijata, \\
Rozsadzę swoje obręcze, & potkin poikki vanteheni, \\
Powybijam dno w mej beczce!'' & ulos pohjani porotan!'' \\
\end{longtable}
}
\begin{spacing}{1.3}
\vspace*{0.5\baselineskip}
\subsection*{Deus ex machina}

Gdy Väinämöinen z~Ilmarinenem przybywają do Pohjoli, Louhi mówi, że Sampo dobrze spełnia swoją rolę w~kraju północnym (XLII:~45--49):
\end{spacing}
{\centering
\small
\begin{longtable}{c c}
,,Jarząbka na dwa nie dzielą, & ,,Ei pyyssä kahen jakoa, \\
Ni pomiędzy trzech wiewiórkę; & oravassa miehen kolmen. \\
Dobrze nam Sampo stukoce, & Hyvä on sampuen hyrätä, \\
Pstrokate wieczko kołace & kirjokannen kahnatella \\
W kamiennej górze w Pohjoli, & Pohjolan kivimäessä, \\
{[...]} & {[...]} \\
\end{longtable}
}
\begin{spacing}{1.3}
%\noindent
Kiedy jednak zaklinacz oświadcza jej (równie łagodnie), że zabiorą oni młynek ze sobą, Louhi wpada w~szał i~w~mgnieniu oka zwołuje całą Pohjolę na zgubę Väinämöinena. Podobną nagłą zmianę sytuacji widzimy, gdy Louhi przylatuje do Ilmarinena i~pyta, czym się aktualnie zajmuje. Po usłyszeniu odpowiedzi --- groźby prometejskiego uwięzienia --- zmienia się ona w~gołąbka pokoju, a~skradzione przez nią słońce i~księżyc wracają natychmiast na nieboskłon (XLIX:~351--378):
\end{spacing}
{\centering
\small
\begin{longtable}{c c}
,,Kuję obrożę na szyję & ,,Taon kaularenkahaista \\
Dla starej jędzy z Pohjoli, & tuolle Pohjolan akalle, \\
Jak zostanie uwiązana & jolla kiinni kytketähän \\
U podnóża litej skały.'' & vaaran vankan liepehesen.'' \\
Louhi, pani na Pohjoli, & Louhi, Pohjolan emäntä, \\
Starka Północy szczerbata, & Pohjan akka harvahammas, \\
Poczuła grożącą zgubę, & jo tunsi tuhon tulevan, \\
Dzień ucisku, co ją czeka, & hätäpäivän päälle saavan. \\
Wnet zerwała się do lotu, & Heti loihe lentämähän, \\
Chcąc ujść cało do Pohjoli. & pääsi poies Pohjolahan. \\
Uwolni księżyc z kamienia, & Laski kuun kivestä irti, \\
Słońce ze skał oswobodzi, & päästi päivän kalliosta. \\
Sama w inną się odmieni, & Itse muuksi muutaltihe, \\
W gołąbka się przedzierzgnęła, & kyhäisihe kyyhkyseksi; \\
Lekkim lotem przyfrunęła & lenteä lekuttelevi \\
Do kuźni Ilmarinena, & sepon Ilmarin pajahan. \\
Jako ptak pod drzwiami siędzie, & Lenti lintuna ovelle, \\
Jako gołąbek na progu. & kyyhkysenä kynnykselle. \\
{[...]} & {[...]} \\ \newpage
,,Takiej mogę tu na progu & ,,Tuota lienen kynnyksellä \\
Udzielić ci odpowiedzi: & sanomata saattamassa: \\
Księżyc wzeszedł już z kamienia, & jopa kuu kivestä nousi, \\
Dzień wydobył się ze skały.'' & päivä pääsi kalliosta.'' \\
\end{longtable}
}
\begin{spacing}{1.3}
\vspace*{0.5\baselineskip}
\subsection*{Kłamstwo}

Louhi ucieka się do kłamstwa, gdy matka Lemminkäinena przybywa do Pohjoli. Zapytana, co z~nim zrobiła, Louhi odpowiada (XV:~69--78):
\end{spacing}
{\centering
\small
\begin{longtable}{c c}
{[...]} & {[...]} \\
Może utonął w przeręblu, & oisko uhkuhun uponnut, \\
Zamarzł między krą na morzu, & meren jäälle jähmettynyt \\
Albo wpadł do paszczy wilka, & vai saanut sutosen suuhun, \\
W straszną paszczękę niedźwiedzia?'' & karhun kauhean kitahan?'' \\
{[...]} & {[...]} \\
,,Kłamiesz, kumo, nadaremnie! & ,,Jo vainen valehtelitki! \\
Wilk nie tknie mojego rodu, & Susi ei syö minun sukua, \\
Niedźwiedź nie ruszy junaka, & karhu ei kaa'a Lemminkäistä: \\
Co w palcach dusił wilczęta, & sormin sortavi sutoset, \\
Dłońmi dławić mógł niedźwiedzie. & käsin karhut kaatelevi. \\
\end{longtable}
}
\begin{spacing}{1.3}
%\noindent
Dopiero po wytknięciu jej kłamstwa i~zagrożeniu śmiercią, Louhi wyjawia prawdę. Później, podczas bitwy o~Sampo, Louhi paradoksalnie wytknie Lemminkäinenowi kłamstwo w~stosunku do matki (miał przecież nie pójść na wojnę przez dziesięć lat). Gdy zaś przybył do niej Ilmarinen, groteskowo kłamie Louhi na temat swojej córki (XIX:~17--32):
\end{spacing}
{\centering
\small
\begin{longtable}{c c}
,,Z panną wynikły przeszkody, & ,,Vaiv' on suuri valvatissa, \\
Z ubiorem nie daje rady: & vaiva valvateltavassa: \\
Jedna noga nie obuta, & jalk' on kesken kenkimistä, \\
Z drugą także nie nadąża. & toinen vieläki kesempi. \\
Przygotuje się dopiero, & Äsken on valmis valvattisi, \\
By cię przyjąć należycie, & oike'in otettavasi, \\
Gdy zorzesz żmijowe pole, & kun sa kynnät kyisen pellon, \\
Zabronujesz zagon węży & käärmehisen käännättelet \\
Tak, by pług się nie poruszał, & ilman auran astumatta, \\
Nie zadygotała brona; & vaarnojen värisemättä. \\
Hiisi chyba je zaorał, & Senpä Hiisi ennen kynti, \\
Diasek bruździł je przed tobą & Lempo varsinki vakoili \\
Sochą z miedzianą grądzielą, & vaarnasilla vaskisilla, \\
Pługiem z lemieszem ze stali, & auralla tuliterällä; \\
Mój zasię synek niebogi & oma poikani poloinen \\
W połowie rzucił robotę.'' & heitti kesken kyntämättä.'' \\
\end{longtable}
}
\vspace*{\baselineskip}
\section{Lemminkäinen}
\vspace*{0.5\baselineskip}
\subsection*{Więź z bóstwem}
\begin{spacing}{1.3}
Lemminkäinen (nazywany również Ahtim oraz Kaukomielim) jest pod wieloma względami powiązany z~bóstwami. Przede wszystkim --- podobnie jak reszta bohaterów \emph{Kalevali} --- zwraca się do Ukko w~potrzebie (XII:~273--288):
\end{spacing}
{\centering
\small
\begin{longtable}{c c}
Kiedy to nie dopomoże, & ,,Kun ei tuosta kyllä liene, \\
Znam ja jeszcze inny sposób, & vielä muistan muunki keinon: \\
Wyżej westchnienia skieruję, & ylemmäksi huokoaime \\
Tam, do Ukko na niebiosach, & tuolle taivahan Ukolle, \\
Który podtrzymuje chmury, & joka pilviä pitävi, \\
Panuje nad obłokami. & hattaroita hallitsevi. \\
O, Ukko, najwyższy Boże, & ,,Oi Ukko, ylijumala, \\
Dziadku stary na niebiosach, & taatto vanha taivahinen, \\
Co przemawiasz spoza chmury, & puhki pilvien puhuja, \\
Skroś powietrze się odzywasz! & halki ilman haastelija! \\
Przynieś do mnie miecz ognisty & Tuo mulle tulinen miekka \\
W płomienistej pochwie skryty, & tulisen tupen sisässä, \\
Którym nieszczęściu przeszkodzę, & jolla haittoja hajotan, \\
Szyki złoczyńców roztrącę, & jolla riitsin rikkehiä, \\
Wrogów na ziemi poskromię, & kaa'an maalliset katehet, \\
Wodnych wieszczków przezwyciężę, & ve'elliset velhot voitan \\
{[...]} & {[...]} \\
\end{longtable}
}
\begin{spacing}{1.3}
Elementy żywiołów, powiązane z~bóstwami dostrzegamy, gdy Kaukomieli wypowiada się na temat swojego pochodzenia (XI:~277--284):
\end{spacing}
{\centering
\small
\begin{longtable}{c c}
Jeślim rodu nieznacznego & ,,Jos en oo su'ulta suuri \\
I z niewysokiego domu, & enkä korkea ko'ilta, \\
Wszak mam miecz z ognistą klingą, & on mulla tulinen miekka, \\
Krzeszące iskry żelazo; & säkenevä säilärauta. \\
On jest moim możnym rodem, & Se onpi sukua suurta, \\
Moją rozległą rodziną; & laajoa lajipereä: \\
W domu diaska hartowany & onp' on Hiiessä hiottu, \\
I u bogów wyostrzony; & jumaloissa kirkastettu. \\
{[...]} & {[...]} \\
\end{longtable}
}
\begin{spacing}{1.3}
W kwestii więzi z~bóstwami Lemminkäinen charakteryzuje się jako jedyny podobieństwem do trickstera Lokiego. Niczym nordycki szachraj przybywa Ahti na biesiadę w~Pohjoli (rzecz jasna spóźniony), po czym wstępuje na środek sali i~ostro krytykuje zabawę (czego omal nie przypłaca życiem) (XXVII:~15--96):
\end{spacing}
{\centering
\small
\begin{longtable}{c c}
Lekkomyślny Lemminkäinen, & Tuop' on lieto Lemminkäinen, \\
Chłopak dziarski i niecnota, & poika, veitikkä verevä, \\
Skoro jeno wszedł do chaty, & heti kun tuli tupahan, \\
Na sam środek izby wkroczył, & astui keskilattialle: \\
Zadrży lipowa podłoga, & silta liekkui lehmuksinen, \\
Zadudni izba jodłowa. & tupa kuusinen kumahti. \\
{[...]} & {[...]} \\
,,Licho niech przychodzi tutaj & ,,Lempo tänne lähteköhön \\
I pod drzwiami się sadowi, & oven suuhun seisomahan, \\
Żeby kopeć wasz wycierać, & nokianne nuohomahan,\\
Oskrobywać wasze sadze! & karstoja karistamahan!\\
{[...]} & {[...]} \\
Czemuż nie ma tego dla mnie, & Miksip' ei ole minulle \\
Co ojcu przysługiwało?'' & kuin ennen minun isolle?'' \\
{[...]} & {[...]} \\
Och ty, włodarko Pohjoli, & ,,Oi sie Pohjolan emäntä, \\
Długozęba gospodyni! & Pimentolan pitkähammas! \\
Marnie sprawiłaś wesele, & Pi'it häät häjyn tavalla,\\
Liche gody zgotowałaś, & kutsut koiran kunnialla.\\
{[...]} & {[...]} \\
\end{longtable}
}
\vspace*{0.5\baselineskip}
\subsection*{Ambiwalencja}
\begin{spacing}{1.3}
Kaukomieli egzystuje na pograniczach wielu dziedzin życia. Najlepiej widoczną jest granica pomiędzy życiem a~śmiercią, na której Ahti lawiruje (XIII:~149--160):
\end{spacing}
{\centering
\small
\begin{longtable}{c c}
Objeżdżał góry, doliny, & Liukui linnat, liukui lannat, \\
Jakie są za morzem kraje, & liukui maat meren-takaiset; \\
Wszystkie Hiisiego ostępy, & hiihti kaikki Hiien korvet, \\
Wszystkie Kostuchy pustacie, & kaikki Kalman kankahatki, \\
Pod paszczęką Surmy-Śmierci, & hiihti Surman suun e'etse, \\
Pod płotem Kalmy-Kostuchy; & Kalman kartanon perätse. \\
Śmiertka paszczę już rozwarła, & Surma jo suutansa avavi, \\
Pochyliła łeb Kostucha, & Kalma päätä kallistavi \\
Żeby przydybać junaka, & ottoaksensa urosta, \\
Pochłonąć Lemminkäinena, & nielläksensä Lemminkäistä: \\
Nie przymierzyła dokładnie, & ei tarkoin tavannutkana, \\
Przegapiła zdobycz snadnie. & ennättänyt ensinkänä. \\
\end{longtable}
}
\begin{spacing}{1.3}
%\noindent
W końcu jednak Lemminkäinen żegna się z~tym światem (do czego przykłada rękę jego głupota oraz lekkomyślność) (XIV:~413--460), a~ratować go z~Tuoneli musi jego własna matka.

Drugie pogranicze to jego niezdecydowanie co do pojęcia dobra i~zła. Widzimy wyraźnie, jak nie może jednoznacznie określić, na jakie zwierzę aktualnie poluje (XIV: 304--312):
\end{spacing}
{\centering
\small
\begin{longtable}{c c}
,,O, Ukko, najwyższy Boże, & ,,Oi Ukko ylijumala, \\
Ukko, który dzierżysz chmury & Ukko, pilvien pitäjä, \\
I kierujesz obłokami! & hattarojen hallitsija! \\
Otwórz upusty niebieskie & Taivas auoksi avaos, \\
I wszystkie okna w przestworzu, & ilma kaikki ikkunoiksi! \\
Z góry spuść żelazne grady, & Sa'a rautaiset rakehet, \\
Krupy lodu zamrożone & laske jäiset jäähyttimet \\
Na grzywę dobrego konia, & harjalle hyvän hevosen, \\
Kłęby źrebca diabelskiego!'' & Hiien laukin lautasille!'' \\
\end{longtable}
}
\newpage
\begin{spacing}{1.3}
Trzecie skrzyżowanie dotyczy głupoty oraz mądrości. Gdy Lemminkäinen znajduje pióra cietrzewi i~zabiera je ze sobą w~myśl zasady ,,wszystko może się przydać'', sytuacja jest raczej komiczna. Niedługo jednak dostrzegamy w~niej metodę, ponieważ właśnie te niepozorne przedmioty przyczyniają się do jego zwycięstwa nad orłem (XXVI:~485--498):
\end{spacing}
{\centering
\small
\begin{longtable}{c c}
Sięgnął do swojej kieszeni, & Tapasip' on taskuhunsa, \\
W sakiewce swojej poszukał, & löihe kukkaroisehensa, \\
Znalazł tam pióra cietrzewi & otti teiren sulkasia; \\
I roztarł je na okruchy & hieroa hitustelevi \\
Pomiędzy obiema dłońmi, & kahen kämmenen välissä, \\
Między palców dziesięciorgiem, & sormen kymmenen sovussa: \\
Wyszło stąd cietrzewi stadko, & siitä syntyi teirikarja, \\
Głuszców zrodziła się zgraja, & koko parvi koppeloita. \\
Zatkał nimi dziób orłowi, & Syöksi ne kokolle suuhun, \\
Wsadził w paszczękę żarłoka, & antoi appajan kitahan, \\
W gardziel orła ognistego, & kulkkuhun kokon tulisen, \\
Iskry sypiącego ptaka; & ikenihin iskulinnun. \\
I tak sobie z nim poradził, & Sillä siitä selvittihe, \\
Przeżył pierwszy dzień podróży. & pääsi päivän ensimäisen. \\
\end{longtable}
}
\vspace*{0.5\baselineskip}
\subsection*{Shapeshifting}
\begin{spacing}{1.3}
Lemminkäinen potrafi zmienić swoją postać --- daje temu dowód, kiedy pod postacią orła ucieka z~wyspy, na której wszyscy chcą go ukatrupić (za jego miłosne ekscesy z~dziewczętami) (XXVIII:~27--34):
\end{spacing}
{\centering
\small
\begin{longtable}{c c}
Tedy letki Lemminkäinen, & Tuossa lieto Lemminkäisen, \\
Ahti z Wyspy lekkomyślny, & tuon on Ahti Saarelaisen, \\
Musiał odmienić oblicze, & täytyi toisiksi ruveta, \\
Przybrać zgoła inną postać: & piti muiksi muutellaita. \\
Jako orzeł wzbił się w górę. & Kokkona ylös kohosi, \\
Unieść chciał się wzwyż, do nieba, & tahtoi nousta taivahalle: \\
Słońce paliło policzki, & päivä poltti poskipäitä, \\
Miesiączek oświecał skronie. & kuuhut kulmia valaisi. \\
\end{longtable}
}
\vspace*{0.5\baselineskip}
\subsection*{Magia}
\begin{spacing}{1.3}
Ahti jest zaklinaczem i~czarownikiem, a~jakość jego czarowania jest raczej przyzwoita (w~przeciwieństwie do skompromitowanego Joukahainena) (XII:~445--450):
\end{spacing}
{\centering
\small
\begin{longtable}{c c}
{[...]} & {[...]} \\
Wyśpiewywać jął zaklęcia, & jop' on loihe loitsijaksi, \\
Zabrał się do czarowania; & laikahtihe laulajaksi: \\
Ogniem mienił mu się kożuch & tulta iski turkin helmat, \\
I skry z oczu się sypały & valoi silmät valkeata \\
Od śpiewu Lemminkäinena, & Lemminkäisen laulaessa, \\
Od jego zaklęć i czarów. & laulaessa, lausiessa. \\
Piewców najlepszych przedzierzgnął & Lauloi laulajat parahat \\
W początkujących śpiewaków, & pahimmiksi laulajiksi; \\
Zatkał usta kamieniami, & kivet suuhun syrjin syösti, \\
Głazów powpychał w gardziele & paaet lappehin lateli \\
Najznakomitszym śpiewakom, & parahille laulajille, \\
Wieszczom najumiejętniejszym. & taitavimmille runoille. \\
\end{longtable}
}
\begin{spacing}{1.3}
Opowiada również Lemminkäinen o~rodowodzie swoich mocy magicznych --- guślarzem stał się na własne życzenie, co czyni go podobnym Väinämöinenowi (który w~tym temacie poszedł nawet dalej i~sam zdecydował o~swoich narodzinach) (XII: 144--174):
\end{spacing}
{\centering
\small
\begin{longtable}{c c}
,,Uroczyli mnie guślarze, & ,,Jo minua noiat noitui, \\
Wiedźmy, żmije mnie jadziły, & noiat noitui, kyyt kiroili; \\
Trzech Lapończyków mnie jednej & koki kolme lappalaista \\
Letniej nocy próbowało & yhtenä kesäisnä yönä, \\
Nagiego na nagiej skale & alasti alakivellä, \\
Bez pasa, bez przyodziewy, & ilman vyöttä, vaattehitta, \\
Ruszyć się nie mogącego; & rikorihman kiertämättä: \\
Tyle ode mnie zyskali, & senpä hyötyivät minusta, \\
To dostało się niezgułom, & sen verran, katalat, saivat, \\
Ile siekierze od skały, & min kirves kivestä saapi, \\
Lub świdrowi od kamienia, & napakaira kalliosta, \\
Zadymce od gołoledzi, & järky jäästä iljanesta, \\
Od pustej izby Kostusze. & Tuoni tyhjästä tuvasta. \\
I drugi raz mi grożono, & ,,Toisinpa oli uhattu, \\
Powtórnie czyhano na mnie, & toisinpa kävi kätehen. \\
W topiel wessać zamierzali, & Mielivät minua panna, \\
{[...]} & {[...]} \\
Ale taki człowiek ze mnie, & Vaan minäpä, mies mokoma, \\
Żem się zbytnio nie przestraszył, & en tuossa kovin hätäillyt; \\
Sam się stałem zaklinaczem, & itse loime loitsijaksi, \\
Przedzierzgnąłem się w guślarza, & sain itse sanelijaksi: \\
Urzekłem żerców czarami, & lauloin noiat nuolinensa, \\
Strzelców razem z ich strzałami, & ampujat asehinensa, \\
Wróżów z nożami z żelaza, & velhot veitsirautoinensa, \\
Zaklinaczy z ich grotami & tietäjät teräksinensä \\
W stromą siklawę Tuoni, & Tuonen koskehen kovahan, \\
W wodny wir przerażający & kinahmehen kauheahan, \\
{[...]} & {[...]} \\
\end{longtable}
}
\begin{spacing}{1.3}
Ahti potrafi również wyczarowywać jarzębiny, dęby i~kukułki eksplodujące złotem, srebrem i~miedzią (por.~XXIX), jak również zaklinać wodę (XL:~23--32):
\end{spacing}
{\centering
\small
\begin{longtable}{c c}
,,Przestań pienić się, siklawo, & ,,Heitä, koski, kuohuminen, \\
Skończ wirować, wodo wrąca! & vesi vankka, vellominen! \\
Córo progu, panno piany, & Kosken tyttö, kuohuneiti! \\
Usiądź na skwierczącej skale, & Istuite kihokivelle, \\
Na głazie pośród bulgotu; & kihopaaelle paneite! \\ \newpage
W objęcia pochwyć grzywacze, & Sylin aaltoja aseta, \\
Wychlusty okrąż rękoma, & käsin kääri käppyröitä, \\
Garściami zagarniaj piany, & kourin kuohuja kohenna, \\
By nam nie bryzgały w piersi, & jottei riusko rinnoillemme \\
Naszych głów nie zagubiły! & eikä päällemme päräjä! \\
\end{longtable}
}
\vspace*{0.5\baselineskip}
\subsection*{Sacrofanum}
\begin{spacing}{1.3}
Podobnie jak inni bohaterowie \emph{Kalevali}, Lemminkäinen ma wielką słabość do kobiet (XI:~15--20):
\end{spacing}
{\centering
\small
\begin{longtable}{c c}
{[...]} & {[...]} \\
Wszak nieznaczną miał przywarę, & vaan tuli vähän vialle, \\
Nie nazbyt chlubne zwyczaje: & tavoiltansa turmiolle: \\
Biegał za białogłowami, & ain' oli naisissa eläjä, \\
Noc w noc z nimi miał zajęcie, & yli öitä öitsilöissä, \\
Gdziekolwiek się zabawiały, & noien impien iloissa, \\
Kołyszące warkoczami. & kassapäien karkeloissa. \\
\end{longtable}
}
\begin{spacing}{1.3}
%\noindent
Dość swobodnie folguje on swoim potrzebom seksualnym (XI:~150--156):
\end{spacing}
{\centering
\small
\begin{longtable}{c c}
{[...]} & {[...]}\\
Ówże gładki Kaukomieli, & itse kaunis Kaukomieli, \\
Przepłoszył kobiece kpiny, & jopa hääti naisen naurun, \\
Uciszył żeńskie języki, & piätteli piian pilkan. \\
Nie było tutaj córeczki, & Ei ollut sitä tytärtä, \\
Panny nawet najcnotliwszej, & piikoa pyhintäkänä, \\
By zostawił ją nietkniętą, & kuta hän ei kosketellut, \\
Żeby jej nie towarzyszył. & jonk' ei vieressä venynyt. \\
\end{longtable}
}
\begin{spacing}{1.3}
Tricksteryzmu Kaukomielego dowodzi również fakt, iż przekonał on swojego kompana Tierę do wyprawy na wojnę, chociaż ten nie zaznał jeszcze rozkoszy nocy poślubnej, co dobitnie komentuje jego rodzina  (XXX:~85--92):
\end{spacing}
{\centering
\small
\begin{longtable}{c c}
,,Nie ruszy na wojnę Tiera, & ,,Ei Tiera sotahan joua, \\
Z łomem swym na bijatykę, & Tieran tuura tappelohon! \\
Dopiero co targu dobił, & Tiera on tehnyt kuulun kaupan, \\
Zawarł umowę na wieki, & ikikaupan iskenynnä: \\
Tiera pannę dziś poślubił, & vast' on nainut naisen nuoren, \\
Znalazł sobie gospodynią, & ottanut oman emännän; \\
Jeszcze draźniąt jej nie tykał, & viel' on nännit näppimättä, \\
Piersi młodej nie utrudził.'' & rinnat riuahuttamatta.'' \\
\end{longtable}
}
%\vspace*{\baselineskip}
\newpage
\subsection*{Deus ex machina}
\begin{spacing}{1.3}
Lemminkäinen dokonuje gwałtownej zmiany sytuacji m.in. wtedy, gdy porywa swoją wybrankę Kylikki (XI:~203--206):
\end{spacing}
%\newpage
{\centering
\small
\begin{longtable}{c c}
{[...]} & {[...]}\\
Porwał Kylikki na sanie, & reutoi Kyllikin rekehen, \\
Zgarnął pannę do wasąga, & koppoi neien korjahansa, \\
Owinął w futro niedźwiedzie, & tuon asetti taljallensa, \\
I na dnie sań usadowił. & liitti liistehyisillensä. \\
\end{longtable}
}
\begin{spacing}{1.3}
Ze względu na swoją niestałość w~zamierzeniach i~postanowieniach \emph{deus ex machina} jest widoczny również, gdy Ahti oświadcza nagle, że wybiera się na wojnę (XII:~57--64):
\end{spacing}
{\centering
\small
\begin{longtable}{c c}
,,Ja nie wierzę w sny niewieście & ,,En usko unia naisten \\
Ani w białogłów zaklęcia; & enkä vaimojen valoja. \\
Matko, któraś mnie nosiła, & Oi emoni, kantajani! \\
Przynieś mi moją kolczugę, & Tuo tänne sotisopani, \\
Mój rynsztunek wojownika! & kanna vainovaatteheni! \\
Ogarnęła mnie ochota & Mieleni minun tekevi \\
Wojennego wypić trunku, & juomahan soan olutta, \\
Miodu wojny zakosztować.'' & soan mettä maistamahan.'' \\
\end{longtable}
}
\begin{spacing}{1.3}
%\noindent
Do nagłych zmian sytuacji należy również przyporządkować fakt, iż jest to jedyna postać w~\emph{Kalevali} która zmartwychwstaje --- dzięki pomocy swojej matki, która przeszukuje ziemię i~niebiosa w~poszukiwaniu szczątków syna na kształt Izydy żeglującej za swoim martwym mężem Ozyrysem.

Lemminkäinen nagle zmienia sytuację również wtedy, gdy udaje mu się pojmać diabelskiego źrebca --- nie przemocą, lecz łagodnością (XIV:~327--342):
\end{spacing}

{\centering
\small
\begin{longtable}{c c}
,,Dobry koniu z diablej ziemi, & ,,Hiitolan hyvä hevonen, \\
Źrebcze z gór pokryty pianą, & vuoren varsa vaahtileuka! \\
Daj teraz swój pysk złocisty, & Tuo nyt kultaturpoasi, \\
Wsadź srebrzystą swoją głowę & pistä päätäsi hopea \\
Pomiędzy tręzelki złote, & kultaisihin koltuskoihin, \\
Pomiędzy srebrne dzwoneczki. & hope'isihin helyihin! \\
Źle nie obejdę się z tobą, & En sua pahoin pitäne, \\
Nie obciążę ponad miarę, & aivan anke'in ajane: \\
Pojadę drogą najkrótszą, & ajan tietä pikkuruisen, \\
Podróż potrwa nie za długo & matkoa ani vähäisen, \\
Do posępnych chat Pohjoli, & tuonne Pohjolan tuville, \\
Do mojej świekry surowej; & ankaran anopin luoksi. \\
Po co miałbym bić rzemieniem & Minkä siimalla sivallan \\
Albo biczem cię częstować, & eli vitsalla vetelen, \\
Jeno musnę cię jedwabiem, & senpä silkillä sivallan, \\
Machnę wełną mięciuteńką.'' & veran äärellä vetelen.'' \\
\end{longtable}
}

\vspace*{0.5\baselineskip}
\subsection*{Kłamstwo}
\begin{spacing}{1.3}
Lemminkäinen ucieka się do kłamstwa, gdy Väinämöinen każe mu się rozejrzeć w~poszukiwaniu żądnych zemsty (po kradzieży Sampo) mieszkańców Pohjoli (XLIII: 45--48):
\end{spacing}
{\centering
\small
\begin{longtable}{c c}
,,Przed nami jest jasne niebo, & ,,Selvänä etinen ilma, \\
Za nami za to gęstnieje: & taakea takainen taivas: \\
Mała chmurka mknie z północy, & pieni on pilvi pohjosessa, \\
Skrawek z północo-zachodu.'' & pilven lonka luotehessa.'' \\
\end{longtable}
}
\begin{spacing}{1.3}
%\noindent
Väinämöinen wytyka mu głupotę i~oszustwo. Kiedy indziej Ahti zarzeka się, iż stronił od kobiet, chociaż romansował prawie ze wszystkimi, jakie tylko spotkał.

\vspace*{0.5\baselineskip}
\subsection*{Idiotyzm}

Fakt, iż Lemminkäinen jest głupcem, widoczny jest już w~pierwszym runie --- ukazany jest Ahti jako ten, który ,,zatarł się z~krotochwilą'' (I:~45--50). Louhi nazywa go nicponiem, któremu nie może oddać swojej córki (XIII:~24--30). Jego głupota jest powszechnie znana, chociaż on sam nie zdaje sobie z~niej sprawy (XI:~111--124):
\end{spacing}
{\centering
\small
\begin{longtable}{c c}
Drwią z Lemminkäinena panny, & Nauroi naiset Lemminkäistä, \\
Docinały mu dziewczyny, & piiat pisti pilkkojansa, \\
Gdy ulicą osobliwie & kun ajoi kummasti kujalle, \\
Jechał niby strach do dworu, & kamalasti kartanolle: \\
Kiedy, wywracając sanie, & ajoi korjansa kumohon, \\
Z rozmachem wleciał we wrota. & veräjähän vierähytti. \\
{[...]} & {[...]} \\
,,Tego nie widziałem jeszcze, & ,,En ole tuota ennen nähnyt, \\
Nie widziałem, nie słyszałem, & en ole nähnyt enkä kuullut \\
By się baby ze mnie śmiały, & naisen nauravan minulle, \\
Szydziły ze mnie dziewczyny.'' & piian pilkkoja suannut.'' \\
\end{longtable}
}
\begin{spacing}{1.3}
%\noindent
Niemądre jest również to, iż zaczarowuje on w~walce wszystkich swoich przeciwników, za wyjątkiem jednego pastucha --- zapłaci za to później najwyższą cenę (XII:~470--476):
\end{spacing}
{\centering
\small
\begin{longtable}{c c}
{[...]} & {[...]} \\
Zawróżył mężów z mieczami, & lauloi miehet miekkoinensa, \\
Zbrojną czeladź z jej orężem, & urohot asehinensa; \\
Młodych, starych zaczarował, & lauloi nuoret, lauloi vanhat, \\
Z nimi wraz tych w sile wieku; & lauloi kerran keskilaaun; \\
Jednego zakląć zaniedbał: & yhen heitti laulamatta: \\
Krów pastucha pośledniego, & karjapaimenen pahaisen, \\
Jednookiego dziadygę. & ukko vanhan umpisilmän. \\
\end{longtable}
}
\newpage
\begin{spacing}{1.3}
Idiotyczne w~przypadku Kaukomielego jest również to, iż rozpacza nad błahymi sprawami --- podczas gdy nie wypełnia próby, która miała zapewnić mu małżonkę, lamentuje nad złamanymi nartami (XIII: 259--270):
\end{spacing}
{\centering
\small
\begin{longtable}{c c}
Tedy letki Lemminkäinen & Siinä lieto Lemminkäinen \\
Ze smętnie schyloną głową & alla päin, pahoilla mielin \\
Popatrzał na swój dobytek & kalujansa katselevi. \\
I te słowa wypowiada: & Itse tuon sanoiksi virkki: \\
,,Niechaj odtąd, póki życia, & ,,Elköhön sinä ikänä \\
Żaden spośród naszych mężów & menkö toinen miehiämme \\
Nie waży się wdawać w lesie & uhalla metsän ajohon, \\
W pogoń za diabelskim łosiem, & Hiien hirven hiihäntähän, \\
Jak się wdałem ja, nieborak, & kuin menin minä, poloinen! \\
Dobre deski zmarnowałem & Hävitin hyvät sivakat, \\
Oraz kij do nart najlepszy, & sauvan kaunihin kaotin, \\
Nadto włócznię najcelniejszą!'' & kiihke'immän keihojani!'' \\
\end{longtable}
}
\begin{spacing}{1.3}
Przez głupotę Lemminkäinena dostaje się również jego kompanom --- gdyby podczas powrotu z~Pohjoli nie zaczął śpiewać (ostrzegano go przed tym nie raz), nie doszłoby do tricksterskiego starcia Väinämöinena z~Louhi, a~Sampo by ocalało (XLII: 278--292):
\end{spacing}
{\centering
\small
\begin{longtable}{c c}
{[...]} & {[...]} \\
Urodziwy Kaukomieli, & tuo on kaunis Kaukomieli, \\
Usta sobie przysposobi, & suutansa sovittelevi, \\
Nabrał tchu, nastroił gardło; & säveltänsä säätelevi. \\
Sam jeden jął wyśpiewywać, & Sai itse hyräilemähän, \\
Począł, psotnik, pohukiwać & loihe, kurja, kukkumahan \\
Swoim zachrypniętym głosem, & äreällä äänellänsä, \\
Popisywać się hałasem. & käreällä kulkullansa. \\
Śpiewał letki Lemminkäinen, & Lauloi lieto Lemminkäinen, \\
Porykiwał Kaukomieli, & karjahteli Kaukomieli; \\
Trzęsła mu się broda, wargi, & suu liikkui, järisi parta, \\
Paszczęki mu się rozwarły, & leukapielet lonkaeli. \\
Śpiew rozlegał się daleko, & Laulu kuului loitommalle, \\
Zawodzenie nad wodami & vierähys vesien poikki, \\
Do wsi szóstej doleciało & kuului kuutehen kylähän, \\
Poprzez siedem morskich toni. & seitsemän selän ylitse. \\
\end{longtable}
}
\newpage
\section{Kullervo}
\vspace*{0.5\baselineskip}
\subsection*{Więź z bóstwem}
\begin{spacing}{1.3}
Więź, jaka łączy Kullervo z~bóstwami, opiera się na modlitwach i~lamentowaniu, gdyż ten pobożny głupiec przez całe życie czuje się samotny (XXXIII:~19--46):
\end{spacing}
{\centering
\small
\begin{longtable}{c c}
,,Przyświecaj, Boże słoneczko, & ,,Paistapa, Jumalan päivä, \\
Użycz ciepła, kłębku Pana, & Herran kehrä, hellittele \\
Niebogiemu pastuszkowi & sepon karjan kaitsijalle,\\
Strzegącemu krów kowala, & poloiselle paimenelle,\\
{[...]} & {[...]} \\
Biegnij, słońce, tocz się, bochnie, & ,,Mene, päivä, viere, vehnä, \\
Prędko mijaj, Boży dzionku! & alene, Jumalan aika! \\
Skryj się, słońce, za świerczyną, & Kule, päivä, kuusikolle, \\
Za las biegnij, bochnie pszenny, & viere, vehnä, vitsikölle, \\
Za jałowczyska uciekaj, & karkoa katajikolle, \\
Zniż się na równi z olchami, & lennä leppien tasalle! \\
Puść już do domu pasterza, & Päästä paimenta kotihin \\
By mógł masła kęs ukroić, & voivatia vuolemahan, \\
Rozerwać bochenek przaśny, & rieskoa repäisemähän, \\
Wygrzebać z niego ośródkę!'' & kakkaroita kaivamahan!'' \\
\end{longtable}
}
\begin{spacing}{1.3}
\vspace*{0.5\baselineskip}
\subsection*{Ambiwalencja}

Kullervo stara się być dobrym człowiekiem (w~\emph{Kalevali} nazwany jest po śmierci nawet bohaterem, por. XXXVI:~345), jednakże w~jego przypadku sprawdza się powiedzenie, iż dobrymi chęciami wybrukowano piekło. Nieudolne postawienie płotu można uznać za brak doświadczenia, jednakże zabójstwo dziecka już nie --- widać w~tym dalekie echo zabójstwa Baldura przez Lokiego (XXXI:~213--232):
\end{spacing}
{\centering
\small
\begin{longtable}{c c}
Skoro dorósł już Kullervo, & Kun oli Kullervo kohonnut, \\
O piędź mu przybyło wzrostu, & saanut vartta vaaksan verran, \\
Posłano go do roboty, & tuopa työlle työnnetähän, \\
Zadanie mu wyznaczono, & raaolle rakennetahan, \\
Niemowlęcia miał doglądać, & lapsen pienen katsontahan, \\
Huśtać malucha-palucha: & sormi pienen souantahan: \\
,,Pięknie pilnuj niemowlęcia, & ,,Katso lasta kaunihisti, \\
Nakarm je i sam zjedz przy nim, & syötä lasta, syö itseki! \\
Przepłukuj w rzece pieluchy, & Rievut virrassa viruta, \\
Dbaj o czyste kaftaniki.'' & pese pienet vaattehuiset!'' \\
Dzień, drugi doglądał dziecka, & Katsoi lasta päivän, kaksi: \\
Rączkę urwał, wykłuł oko, & käen katkoi, silmän kaivoi. \\
Aż pod koniec dnia trzeciego & Siitä kohta kolmannella \\
Zemrzeć dziecku dał w chorobie, & lapsen tauilla tapatti, \\
Z biegiem wód puścił pieluchy, & rievut viskoi virran vieä, \\
W ogniu spalił kołyseczkę. & kätkyen tulella poltti. \\
Pomyślał sobie Untamo: & Untamo ajattelevi: \\ \newpage
,,Nie ma widocznie zacięcia, & ,,Ei tämä tähän sopiva \\
By pilnować niemowlęcia, & lapsen pienen katsontahan, \\
Huśtać malucha-palucha; & sormi pienen souantahan! \\
{[...]} & {[...]} \\
\end{longtable}
}
\begin{spacing}{1.3}
%\noindent
W~tym incydencie widać również podobieństwo Kullervo do trickstera Winnebagów, pod którego opieką znajdują się dzieci jego brata. Noszone w~pęcherzu i~nieumiejętnie karmione umierają, a~jedyną reakcją boskiego szaleńca na to nieszczęście jest sarkazm (Radin 2010:~35--36):
\end{spacing}

\begin{adjustwidth}{1,5cm}{1,5cm}
\small{
,,Po krótkiej chwili sam zgłodniał. »Te małe dzieci należy karmić raz w~miesiącu, powiedziano mi« --- pomyślał. Ale teraz sam był bardzo głodny. Wobec tego powiedział znowu: »Och, och! Chyba już pora, żeby moje kochane dzieci znowu były głodne. Muszę zdobyć dla nich coś do jedzenia«. [...] Kiedy mięso ostygło, odpiął od pasa pęcherz i~otworzył go. Ku jego zdumieniu dzieci były martwe. »Kochane małe dzieci! Cóż za nieszczęście, że umarły!«''}\\
\end{adjustwidth}
\vspace*{0.5\baselineskip}
\subsection*{Magia}
\begin{spacing}{1.3}
Kullervo od urodzenia posiada zdolności magiczne i~nadprzyrodzone, m.in. nadludzką siłę (XXXI:~87--98):
\end{spacing}
{\centering
\small
\begin{longtable}{c c}
Bujał się w kolebce dzieciak, & Liekkui lapsi kätkyessä, \\
Kołysał się, trząsł czupryną, & lapsi liekkui, tukka löyhki. \\
Kołysał się dzionek, drugi, & Liekkui päivän, liekkui toisen; \\
Dnia trzeciego na ostatek, & jopa kohta kolmantena, \\
Jak nie wierzgnie malec nogą, & kun tuo poika potkaisihe, \\
Jak kopnie, jak się natęży, & potkaisihe, ponnistihe, \\
Rozdarł swoje powijaki, & katkaisi kapalovyönsä, \\
Z koca w mig się oswobodził, & pääsi päälle peittehensä, \\
Strzaskał lipową kołyskę, & särki liekun lehmuksisen, \\
Potargał do cna pieluchy. & kaikki riepunsa revitti. \\
Stwierdzono, że zuch zeń będzie, & Nähtihin hyvä tulevan, \\
Że się nieźle zapowiada; & keksittihin kelpoavan. \\
{[...]} & {[...]} \\
\end{longtable}
}
\begin{spacing}{1.3}
Potrafi on już mówić, pałać chęcią zemsty i~unikać śmierci w~iście groteskowy sposób, panując nad naturą i~bawiąc się żywiołem (XXXI:~109--192):
\end{spacing}
{\centering
\small
\begin{longtable}{c c}
,,Jak jeno dorosnąć zdołam & ,,Kunpa saisin suuremmaksi, \\
I okrzepnie moje ciało, & vahvistuisin varreltani, \\
Tedy pomszczę rany ojca, & kostaisin isoni kohlut, \\
Matki mojej łzy odpłacę!'' & maksaisin emoni mahlat!'' \\
{[...]} & {[...]} \\
Do antałka go włożono, & Pannahanpa puolikkohon, \\
Posadzono go do beczki, & työnnetähän tynnyrihin; \\
Zaniesiono ją do wody & siitä vieähän vetehen, \\
I opuszczono na fale. & lasketahan lainehesen. \\
{[...]} & {[...]} \\
W beczce nie skonało dziecię, & kuollut poika puolikkohon! \\
Chłopak z beczki się wydostał, & Poika oli pääsnyt puolikosta \\
Siedział na grzywacza grzbiecie, & --- istui aaltojen selässä \\
W ręku wędę miał miedzianą, & vapa vaskinen käessä,\\
Linewkę przy niej jedwabną, & siima silkkinen perässä;\\
Którą ryby łowił w morzu & onkivi meren kaloja,\\
I głębokość wody mierzył: & merivettä mittoavi:\\
{[...]} & {[...]} \\
Zebrano, nagromadzono & Koottihin, keräeltihin \\
Suche brzozowe bierwiona, & koivuja, kovia puita, \\
Sosny z gałęziami stoma, & honkia satahavuja, \\
Nasyconymi żywicą, & tiettäviä tervaksia, \\
Tysiąc sań brzozowej kory, & tuohia tuhat rekeä, \\
Sto sążni drew jesionowych; & sata syltä saarnipuita. \\
Ogień pod nie podłożono, & Tuli puihin tuiskattihin, \\
Rozbiegły się w stosie iskry, & roviohon roiskattihin, \\
I wrzucono weń chłopaka & siihen poika paiskattihin \\
Do ognia gorejącego. & keskelle tulen palavan. \\
Paliło się dzień i drugi, & Paloi päivän, tuosta toisen, \\
Paliło się przez dzień trzeci, & paloi päivän kolmannenki. \\
Wybrano się, by zobaczyć; & Käytihin katsastamahan: \\
Po kolana tkwił w popiele, & poik' oli porossa polvin, \\
Po łokcie głęboko w prochu, & kypenissä kyynäsvarsin, \\
Z pogrzebaczem długim w dłoni, & hiilikoukkunen käessä, \\
Którym płomię podtrzymywał, & millä tulta kiihottavi, \\
Węgielki nim przegartywał & hiiliä kokoelevi, \\
I ni włos mu nie spadł z głowy, & katomatta karvankana, \\
Nie zawinął się kędziorek. & kutrisen kähertymättä! \\
{[...]} & {[...]} \\
Powieszono go na drzewie, & Poika puuhun hirtetähän, \\
Żeby dyndał się na dębie. & tammehen ripustetahan. \\
Dwie i trzy minęły noce, & Kului yötä kaksi, kolme, \\
I tyleż dni za koleją, & saman verran päiviäki. \\
Pomyślał sobie Untamo: & Untamo ajattelevi: \\
,,Pora byłaby popatrzeć, & ,,Aik' on käyä katsomahan, \\
Zali przepadł już Kullervo, & joko Kullervo katosi, \\
Poniósł śmierć na szubienicy.'' & kuoli poika hirsipuuhun.'' \\
Posłał sługę, by zobaczył, & Laittoi orjan katsomahan. \\
Rzecze po powrocie sługa: & Orja toi sanan takaisin: \\
,,Nie przepadł zgoła Kullervo, & ,,Ei ole Kullervo kaonnut, \\
Nie skonał na szubienicy, & kuollut poika hirsipuuhun! \\
Wyrzezał wzory na drzewie, & Poika puuta kirjoittavi \\
Mając w garści kozik mały, & pieni piikkonen käessä. \\
Całe drzewo jest we wzorki, & Koko puu kuvia täynnä, \\
Pełen cały dąb obrazków: & täynnä tammi kirjoitusta: \\
Są tu ludzie, są tu miecze, & siinä miehet, siinä miekat, \\
Włócznie mają też przy boku.'' & siinä keihä'ät sivulla.'' \\
\end{longtable}
}
\begin{spacing}{1.3}
%\noindent
Kullervo na drzewie przywołuje obraz Odyna wiszącego na Yggdrasilu, który przez swoją ofiarę miał nauczyć się magicznych pieśni i~run. Kalevaliczny szaleniec jednak żadnej mądrości nie zdobywa, chociaż w~iście magiczny sposób ratuje się od śmierci.
\newpage
Magia praktyczna w~wykonaniu Kullervo ma miejsce, gdy przemienia on wilki i~niedźwiedzie w~zwierzęta gospodarskie (XXXIII:~129--138):
\end{spacing}
{\centering
\small
\begin{longtable}{c c}
Wyciął witkę z zagajnika, & Otti vitsan viiakosta, \\
Bat na bydło jałowcowy, & katajaisen karjanruoskan; \\
Wpędził krowie stado w błoto, & sorti suohon lehmäkarjan, \\
W chruśniaku rozproszył woły, & härät murtohon murenti \\
Połowę na pastwę wilkom, & puoliksi susien syöä, \\
Niedźwiedziom z boru połowę: & puolen korven kontioille. \\
Wilki zaczarował w krowy, & Suet lausui lehmäsiksi, \\
W bydełko zmienił niedźwiedzie, & karhut karjaksi rakenti, \\
Jedne przezwał Maluśkimi, & minkä pisti Pienikiksi, \\
Drugie zasię Białuchami. & kunka Kyytäksi kyhäisi. \\
\end{longtable}
}
\begin{spacing}{1.3}
Do magicznych zdolności zalicza się w~przypadku Kullervo również groteskową rozmowę z~przedmiotem nieożywionym --- własnym mieczem (XXXVI:~319--334):
\end{spacing}
{\centering
\small
\begin{longtable}{c c}
Kullervo z Kalervy rodu & Kullervo, Kalervon poika, \\
Wyciągnął miecz wyostrzony, & tempasi terävän miekan; \\
Oglądał go i obracał, & katselevi, kääntelevi, \\
Zapytywał go i badał; & kyselevi, tietelevi. \\
Zapytał swój miecz o zdanie, & Kysyi mieltä miekaltansa, \\
Zali miałby on ochotę & tokko tuon tekisi mieli \\
Grzesznego ciała skosztować & syöä syyllistä lihoa, \\
I napić się krwi niedobrej. & viallista verta juoa. \\
Miecz rozważył myśli męża, & Miekka mietti miehen mielen, \\
Zrozumiał mowę junaka, & arvasi uron pakinan. \\
Odpowie tymi słowami: & Vastasi sanalla tuolla: \\
,,Czemuż nie miałbym ochoty & ,,Miks' en söisi mielelläni, \\
Skosztować grzesznego ciała & söisi syyllistä lihoa, \\
I napić się krwi niedobrej? & viallista verta joisi? \\
Kosztowałem ciał bezgrzesznych, & Syön lihoa syyttömänki, \\
Krwi niewinnej się napiłem.'' & juon verta viattomanki.'' \\
\end{longtable}
}
\begin{spacing}{1.3}
\vspace*{0.5\baselineskip}
\subsection*{Sacrofanum}

Podobnie jak inni bohaterowie fińskiego eposu, Kullervo ma również słabość do płci pięknej. Zaleca się on do trzech dziewcząt, a~gdy po raz trzeci uzyskuje odpowiedź odmowną, porywa niewiastę (XXXV:~141--166):
\end{spacing}
{\centering
\small
\begin{longtable}{c c}
,,Chodź do moich sanek, dziewczę, & ,,Käy, neito, rekoseheni, \\
Pod wojłoki moje, miła, & armas, alle vilttieni, \\
Jeść będziemy moje jabłka & syömähän omeniani, \\
I zagryzać orzechami!'' & puremahan päähkeniä!'' \\
Odpowiedziała dziewczyna, & Neiti vastahan sanovi, \\
Cynopierśna odburknęła: & tinarinta riuskuttavi: \\
,,Gwiżdżę, głupcze, na twe drwalki, & ,,Sylen, kehno, kelkkahasi, \\ \newpage
Na twoje sanki, hultaju! & retkale, rekosehesi! \\
Zimno jest pod wojłokami, & Vilu on olla viltin alla, \\
Obrzydliwie w twoich saniach!'' & kolkko korjassa eleä.'' \\
Kullervo, Kalervy synek, & Kullervo, Kalervon poika, \\
Czarci syn w modrej pończosze, & sinisukka äijön lapsi, \\
Porwał pannę do sań swoich, & koppoi neion korjahansa, \\
Wrzucił do swego wasąga, & reualti rekosehensa, \\
Na futrach ją usadowił, & asetteli taljoillensa, \\
Jął ją wciągać pod wojłoki. & alle viltin vieretteli. \\
Tak mu na to rzecze panna, & Neiti tuossa noin sanovi, \\
Przeciwi się cynopierśna: & tinarinta riitelevi: \\
,,Natychmiast mnie stąd wyswobódź, & ,,Päästä pois minua tästä, \\
Wolność przywróć mi, dziecinie, & laske lasta vallallensa \\
Bym słów zwodnych nie słyszała, & kunnotointa kuulemasta, \\
Nie uległa złej namowie, & pahalaista palvomasta, \\
Albo w saniach dno wybiję, & tahi potkin pohjan puhki, \\
Porozrzucam wszystkie deski, & levittelen liistehesi, \\
Porozbijam je na trzaski, & korjasi pilastehiksi, \\
Na wióry ten włok rozniosę!'' & rämäksi re'en retukan!'' \\
\end{longtable}
}
\begin{spacing}{1.3}
%\noindent
Następnie uwodzi ją i~--- po raz pierwszy w~\emph{Kalevali} --- mamy do czynienia ze sceną seksu: (XXXV:~169--188):
\end{spacing}
{\centering
\small
\begin{longtable}{c c}
{[...]} & {[...]} \\
Otwarł ze skarbami skrzynię, & aukaisi rahaisen arkun, \\
Odchylił wieczko ozdobne, & kimahutti kirjakannen; \\
Pokazywał srebro swoje, & näytteli hope'itansa, \\
Postawy sukna rozłożył, & verkaliuskoja levitti, \\
Złotem przetkane pończochy, & kultasuita sukkasia, \\
Pasy srebrem wyszywane. & vöitänsä hopeapäitä. \\
Sukna myśli panny zmącą, & Verat veivät neien mielen, \\
Skarby młódkę odmieniły, & raha muutti morsiamen, \\
Srebro ciągnie ją do zguby, & hopea hukuttelevi, \\
Wreszcie nakłoniło złoto. & kulta kuihauttelevi. \\
{[...]} & {[...]} \\
Nakłaniał ją i przywabiał, & kuihutteli, kutkutteli, \\
Z ręką na wodzach rumaka. & käsi orosen ohjaksissa, \\
Drugą na piersiach dziewczyny, & toinen neitosen nisoissa. \\
Że wreszcie dziewczynę skusił, & Siinä neitosen kisasi, \\
Cynopierśną sił pozbawił, & tinarinnan riu'utteli \\
Pod wojłokiem wyszywanym, & alla vaipan vaskikirjan, \\
Pośród cętkowanych futer. & päällä taljan taplikkaisen. \\
\end{longtable}
}
\begin{spacing}{1.3}
%\noindent
Dowiaduje się potem Kullervo, iż złamał święte tabu --- współżył z~własną siostrą. Po tym wydarzeniu płacze cała natura, a~Kullervo wpada w~ciężką depresję.

\newpage
\subsection*{Idiotyzm}
Kullervo jest największym głupcem w~całej \emph{Kalevali}. Każda praca, jakiej się podejmuje, przynosi katastrofalne skutki. Gdy ma postawić płot lub młócić ziarno, robi to na opak (XXXI:~334--348):
\end{spacing}
{\centering
\small
\begin{longtable}{c c}
,,Na nic zda się ta robota, & ,,Ei tämä tähän sopiva! \\
Nie do przejścia ten płot sklecił, & Pani aian aukottoman, \\
Pozostawił bez wierzei, & veräjättömän kyhäsi, \\
Wzniósł go do samego nieba, & tuon on nosti taivosehen, \\
Wypiętrzył ponad obłoki, & ylös pilvihin kohotti: \\
Nie przejdę przez niego górą & en tuosta ylitse pääse \\
Ani przez szparę do środka; & enkä reiästä sisälle! \\
{[...]} & {[...]} \\
Kullervo, Kalervy synek, & Kullervo, Kalervon poika, \\
Wziął się do młócenia zboża: & jo oli puimassa rukihit: \\
W proch ziarno z plewami zmłócił & pui rukihit ruumeniksi, \\
I słomę w sieczkę obrócił. & olet kaunaksi kaotti. \\
\end{longtable}
}
\begin{spacing}{1.3}
%\noindent
Jego idiotyzm jest zaś notoryczny i~nic nie zapowiada jakiejkolwiek poprawy (XXXV: 3--10):
\end{spacing}
{\centering
\small
\begin{longtable}{c c}
{[...]} & {[...]} \\
Od tej chwili mógł przebywać & sai tuosta elelemähän \\
Pod rodzicielską opieką; & alla varjon vanhempien; \\
Ale nie nabrał rozumu, & ei saanut älyämähän, \\
Nie pojmował myśli męża, & miehen mieltä ottamahan, \\
Bowiem licho był chowany, & kun oli kaltoin kasvateltu, \\
Na głupca wykołysany & tuhmin lasna tuuiteltu \\
U lichego wychowawcy & luona kalton kasvattajan, \\
I piastunki niedorzecznej. & tuon on tuhman tuuittajan. \\
\end{longtable}
}
\vspace*{\baselineskip}
\section{Ukko}
\vspace*{0.5\baselineskip}
\subsection*{Więź z bóstwem}
\begin{spacing}{1.3}

W~kalevalicznym panteonie bogów Ukko jest tym najwyższym i~najważniejszym (II:~317--320):
\end{spacing}
{\centering
\small
\begin{longtable}{c c}
Oj, Ukko, najwyższy Boże, & ,,Oi Ukko, ylijumala \\
Ojcze władający w niebie, & tahi taatto taivahinen, \\
Dżdżów i chmur mocarny sprawco & vallan pilvissä pitäjä, \\
I opiekunie obłoków! & hattarojen hallitsija! \\
\end{longtable}
}
\newpage
\begin{spacing}{1.3}
%\noindent
To on stworzył podwaliny całego świata (z~wyłączeniem nieba, które wykuł Ilmarinen) (IX:~33--46):
\end{spacing}
{\centering
\small
\begin{longtable}{c c}
Ukko, najwyższy stworzyciel, & ,,Tuo Ukko, ylinen luoja, \\
Sam bóg powietrznych przestworzy, & itse ilmojen jumala, \\
Oddzielił od niebios wodę, & ilmasta ve'en eroitti, \\
Z wody stały ląd wyłonił, & veestä maati manterehen. \\
Jeno żałosne żelazo & Rauta on raukka syntymättä, \\
Pozostało nie zrodzone. & syntymättä, kasvamatta. \\
Ukko, bóg ponad przestworzem, & ,,Ukko, ilmoinen jumala, \\
Zacierał obiedwie dłonie, & hieroi kahta kämmentänsä, \\
Zaciskał je z całej siły & mykelti molempiansa \\
Na swoim lewym kolanie; & vasemmassa polven päässä. \\
Narodziły się trzy panny, & Siitä syntyi kolme neittä, \\
Trzy stworzycielki natury, & koko kolme luonnotarta \\
Macierze rudy żelaznej, & rauan ruostehen emoiksi,\\
Modrej stali trzy sprawczynie. & suu sinervän siittäjiksi.\\
\end{longtable}
}
\vspace*{0.5\baselineskip}
\subsection*{Ambiwalencja i magia}
\begin{spacing}{1.3}
Tricksterskie połączenie dobra i~zła w~przypadku Ukko jest wyjątkowo dobrze widoczne i nie pozostawia żadnych wątpliwości. Najwyższy kalevaliczny bóg jest sprawcą wielu pomyślności (II:~331--348; IX:~403--416, 567--576):
\end{spacing}
{\centering
\small
\begin{longtable}{c c}
Ukko, najwyższy wśród bogów, & Tuo Ukko, ylijumala, \\
Ojciec władający w niebie & taatto taivon valtiainen, \\
Rozkazał zebrać się chmurom, & piti pilvissä keräjät, \\
Zgromadził je na naradę, & sekehissä neuvot selvät. \\
{[...]} & {[...]}\\
Jak najściślej wszystkie stłoczył, & lomituksin loukahutti. \\
Mżenie sączył jął z niebiosów, & Vihmoi vettä taivosesta,\\
Pokrapiać miodem z obłoków & mettä pilvistä pirotti\\
{[...]} & {[...]}\\
Aż wyrosła ruń oścista, & Nousipa oras okinen, \\
Wystrzeliła bujnowłosa & kannonkarvainen yleni \\
Z pulchnej roli, uprawionej & maasta pellon pehmeästä, \\
Mozołem Väinämöinena. & Väinämöisen raatamasta.\\
{[...]} & {[...]}\\
O Ukko, stwórco najwyższy, & ,,Oi Ukko, ylinen luoja, \\
Boże, któryś jest w niebiosach! & taivahallinen jumala! \\
Zawitaj tutaj w potrzebie, & Tule tänne tarvittaissa, \\
Przybądź tu na zaproszenie, & käy tänne kutsuttaessa! \\
Dotknij swoją pewną ręką, & Tunge turpea kätesi, \\
Naciśnij swym grubym kciukiem, & paina paksu peukalosi \\
By zatkać otwór niesforny, & tukkeheksi tuiman reiän, \\
Zawrzeć te oporne wrota, & paikaksi pahan veräjän! \\
Przynieś listek delikatny, & Veä päälle lemmen lehti, \\
Żółtego grążela płatek, & kultalumme luikahuta \\
Żeby krwi zagrodzić drogę, & veren tielle telkkimeksi, \\
Zatamować uchodzącą, & tulevalle tukkeheksi, \\
By nie tryskała na brodę, & jottei parsku parralleni, \\
Nie lała się na odzienie.'' & valu vaaterievuilleni!'' \\
{[...]} & {[...]}\\
,,Stąd zawsze łaski przychodzą, & ,,Tuoltapa aina armot käyvät, \\
Pomoc przychodzi prawdziwa & turvat tuttavat tulevat  \\
Ze sklepienia wysokiego & ylähältä taivahasta,  \\
Od Stwórcy wszechmogącego. & luota Luojan kaikkivallan.  \\
Pozdrowiony bądź, Jumala, & ,,Ole nyt kiitetty, Jumala,  \\
Ty sam, Stwórca, uwielbiony, & ylistetty, Luoja, yksin,  \\
Boś udzielił mi pomocy, & kun annoit avun minulle,  \\
Przyniosłeś pewną opiekę & tuotit turvan tuttavasti  \\
W owych boleściach okrutnych, & noissa tuskissa kovissa,  \\
Gdy mnie ostra stal zraniła!'' & terän rauan raatamissa!''  \\
\end{longtable}
}
\begin{spacing}{1.3}
Jednakże Ukko potrafi być również niszczycielski, co przejawia się w~spełnianiu przez niego próśb bohaterów, modlących się o~zesłanie z~niebios konkretnego zniszczenia. Mroczna moc Ukko może być skierowana na pojedynczą jednostkę (XIV:~313--322):
\end{spacing}
{\centering
\small
\begin{longtable}{c c}
Ukko, najwyższy stworzyciel, & Tuo Ukko, ylinen luoja, \\
Główny bóg nad obłokami, & pilven-päällinen jumala, \\
Rozdarł na poły przestworza, & ilman riehoiksi revitti, \\
Na dwoje nieba sklepienie, & taivon kannen kahtaloksi; \\
Sypnął szronem, sypnął lodem, & satoi hyytä, satoi jäätä, \\
Posypał żelaznym gradem, & satoi rauaista raetta, \\
Mniejszym niźli końska głowa, & pienemmät hevosen päätä, \\
Większym niż głowa człowieka, & päätä ihmisen isommat, \\
Na grzywę dobrego konia, & harjalle hyvän hevosen, \\
Kłęby źrebca diabelskiego. & Hiien laukin lautasille. \\
\end{longtable}
}
\begin{spacing}{1.3}
Może być ona również przeznaczona całemu narodowi, o~co prosi Louhi (XLV: 15--22), a~co skutkuje porodem Loviatar --- Ukko pomaga w~przygotowaniach do narodzenia się zła na świecie (XLV:~88--96):
\end{spacing}
{\centering
\small
\begin{longtable}{c c}
{[...]} & {[...]} \\
Stwórca z niebiosów wyrzecze: & lausui Luoja taivahalta:\\
,,Stoi szałas na trzech żerdziach & ,,Tuoll' on suolla kolmisoppi \\
Na bagnie, nad brzegiem morza, & rannalla meryttä vasten, \\
W pokrytej mrokiem Pohjoli, & pimeässä Pohjolassa, \\
Mgłami spowitej Sarioli; & sangassa Sariolassa. \\
Tam pójdziesz dziatwę porodzić, & Mene sinne poikimahan, \\
Sprawisz ulgę swemu łonu, & kohtusi keventämähän! \\
Tam cię właśnie potrzebują, & Siellä silma tarvitahan, \\
Na twą czeladkę czekają!'' & väkeäsi vuotetahan.'' \\
\end{longtable}
}
\begin{spacing}{1.3}
Ukko posiada w~sobie również pierwiastek magiczny --- potrafi rozniecić ogień w~niebie za pomocą miecza, jak również strzelać nim z~paznokci. Schowany do kapciucha ogień służy później do stworzenia nowego słońca i~nowego księżyca. W~iście tricksterski sposób wymyka się ognik Ukko i~spada na ziemię poprzez dziewięć niebios, a~następnie dokonuje niemałych zniszczeń (XLVII:~191--202):
\end{spacing}
\newpage
{\centering
\small
\begin{longtable}{c c}
{[...]} & {[...]}\\
Pannom piersi pokaleczył, & rikkoi rinnat tyttäriltä, \\
Poparzył sutki dziewczętom, & neitosilta nännit näppi, \\
Poranił chłopcu kolana, & turmeli pojalta polvet, \\
Spalił brodę gospodarza. & isännältä parran poltti. \\
{[...]} & {[...]} \\
Spalił dzieciątko w kołysce, & poltti lapsen kätkyestä, \\
Spalił piersi rodzicielce; & poltti paarmahat emolta. \\
Poszło dzieciątko do Many, & Se lapsi meni Manalle, \\
Chłopaczek do Tuoneli, & toki poika Tuonelahan, \\
{[...]} & {[...]} \\
\end{longtable}
}
\begin{spacing}{1.3}
%\noindent
To właśnie przez niezdarność bóstwa cierpią zwierzęta oraz płonie mityczne jezioro Alue. Ukko przez szlachetne zamierzenia czyni ostatecznie zło, co dowodzi jego ambiwalencji.

\vspace*{0.5\baselineskip}
\subsection*{Sacrofanum}

W~kontekście łączenia sfery \emph{sacrum} ze sferą \emph{profanum}, należy spojrzeć na Ukko podobnie jak na Väinämöinena --- obaj lubią otaczać się kobietami. Mówi o~tym żona Ilmarinena, kierując swoje słowa do Ukko (XXXII:~79--82):
\end{spacing}
{\centering
\small
\begin{longtable}{c c}
Sporo dziewcząt masz u siebie, & Paljo on piikoja sinulla, \\
Setki czekają rozkazu, & saoin käskyn kuulijoita, \\
Które pod niebem mieszkają, & eläjiä ilman alla, \\
Panny przychylne stworzeniom. & luonnottaria hyviä. \\
\end{longtable}
}
\begin{spacing}{1.3}

\vspace*{0.5\baselineskip}
\subsection*{Deus ex machina}

Ukko jako \emph{bóg z~maszyny} ukazuje się czytelnikowi, gdy spełnia prośby kalevalicznych bohaterów. To za jego pomocą nagle dokonała żywota żona Ilmarinena (na prośbę Kullervo). Wszelkie zesłania żywiołowych sił natury i~wahania pogodowe również świadczą o~jego naturze jako tego, który odwraca sytuację.

\vspace*{0.5\baselineskip}
\subsection*{Idiotyzm}

Pomimo swojej wszechpotęgi i~statusu głównego boga, Ukko bywa również głupi. Gdy Louhi zamyka słońce i~księżyc w kamiennej górze Pohjoli, ciemności zapadają zarówno nad ziemią, jak i~nad niebiosami (XLVII:~33--40). Wtedy to Ukko dość długo się zastanawia, co może być przyczyną takiego stanu rzeczy (XLVII:~47--58):
\end{spacing}
\newpage
{\centering
\small
\begin{longtable}{c c}
Nieporęcznie żyć bez ognia, & Tukela on tuletta olla, \\
Zgoła trudno bez płomienia, & vaiva suuri valkeatta, \\
Cni się bezeń człowiekowi, & ikävä inehmisien, \\
A nawet samemu Ukko. & ikävä itsen Ukonki. \\
Ukko, najwyższy wśród bogów, & Tuo Ukko, ylijumala, \\
Sam wielki nieba stworzyciel, & itse ilman suuri luoja, \\
Jął się wielce tym zdumiewać, & alkoi tuota ouostella. \\
Głowił się i zastanawiał, & Arvelee, ajattelevi, \\
Co za stwór sięgnął miesiąca, & mikä kumma kuun e'essä, \\
Co za cień na drodze słońca, & mikä terhen päivän tiessä, \\
Skoro miesiąc nie migoce, & kun ei kuu kumotakana \\
Ani słoneczko nie świeci. & eikä päivä paistakana. \\
\end{longtable}
}
\vspace*{\baselineskip}
\section{Ilmatar}
\begin{spacing}{1.3}
W~przypadku Ilmatar (zwanej również Luonnotar) tricksterska domena bliskości bóstwa ustępuje miejsca domenie bycia samym bóstwem, a~następnie udzielania potrzebującym swojej pomocy. Wielka córa powietrza jest matką siedmiu synów (sił natury), zrodziła również za sprawą morza i~wiatru zaklinacza Väinämöinena, można więc nazwać ją matką trickstera. Jest to najbardziej widoczny związek tej postaci z~boskim szelmą, który często ją przywołuje w~potrzebie (II:~101--106):
\end{spacing}
{\centering
\small
\begin{longtable}{c c}
Przeto stary Väinämöinen & Siitä vanha Väinämöinen \\
Sam ją przyzywał zaklęciem: & itse tuon sanoiksi virkki: \\
,,Panno stworzenia, boginko & ,,Kave äiti, kantajani, \\
Macierzy i opiekunko, & luonnotar, ylentäjäni! \\
Wiele w morzu masz czeladzi, & laitapa ve’en väkeä \\
Może z niej który poradzi, & --- veessä on väkeä paljo --- \\
{[...]} & {[...]}\\
\end{longtable}
}
\begin{spacing}{1.3}
Pod koniec eposu mamy do czynienia z~nietypową sytuacją, w~której Väinämöinen i~Ilmarinen spotykają Ilmatar podczas podróży łódką. Bogini zadaje im pytanie (XLVII:~145--146):
\end{spacing}
{\centering
\small
\begin{longtable}{c c}
,,Coście za jedni, mężowie, & ,,Mitä miehiä olette, \\
Jakie nosicie imiona?'' & kuinka teitä kutsutahan?'' \\
\end{longtable}
}
\begin{spacing}{1.3}
%\noindent
Sytuacja jest dość groteskowa, ponieważ wielka bogini-matka nie poznaje własnego syna, co wskazuje na to, iż do boskiej mądrości trochę jednak jej brakuje. Co ciekawe, Väinämöinen również nie rozpoznaje matki. Podczas gdy przedstawia on siebie i~swojego kompana z~imienia, Ilmatar przedstawia się następująco, podkreślając wyniośle swoje pochodzenie (XLVII:~154--158):
\end{spacing}
\newpage
{\centering
\small
\begin{longtable}{c c}
,,Spośród niewiast jam najstarszą & ,,Minä olen vanhin vaimoksia,\\
Spomiędzy panien powietrza. & vanhin ilman impilöitä,\\
Pierwszą ze wszystkich macierzą, & ensin emä itselöitä.\\
W cenie pięciu żonom równa & joll’ on vihki viien vaimon,\\
I sześciu pannom urodą; & muoto kuuen morsiamen.\\
{[...]} & {[...]} \\
\end{longtable}
}

\vspace*{\baselineskip}
\section{Marjatta}
\begin{spacing}{1.3}
Skromna dziewczyna za sprawą borówki zachodzi w~ciążę i~rodzi syna, który zaraz potem ośmiesza Väinämöinena i~zostaje ochrzczonym królem Karelii (por.~L). Przez to wydarzenie można intepretować postać Marjatty jako ambiwalentną --- sprowadza ona szczęście do Kalevali, jednocześnie (pośrednio) skazując na banicję jej odwiecznego pieśniarza. Fakt urodzenia się potomka Marjatty zahacza również o~magię, jak i~sam w~sobie działa jako \emph{deus ex machina} (tutaj również dosłownie) dla całej historii --- oto mają nadejść czasy nowej wiary (oczywista bliskość boga), chociaż Väinämöinen obiecuje, że jeszcze do Kalevali powróci --- tutaj z~kolei mamy do czynienia z~połączeniem \emph{sacrum} i~\emph{profanum} na linii wiary pogańskiej i~chrześcijańskiej. Marjatta jest zatem tricksterem ukrytym i~skoncentrowanym, ponieważ jej historii poświęcono tylko jedno runo eposu --- ostatnie.
\end{spacing}



\chapter{Podsumowanie}
\begin{spacing}{1.3}
Analiza zachowań bohaterów fińskiego eposu narodowego potwierdza występowanie u~nich cech szalbierskich w~ilości wystarczającej, by nazwać ich odbiciami prototypowego trickstera na płaszczyźnie kalevalicznej. W~sylwetkach postaci z~łatwością odnajdujemy domeny kłamstwa, ambiwalencji, więzi z~bogami, idiotyzmu, magii, nagłego odwracania sytuacji i~rozwiązywania akcji, zmiennokształtności oraz głupoty i~łączenia sakralności ze sprośnością.

Przyporządkowanie postaciom \emph{Kalevali} domen tricksterskich prowadzi do wyodrębnienia się wśród nich dwóch grup: tricksterów \emph{większych} i~tricksterów \emph{mniejszych}. Tricksterzy więksi to postacie, do których ma zastosowanie sześć lub więcej domen boskich szaleńców --- są to prototypowi tricksterzy Väinämöinen i~Lemminkäinen (posiadający po 8~domen) oraz Ilmarinen i~Ukko (po 6~domen).
\end{spacing}
\vspace{\baselineskip}
\begin{table}[!h]
\renewcommand{\arraystretch}{1.18}
\centering
\begin{center}
\begin{tabu}{l |[1.5pt] c|c|c|c} %
                        & Väinämöinen & Lemminkäinen & Ilmarinen & Ukko      \\\tabucline[1.5pt]{-}
{więź z bóstwem}        & \ding{51}   & \ding{51}    & \ding{51} & \ding{51} \\\hline
{ambiwalencja}          & \ding{51}   & \ding{51}    & ---       & \ding{51} \\\hline
{shapeshifting}         & \ding{51}   & \ding{51}    & ---       & ---       \\\hline
{magia rytualna}        & \ding{51}   & \ding{51}    & \ding{51} & \ding{51} \\\hline
{sacrofanum}            & \ding{51}   & \ding{51}    & \ding{51} & \ding{51} \\\hline
{deus ex machina}       & \ding{51}   & \ding{51}    & \ding{51} & \ding{51} \\\hline
{kłamstwo}              & \ding{51}   & \ding{51}    & \ding{51} & ---       \\\hline
{idiotyzm}              & \ding{51}   & \ding{51}    & \ding{51} & \ding{51} \\

\end{tabu}
\end{center}
\caption{Tricksterzy więksi w \emph{Kalevali}}
\end{table}
\newpage
\vspace*{3.5\baselineskip}
\begin{spacing}{1.3}
Tricksterzy mniejsi posiadają przypisanych pięć lub mniej szaleńczych domen. Są to kolejno: Louhi i~Kullervo (po 5~domen), Marjatta (4~domeny) oraz Joukahainen i~Ilmatar (po 3~domeny).
\end{spacing}
\vspace{\baselineskip}
\begin{table}[!h]
\renewcommand{\arraystretch}{1.18}
\centering
\begin{center}
\begin{tabu}{l |[1.5pt] c|c|c|c|c}
                            & Louhi       & Kullervo     & Marjatta  & Joukahainen & Ilmatar     \\\tabucline[1.5pt]{-}
{więź z bóstwem}        & ---         & \ding{51}    & \ding{51} & ---         & \ding{51}   \\\hline
{ambiwalencja}          & \ding{51}   & \ding{51}    & \ding{51} & ---         & ---         \\\hline
{shapeshifting}         & \ding{51}   & ---          & ---       & ---         & ---         \\\hline
{magia rytualna}        & \ding{51}   & \ding{51}    & \ding{51} & \ding{51}   & \ding{51}   \\\hline
{sacrofanum}            & ---         & \ding{51}    & ---       & ---         & ---         \\\hline
{deus ex machina}       & \ding{51}   & ---          & \ding{51} & ---         & ---         \\\hline
{kłamstwo}              & \ding{51}   & ---          & ---       & \ding{51}   & ---         \\\hline
{idiotyzm}              & ---         & \ding{51}    & ---       & \ding{51}   & \ding{51}   \\

\end{tabu}
\end{center}
\caption{Tricksterzy mniejsi w \emph{Kalevali}}
\end{table}
\vspace{\baselineskip}
\begin{spacing}{1.3}
Najczęściej występującą cechą tricksterską w~\emph{Kalevali} jest magia (oraz związane z~nią zaklinania) --- mają z nią do czynienia wszyscy bohaterowie. Na drugim miejscu plasują się \emph{ex aequo} domeny: więzi z~bóstwem i~idiotyzmu (siedmiu przedstawicieli). Sześciu reprezentantów ma domena nagłego odwrócenia sytuacji (\emph{deus ex machina}), jak również domena ambiwalencji (egzystowania ,,na granicy''). Kłamstwo i~sacrofanum to domeny występujące pięciokrotnie. Najrzadziej występująca u~bohaterów \emph{Kalevali} domena to shapeshifting (zmiennokształtność) --- cechą tą mogą poszczycić się jedynie: Väinämöinen, Lemminkäinen i~Louhi.
\end{spacing}
\newpage
\vspace*{3.5\baselineskip}
\begin{spacing}{1.3}
Z powyższych danych wynikają następujące wnioski:
\vspace{\baselineskip}
\begin{enumerate}[1)]   %[\ding{51}]
  \item kalevaliczny trickster para się magią, śpiewaniem lub zaklinaniem,
  \item ponad $\sfrac{3}{4}$ kalevalicznych tricksterów posiada koneksje z~bóstwem/bóstwami,
  \item ponad $\sfrac{3}{4}$ kalevalicznych tricksterów cechuje głupota (często groteskowa),
  \item dwie na trzy postacie potrafią nagle odwrócić daną sytuację,
  \item dwie na trzy postacie są ambiwalentne i żyją ,,na granicy”,
  \item ponad połowa głównych bohaterów fińskiego eposu narodowego to kłamcy,
  \item ponad połowa głównych bohaterów \emph{Kalevali} łączy świętość ze sprośnością,
  \item jedna na trzy postacie może w zależności od potrzeby przybrać inne oblicze.\\
\end{enumerate}
\end{spacing}
%\vspace{\baselineskip}
\begin{spacing}{1.3}
Przyjmując podobny margines jak w~przypadku wyodrębnienia tricksterów większych z~ogólnej puli (poprzez wystąpienie minimum 6~z~8~domen) w~przypadku częstotliwości występowania cech boskich szaleńców (czyli przyjmując, że cecha występuje u minimum $\sfrac{3}{4}$ bohaterów) możemy stwierdzić, że stereotypowy fiński trickster:

\begin{enumerate}[1)]
  \item jest magiem, szamanem, śpiewakiem lub zaklinaczem,
  \item jest powiązany z bóstwami lub sam jest bóstwem,
  \item jest idiotą.
\end{enumerate}
\end{spacing}


\clearpage{\pagestyle{empty}\cleardoublepage}
\part{PRAESTIGIATOR NOVUS KALEVALANUS}
\clearpage{\pagestyle{empty}\cleardoublepage}

\chapter{Nowy trickster kalevaliczny}
\begin{spacing}{1.3}
Tricksterzy są obecni w~mitologiach i~opowieściach całego świata. Poprzez świadkowanie różnym kulturom i~poprzez różne sposoby utrwalania opowieści dla potomnych ulegają szalbierze mniej lub bardziej znaczącym przekształceniom, jak również nabierają oni cech specyficznych, właściwych danemu regionowi czy grupie etnicznej (nie bez kozery indiańskim szaleńcem jest Zając, a~nordyckim --- olbrzym-jötunn).

W~przypadku kanwy kalevalicznej typowe cechy tricksterskie zostały już szczegółowo przedstawione. Patrząc jednak z~perspektywy całości na fiński epos narodowy, możemy w~nim dostrzec szereg powtarzających się cech charakterystycznych, które --- pomimo tego, iż nie należą stricte do portretu psychologicznego boskiego błazna --- grają znaczącą rolę w~konstruowaniu tych epickich postaci. Biorąc pod uwagę te właściwości, oraz dodając je do ogólnej puli omówionych już wcześniej tricksterskich punktów, wyłania się nam obraz nowego mistyfikatora --- \emph{nowego trickstera kalevalicznego}.
\end{spacing}
%\newpage
%\vspace{3\baselineskip}
\section{Antyprometeusz}
\begin{spacing}{1.3}
Prometeusz jest jednym z~wielu przykładów książkowego trickstera --- ambiwalentny heros, który przeciwstawia się ustanowionym przez bogów zasadom w~imię dobra ludzkości, kradnie dla nich ogień, za co zostaje następnie dotkliwie ukarany. \emph{Nowy trickster kalevaliczny} jest często odwrotnością Prometeusza --- z~własnej woli odbiera szczęście ogółu, lub nie robi niczego, żeby takiemu rozwojowi wypadków zapobiec. Za przykład może posłużyć Joukahainen, który --- nie bacząc na słowa matki --- pragnie jedynie dokonać zemsty na Väinämöinenie (VI:~135--210):
\end{spacing}
\newpage
{\centering
\small
\begin{longtable}{c c}
Tymi słowy na ostatek & Virkki viimeinki sanoiksi,\\
Odezwie się, tak wypowie: & itse lausui, noin nimesi: \\
,,Niechaj nawet dwakroć ujdzie & ,,Kaotkohot jos kahesti\\
Wszelka radość z tego świata & kaikki ilmaiset ilomme,\\
I wszystkie pieśni przepadną, & kaikki laulut langetkohot!\\
Strzelę pewnie, bez wahania.'' & Varsin ammun, en varanne.''\\
{[...]} & {[...]}\\
Wówczas młody Joukahainen & Siinä nuori Joukahainen \\
Takim chełpił się językiem: & itse kielin kerskaeli:\\
,,Już, stary Väinämöinenie, & ,,Et sinä, vanha Väinämöinen,\\
Nie zobaczysz nigdy w świecie, & enämpi elävin silmin\\
Nie ujrzysz na własne oczy & sinä ilmoisna ikänä,\\
Jasnej poświaty miesiąca & kuuna kullan valkeana\\
Na wrzosowiskach Väinöli, & astu Väinölän ahoja,\\
Na polanach Kalevali! & Kalevalan kankahia!\\
Sześć lat bujaj tu na fali & ,,Kupli nyt siellä kuusi vuotta,\\
I siedem kolejnych wiosen, & seuro seitsemän kesyttä,\\
Pokołuj sobie lat osiem & karehi kaheksan vuotta\\
Na owych wodach przestronnych, & noilla väljillä vesillä,\\
Na rozkołysanych falach. & lake'illa lainehilla:\\
Przez sześć lat jako świerczyna, & vuotta kuusi kuusipuuna,\\
Siedem lat jako choina, & seitsemän petäjäpuuna,\\
Osiem jako pień zrąbany!'' & kannon pölkkynä kaheksan!''\\
\end{longtable}
}
%\vspace{0.5\baselineskip}
\begin{spacing}{1.3}
%\noindent
Pod względem chęci zniszczenia niemal równy Joukahainenowi jest szalony Kullervo (XXXI: 279--294):
\end{spacing}
%\newpage
{\centering
\small
\begin{longtable}{c c}
,,Niech się wali las dokoła, & ,,Sini kaski kaatukahan,\\
Niechaj runą grube brzozy, & koivu solki sortukahan,\\
Dokąd mój głos się rozlega, & kuni ääni kuulunevi,\\
Gwizdanie moje dobiega! & kuni vierrevi vihellys!\\
Ani pęd niech się nie płoży, & ,,Elköhön vesa venykö,\\
Ani źdźbło się nie uniesie, & elköhön koretko korsi\\
Póki życia na tej ziemi & sinä ilmoisna ikänä,\\
I miesiączek złoty świeci, & kuuna kullan valkeana\\
Tu, gdzie rąbał syn Kalervy, & kaskessa Kalervon poian,\\
Mąż karczował znakomity! & otoksessa oivan miehen!\\
Gdyby tu jęczmień zasiano & ,,Ottaisiko maa orahan,\\
I zakiełkowało zboże, & nousisiko nuori laiho,\\
Gdyby źdźbła się zawiązały & sekä korsi korteuisi,\\
I wystrzeliły łodygi, & jotta varsi varteuisi,\\
Niech z nich kłos nie napęcznieje & elköhön tereä tehkö,\\
Ani ziarno nie dojrzeje!'' & varsi päätä valmistako!''\\
\end{longtable}
}
%\vspace{0.5\baselineskip}
%\newpage
\begin{spacing}{1.3}
Najlepszym jednak przykładem antyprometejskiej postawy \emph{nowego trickstera kalevalicznego} jest pałająca żądzą zysku i~zemsty Louhi. Pierwsza pobudka widoczna jest, gdy świeżo co wykute przez Ilmarinena Sampo zostaje zamknięte w~kamiennej górze Pohjoli (X:~423--432):
\end{spacing}
\newpage

%\vspace{1cm}

{\centering
\small
\begin{longtable}{c c}
Cieszy się Pani Północy, & Niin ihastui Pohjan akka;\\
W mig pokaźne Sampo schowa & saattoi sitte sammon suuren\\
W kamiennej górze Pohjoli, & Pohjolan kivimäkehen,\\
Wewnątrz kopca miedzianego, & vaaran vaskisen sisähän\\
Za dziewięcioma skoblami; & yheksän lukon ta'aksi.\\
Zakorzeniło się Sampo & Siihen juuret juurrutteli\\
Na dziewięć sążni głęboko, & yheksän sylen syvähän:\\
Jeden korzeń w matce-ziemi, & juuren juurti maaemähän,\\
Drugi korzeń w wirze wodnym, & toisen vesiviertehesen,\\
Trzeci na wzgórzu pod chatą. & kolmannen kotimäkehen.\\
\end{longtable}
}
%\vspace{0.5\baselineskip}
\begin{spacing}{1.3}
%\noindent
Pragnienie zemsty natomiast wyraża Louhi modlitwą do Ukko (XLV:~15--22), na skutek której (z~bożą pomocą) bogini Loviatar rodzi swoje potworne potomstwo, mające zniszczyć krainę Kalevy. Gdy jednak Väinämöinen, będący znachorem, uzdrawia (również z~bożą pomocą) chorych i~niedomagających, Louhi ucieka się do ostatniej próby sprowadzenia ciemności i~smutku na świat i~zabiera z~nieboskłonu słońce i~księżyc, co --- jak wykazano już wcześniej --- daje się we znaki nawet najwyższemu bóstwu.
\end{spacing}
%\newpage
\section{Antytrickster}
\begin{spacing}{1.3}
Trickstera cechuje pozorna głupota i~chaotyczne działanie mające doprowadzić go do celu, który okaże się w~ogólnym podsumowaniu mądry (i~często też szlachetny). \emph{Nowy trickster kalevaliczny} natomiast cechuje się często działaniem przeciwnym --- w~jego przypadku skutek działań jest często opłakany, a~powiedzenie, iż dobrymi uczynkami wybrukowano piekło, daje znać o~swojej prawdziwości. Boski kowal Ilmarinen daje się głupio oszukać szerszeniowi, co skutkuje tym, iż nieposkromione jeszcze żelazo obraca się przeciwko człowiekowi (IX:~259--266):
\end{spacing}
%\newpage
{\centering
\small
\begin{longtable}{c c}
{[...]} & {[...]}\\
Stal zabulgotała gniewnie, & ,,Sai siitä teräs pahaksi,\\
Szaleć zaczęło żelazo, & rauta raivoksi rupesi,\\
Nie dochowało przysięgi, & petti, vaivainen, valansa,\\
Swoją cześć jak pies pożarło, & söi kuin koira kunniansa:\\
Brata rąbało, nieszczęsne, & veisti, raukka, veljeänsä,\\
Własny ród kłami kąsało. & sukuansa suin piteli,\\
Jęło przelewać posokę, & veren päästi vuotamahan,\\
Powodować krwi rozlanie.'' & hurmehen hurahtamahan.''\\
\end{longtable}
}
%\vspace{0.5\baselineskip}
\newpage
\begin{spacing}{1.3}
Kullervo z~kolei jest przykładem antytrickstera, którego złe uczynki prowadzą jedynie do coraz większego zła. Gdy ma się zaopiekować dzieckiem --- traci ono życie (po drodze jeszcze oprawca wyrywa mu rękę i~wyłupia oko --- XXXI:~213--228), a~gdy pała chęcią zemsty w~myśl zasady ,,oko za oko'' --- bez zastanowienia oddaje się ferworowi zniszczenia (XXXIII:~143--160):
\end{spacing}
{\centering
\small
\begin{longtable}{c c}
Złośliwy pastuch niecnota, & Tuo pahainen paimen raiska,\\
Kullervo, Kalervy synek & Kullervo, Kalervon poika,\\
Przywiedzie pod dom niedźwiedzie, & ajoi kontiot kotihin,\\
Wilczą zgraję do zagrody. & susikarjan kartanolle.\\
Jeszcze poduszczał kudłacze, & Vielä neuvoi karhujansa,\\
Wilkom tłumaczył wyraźnie: & susillensa suin puheli:\\
,,Rwijcie biodra gospodyni, & ,,Repäise emännän reisi,\\
Poprzegryzajcie jej łydki, & pure puoli pohkeata,\\
Kiedy przyjdzie was zobaczyć, & kun tulevi katsomahan,\\
Do dojenia się nachyli!'' & lyykistäikse lypsämähän!''\\
Wziął piszczałę z krowiej kości, & Teki luikun lehmän luista,\\
Wybrał wołu róg donośny, & härän sarvesta helinän,\\
Trąbkę z nogi Jarzębatej, & torven Tuomikin jalasta,\\
Na fujarkę staw Łaciatej; & pillin Kirjon kinterestä.\\
Popiskiwał swą piszczałą, & Lujahutti luikullansa,\\
Zatrąbił na swoim rogu & toitahutti torvellansa\\
Trzykroć z pagórka w pobliżu, & kolmasti kotimäellä,\\
A w opłotkach sześciokrotnie. & kuuesti kujosten suussa.\\
\end{longtable}
}
%\vspace{0.5\baselineskip}
\begin{spacing}{1.3}
Co się tyczy Lemminkäinena --- jego głupota ściąga na wracających ze skradzionym Sampo kompanów masę kłopotów. Głupie zachowanie, polegające na śpiewaniu w~łodzi (przed czym przestrzega go Väinämöinen) skutkuje tym, że przerażony ptak leci do Pohjoli i~budzi Louhi, która szybko orientuje się, co zaszło. Ostatecznie bohaterowie wychodzą cało z~opresji, ale Sampo jest stracone.
%\newpage
Działania Ukko również przynoszą ciężkie szkody, czego dobrym przykładem jest utrata przez Väinämöinena jego ukochanego instrumentu (XLII:~496--502) na skutek burzy zesłanej przez bóstwo. Jeszcze lepszą emanacją Ukko jako postaci w~duchu antytricksteryzmu jest rozniecenie przez niego ognia w~niebie, który następnie spada na ziemię i~dokonuje potwornych zniszczeń --- doprowadza m.in. do spalenia mitycznego jeziora Alue, jak również rani dziewczęta i~dzieci (XLVII:~191--202):
\end{spacing}
{\centering
\small
\begin{longtable}{c c}
{[...]} & {[...]}\\
Pannom piersi pokaleczył, & rikkoi rinnat tyttäriltä,\\
Poparzył sutki dziewczętom, & neitosilta nännit näppi,\\
Poranił chłopcu kolana, & turmeli pojalta polvet,\\
Spalił brodę gospodarza. & isännältä parran poltti.\\
{[...]} & {[...]}\\
Spalił dzieciątko w kołysce, & poltti lapsen kätkyestä,\\
Spalił piersi rodzicielce; & poltti paarmahat emolta.\\
Poszło dzieciątko do Many, & Se lapsi meni Manalle,\\
Chłopaczek do Tuoneli, & toki poika Tuonelahan,\\
{[...]} & {[...]}\\
\end{longtable}
}
%\vspace{0.5\baselineskip}
\newpage
\section{Pycha}
\begin{spacing}{1.3}
\emph{Nowi tricksterzy kalevaliczni} są egoistyczni i~pyszni, a~ich ego jest nazbyt wybujałe. Ilmarinen wiele razy podkreśla swój kowalski kunszt, okraszając swoje wypowiedzi iście szelmowską pewnością siebie (X:~213--216, 271--280; XLIX:~341--344):
\end{spacing}
{\centering
\small
\begin{longtable}{c c}
,,Zdarzyło mi się oglądać & ,,Lienen tullut tuntemahan\\
Kowala Ilmarinena, & tuon on seppo Ilmarisen,\\
Skoro właśnie ja nim jestem, & kun olen itse Ilmarinen,\\
Owym biegłym rzemieślnikiem.'' & itse taitava takoja.''\\
{[...]} & {[...]}\\
,,Jestem w stanie wykuć Sampo, & ,,Saattanen takoa sammon,\\
Wyklepać wieczko ozdobne, & kirjokannen kalkutella\\
Z pióra w ogonie łabędzia, & joutsenen kynän nenästä,\\
Z kropli mleka dojnej krowy, & maholehmän maitosesta,\\
Z małego ziarnka jęczmienia, & ohran pienestä jyvästä,\\
Z pasma wełny młodej jarki, & kesäuuhen untuvasta,\\
Skoro wykułem niebiosa, & kun olen taivoa takonut,\\
Nieba spoiłem sklepienie, & ilman kantta kalkuttanut\\
Jeszcze przed początkiem świata, & ilman alkusen alutta,\\
Bez śladu szwu najlżejszego.'' & riporihman tehtyisettä.''\\
{[...]} & {[...]}\\
,,A cóż niby w tym dziwnego, & ,,Ei tuo kumma ollekana,\\
Że porządnie kuć potrafię, & jos olen takoja tarkka,\\
Skoro wykułem niebiosa, & kun olen taivoa takonut,\\
Wieczko powietrza sklepałem?'' & ilman kantta kalkutellut.''\\
\end{longtable}
}
%\vspace{0.5\baselineskip}
%\newpage
\begin{spacing}{1.3}
Ilmatar, bogini powietrza, nie omieszkuje podczas spotkania przypomnieć bohaterom o~swojej wspaniałej urodzie równej sześciu pannom (XLVII:~154--158), a~Joukahainen za bardzo wierzy w~siłę swoich zaklęć przed pojedynkiem z~Väinämöinenem (III:~55--60):
\end{spacing}
{\centering
\small
\begin{longtable}{c c}
Skoro zechcę się z nim zmierzyć, & Jos tahon tasalle panna,\\
Iść o lepsze jak mąż z mężem, & miesten verroille vetäitä,\\
Sam zaklęć swoich użyję, & itse laulan laulajani,\\
Zamawiania swe wypowiem, & sanelen sanelijani:\\
Najlepszego piewcę pieśni & laulan laulajan parahan\\
W najlichszego przeinaczę, & pahimmaksi laulajaksi,\\
{[...]} & {[...]}\\
\end{longtable}
}
%\vspace{0.5\baselineskip}
\begin{spacing}{1.3}
Sam Väinämöinen pyszni się przed kolejnymi dziewczętami podkreślając swoje walory (VIII:~88--90), gdyż z~niego ,,mąż nie pierwszy lepszy''. Z~kolei Lemminkäinen podkreśla swoją męskość i~odwagę, co wręcz ocieka groteską (wszakże przez swoje ekscesy szybko musi on z~Północy uciekać). Szelmowsko ocenia potencjalnych ,,mężów'' poprzez pryzmat tego, czy potrafiliby oni naciągnąć cięciwę w~jego łuku --- zupełnie tak, jak postępuje Odyseusz po powrocie z~wojny do Itaki (XXVI:~357--362):
\end{spacing}
\newpage
{\centering
\small
\begin{longtable}{c c}
,,Tego uznałbym za męża & ,,Sen mä mieheksi sanoisin,\\
I miałbym go za mocarza, & urohoksi arveleisin,\\
Kto by mój łuk napiąć zdołał, & joka jouseni vetäisi,\\
Naciągnął na nim cięciwę & kiveräni kiinnittäisi\\
W chatach pochmurnej Pohjoli, & noilla Pohjolan tuvilla,\\
Smętnych siedliskach Sarioli.'' & Sariolan salvoksilla.''\\
\end{longtable}
}
\vspace{1\baselineskip}

\section{Ostateczna porażka}
\begin{spacing}{1.3}
Prototypowy trickster, niekiedy odwracając kłopotliwą sytuację lub uciekając się do szeregu forteli, osiąga swój zamierzony cel. W~przypadku \emph{nowego trickstera kalevalicznego} mamy często do czynienia z~porażką bohatera, co stanowi swoiste przeciwieństwo jego zwycięstw pomniejszych. W~tej materii przoduje Väinämöinen. Wymykający się wielokrotnie śmierci i~obdarowujący Kalevalę wszelkimi dobrami mag i~czarownik w~ostatecznym rozrachunku przegrywa z~nowonarodzonym synem Marjatty i~jest zmuszony do odpłynięcia (podkreśla jednak, że pewnego dnia --- gdy ludzie znów będą go potrzebować --- powróci, by odbudować Sampo) (L:~501--512):
\end{spacing}
{\centering
\small
\begin{longtable}{c c}
Po czym stary Väinämöinen & Siitä vanha Väinämöinen\\
Raźnie odbije od brzegu & laskea karehtelevi\\
Nad burtą miedzianej łodzi, & venehellä vaskisella,\\
W spiżem pokrytej komiedze, & kuutilla kuparisella\\
W stronę górnych ziem-macierzy, & yläisihin maaemihin,\\
Gdzie się z dołem nieba schodzą. & alaisihin taivosihin.\\
Zniknął tam ze swym czółenkiem, & Sinne puuttui pursinensa,\\
Ze swą łódeczką pogrążył, & venehinensä väsähtyi.\\
Zostawił po sobie gęśle, & Jätti kantelon jälille,\\
Muzyckie naczynie Suomi, & soiton Suomelle sorean,\\
Na radość wieczną ludowi, & kansalle ilon ikuisen,\\
Pieśń potężną swojej dziatwie. & laulut suuret lapsillensa.\\
\end{longtable}
}
%\vspace{0.5\baselineskip}
\begin{spacing}{1.3}
Joukahainen, postać przeciwna Väinämöinenowi, również ponosi klęskę. Przegrywa on z~kretesem pojedynek na wiedzę, a~żeby w~ogóle ujść z~niego z~życiem, oddaje przeciwnikowi swoją siostrę jako okup. Wiedziony następnie zemstą dokonuje zamachu na Väinämöinenie, wskutek którego emanacja szczęścia, radości i~dobroci zostaje przebita na wylot strzałą na wysokości łopatki. W~dłuższej perspektywie zamach nie szkodzi zaklinaczowi, który po przejściu rekonwalescencji dobitnie podkreśla porażkę Joukahainena (X:~21--30):
\end{spacing}
\newpage
{\centering
\small
\begin{longtable}{c c}
Tam wyrzecze takie słowa, & Siinä tuon sanoiksi virkki,\\
Takie wypowiedział zdanie: & itse lausui ja pakisi:\\
,,A zduś, wilku, jasnowidza, & ,,Syö, susi, unennäkijä,\\
Bierz, chorobo, Lapończyka! & tapa, tauti, lappalainen!\\
Zapowiedział, że nie wrócę, & Sanoi ei saavani kotihin\\
By oglądać w swoim domu & enämpi elävin silmin\\
Na własne oczy słoneczko, & sinä ilmoisna ikänä,\\
I miesiączek jaśniejący & kuuna kullan valkeana\\
Na wrzosowiskach Väinöli, & näille Väinölän ahoille,\\
Na polanach Kalevali!'' & Kalevalan kankahille.''\\
\end{longtable}
}
%\vspace{0.5\baselineskip}
\begin{spacing}{1.3}
Co się zaś tyczy pozostałych bohaterów --- Kullervo popełnia samobójstwo na wieść o~akcie kazirodztwa, podkreślając wcześniej wielokrotnie swoją samotność i~tragedię. Ilmarinen zostaje wykorzystany do wykucia cudownego młynka (za wykonanie którego miał ożenić się z~córką Louhi) --- kowal panny nie dostaje, a Louhi zabiera Sampo. Pod koniec eposu wykuwa sobie Ilmarinen nawet nową żonę ze złota, co jeszcze wyraźniej potęguje jego porażkę. W~kwestii Sampo --- roztrzaskuje się ono na kawałki, po czym wpada do morza. Wtedy wściekła Louhi próbuje mścić się na całej Kalevali, co jest z~kolei jej ostateczną porażką --- wszakże choroby pani Loviatar zostają pokonane, a~czarownica zmuszona jest zwrócić nieboskłonowi przywłaszczone ciała niebieskie.

Cechy \emph{nowego trickstera kalevalicznego} pozwalają na modyfikację stereotypowego trickstera w~fińskim eposie narodowym --- jest nim ,,złączony specyficzną więzią z~bóstwami i~niegrzeszący mądrością zaklinacz o~zbyt wybujałym ego, który często działa na świadomą szkodę dla innych i~w~ostatecznym rozrachunku ponosi klęskę''.
\end{spacing}

\clearpage{\pagestyle{empty}\cleardoublepage}
\part{ANEKS}
\clearpage{\pagestyle{empty}\cleardoublepage}

\begin{figure}
\centering
\includegraphics[height=20cm]{Eshu}
\caption[\`{E}\d{s}\`{u}-Elegbara]{{\`{E}\d{s}\`{u}-Elegbara} \\ figurka klęcząca, drewno, Yor\`{u}b\'{a}, XIX/XX wiek \\ {\footnotesize {New York, Brooklyn Museum, Arts of Africa, nr 81.102}}}
\end{figure}

\newpage
\begin{figure}[h]
\centering
\includegraphics[width=\textwidth]{Hermes}
\caption[Hermes]{{Hermes} \\ {\footnotesize {Alexander Stuart Murray}, 1895, {\it Manual of Mythology}, Philadelphia: David McKay.}}
\end{figure}

\newpage
\begin{figure}[h]
\centering
\includegraphics[width=\textwidth]{Loki}
\caption[Loki]{{Loge (Loki)} \\ {Arthur Rackham:} {\it The Rhine's Pure-Gleaming Children Told Me of Their Sorrow} \\ {\footnotesize {Richard Wagner}, 1910, {\it The Rhinegold and the Valkyrie}, London: William Heinemann, s.~28.}}
\end{figure}

\newpage
\begin{figure}[h]
\centering
\includegraphics[width=\textwidth]{Loki3}
\caption[Jörmungandr, Fenrir \& Hel]{{Jörmungandr, Fenrir \& Hel} \\ {Willy Pogany:} {\it The Children of Loki} \\ {\footnotesize Padraic Colum, 1920, {\it The Children of Odin}, New York: The Macmillan Company.}}
\end{figure}

\newpage
\begin{figure}[h]
\centering
\includegraphics[width=\textwidth]{Jahwe}
\caption[Jahwe]{{Jahwe} \\ {Albrecht Dürer:} {\it Apokalipsa: św. Jan zjadający księgę} \\ {\footnotesize Karlsruhe, Staatliche Kunsthalle}}
\end{figure}

\newpage
\begin{figure}[h]
\centering
\includegraphics[height=15cm]{VAN}
\caption[Väinämöinen]{{Väinämöinen} \\ {Robert Wilhelm Ekman:} {\it Väinämöisen soitto} \\ {\footnotesize Matti Ruotsalainen, Kuvataiteen keskusarkisto}}
\end{figure}

\newpage
\begin{figure}[h]
\centering
\includegraphics[width=\textwidth]{JK}
\caption[Joukahainen]{{Joukahainen} \\ {Akseli Gallén-Kallela:} {\it Joukahaisen kosto} \\ {\footnotesize Kari Lehtinen, Turun taidemuseo}}
\end{figure}

\newpage
\begin{figure}[h]
\centering
\includegraphics[width=\textwidth]{ILM}
\caption[Ilmarinen]{{Ilmarinen} \\ {Robert Wilhelm Ekman:} {\it Väinämöinen laulaa loitsuja} \\ {\footnotesize Antti Karvinen, Kuvataiteen keskusarkisto}}
\end{figure}

\newpage
\begin{figure}[h]
\centering
\includegraphics[height=20cm]{Louhi}
\caption[Louhi]{{Louhi} \\ {Tapio Tapiovaara:} {\it Louhi, Pohjan akka} \\ {\footnotesize Hannu Pakarinen, Kuvataiteen keskusarkisto}}
\end{figure}

\newpage
\begin{figure}[h]
\centering
\includegraphics[height=15cm]{LEM}
\caption[Lemminkäinen]{{Lemminkäinen} \\ {Carl Eneas Sjöstrand:} {\it Lemminkäinen} \\ {\footnotesize Henri Tuomi, Kuvataiteen keskusarkisto}}
\end{figure}

\newpage
\begin{figure}[h]
\centering
\includegraphics[width=\textwidth]{Kul}
\caption[Kullervo]{{Kullervo} \\ {Akseli Gallén-Kallela:} {\it Katuva Kullervo} \\ {\footnotesize Helsinki, Ateneumin taidemuseo}}
\end{figure}

\newpage
\begin{figure}[h]
\centering
\includegraphics[width=1.005\textwidth]{ILMUKKO}
\caption[Ilmatar \& Ukko]{{Ilmatar \& Ukko} \\ {Robert Wilhelm Ekman:} {\it Ilmatar huutaa Ukkoa} \\ {{\footnotesize Antti Karvinen, Kuvataiteen keskusarkisto}}}
\end{figure}
%\clearpage

\pagebreak
\newpage
%\newpage
\begin{figure}[h]
\centering
\includegraphics[width=1.005\textwidth]{Ilmatar}
\caption[Ilmatar]{{Ilmatar} \\ {Robert Wilhelm Ekman:} {\it Ilmatar} \\ {\small \texttt{\url{www.kantele.net}}}}
\end{figure}

\newpage
\begin{figure}[h]
\centering
\includegraphics[height=20cm]{Marjatta}
\caption[Marjatta]{Marjatta \\ {Akseli Gallén-Kallela:} {\it Marjatta rannalla} \\ {\small \texttt{\url{www.wahooart.com}}}}
\end{figure}

\ChapterOutsidePart
\addtocontents{toc}{\protect\vspace*{0.8\baselineskip}}
\chapter*{Tiivistelmä}
\addcontentsline{toc}{chapter}{Tiivistelmä --- streszczenie pracy w~języku fińskim}
\selectlanguage{finnish}
\begin{spacing}{1.3}
\noindent
Pro gradu -työssäni käsittelen trickster-hahmon prototyyppiä ja stereotyyppiä suomalaisessa kansalliseepoksessa, Kalevalassa. Valitsin tämän aiheen, koska minua kiinnostaa suomalainen kirjallisuus, pohjoisten kansojen uskomukset ja jumalan huiputtaja -hahmo, jonka ominaisuudet löydetään kaikista maailman mytologioista.

Tricksterilla (jota kutsutaan myös \emph{jumalan hulluksi}, \emph{huijariksi}, \emph{petkuttajaksi}, \emph{temppuilijaksi} ja \emph{veijariksi}) on kaksi päätoimintaa. Hän on mytologinen olento, joka toimii kevytmielisen mielikuvituksen symbolina, ja kulttuurisankari, joka voi elää sekä huonontuneena jumalana, että pyhänä syntisenä. Trickster on myös jumalan hullu, joka luo kulttuuria ja moraalia häiritsemällä jumalaa maailman luomisen aikana. Petkuttaja-hahmo tutkittiin ensimmäisen kerran 1890-luvulla. Paul Radinin kirjoittama \emph{The Trick- ster: A~Study in American Indian Mythology} -kirjan takia trickster-hahmo tuli suosituimmaksi maailmalla. Puolassa tricksterista on kirjoittanut mm. Jerzy Wasilewski. Jumalan hullun ominaisuuksia voi löytää useita Kalevalasta, mutta aiheesta ei ole kirjoitettu paljon Suomessa.

Kognitiivismista kertova teoreettinen osa perustuu sekä puolalaisten (mm. Jerzy Bartmińskin, Renata Grzegorczykowan, Henryk Kardelan ja Anna Wierzbickan), että ulkomaalaisten tutkijoiden (kuten George Lakoffin ja Ludwig Wittgensteinin) teoksiin. Kirjoittaessani jumalan hullusta olen käyttänyt lähteinä mm. Paul Radinin ja Carl Gustav Jungin tekstejä. Myös Jacek Sieradzanin \emph{Szaleństwo w religiach świata} -kirja, joka kertoo jumalan tyhmyydestä erilaisissa uskonnollisissa järjestelmissä, on vaikuttanut työhöni merkittävästi. Tutkimusaineistona toimii Jerzy Litwiniukin puolantama \textit{Kalevala}. Tutkimukseni koostuu neljästä pääosasta ja liitteestä, jossa esittelen trickstereiden kuvia.
\newpage
\vspace*{2.5\baselineskip}
Työni ensimmäisessä osassa esittelen kognitiivisen koulukunnan pääajatuksia ja niiden perusteita sekä kognitivismin jaon symboliseen ja subsymboliseen. Seuraavissa luvuissa kerron Ronald Langackerin luomasta kognitiivisesta kielopista ja muista teorioista, jotka ovat siihen vaikuttaneet --- Eleanor Roshin prototyypin teoriasta, George Lakoffin kognitiivisista malleista ja Ludwig Wittgensteinin perheyhtäläisyys-ajatuksesta. Tarkastelen myös stereotyyppejä ja järjestelmiä, niiden eroja ja Zadehin sumeita joukkoja (engl. \emph{fuzzy sets}).

Trickster-hahmo on keskeinen aihe tutkimukseni toisessa osassa. Viidennessä, kuudennessä ja seitsemännessä luvussa tutkin huiputtajan käsitettä antropologisesta ja šamanistisesta näkökulmasta. Käsittelen sekä hänen irvokasta seksuaalisuuttaan (Wak\-djunkaga), että jumalalta saatua lähetystyötään ja kaksiarvoisuuttaan (\`{E}\d{s}\`{u}-Elegbara, Hermes). Näytän myös tricksterin tyhmyyttä ja tietämättömyyttä (Loki). Kappaleessa kahdeksan kuvailen Jahvea rikollisena jumalana, jolta Alf Hiltebeitel mukaan löytyy tricksterin ominaisuuksia.

Kolmas osa toimii työssäni kognitiivisen kieliopin ja trickster-hahmon synteesina. Kuvailen yhdeksän kalevalaista hahmoa --- Väinämoistä, Joukahaista, Ilmarista, Louhea, Lemminkäistä, Kullervoa, Ukkoa, Ilmatarta ja Marjattaa --- soveltaen kahdeksaa tricksterin aluetta, jotka olen erottanut tarkkaan analyysin avulla. Tutkimukseni todistaa, että eepoksen hahmoissa on tricksterin heijastuksia. Sen mukaan, kuinka monta tricksterin aluetta ja mikä merkitys kyseisellä hahmolla on, jaan kalevalaiset tricksterit \emph{suuriin} ja \emph{pieniin} trickstereihin. Analyysin mukaan, suomalainen stereotyyppinen jumalan hullu on tyhmä tietäjä, jolla näkyy myös yhteyttä jumalaan tai jumaloihin.

Neljännessä osassa luon \emph{uusi kalevalainen tricksteri} -käsitteen laajentamalla tricksterin ominaisuusten joukon ja lisäämällä kalevalaisia attribuutteja. Ne attribuutit eivät ehdottomasti ole tricksterin ominaisuuksia, mutta ne ilmenevät säännöllisesti kalevalaisten teoissa. Näitä ovat: Prometheuksen ja tricksterin päinvastaisuus (hahmojen tyhmät toimet johtavat pahaan), ylpeys ja lopullinen epäonnistuminen (esim. Väinämöinen lähtee Kalevalasta naurunalaistettuna).

\end{spacing}

\ChapterOutsidePart
\chapter*{Literatura}
\addcontentsline{toc}{chapter}{Literatura}

\begin{thebibliography}{}

\bibitem[\protect\citename{Anusiewicz \& Bartmiński, }1998]{anubar98}
%{\sc Anusiewicz Janusz \& Bartmiński Jerzy (red.)},
{ANUSIEWICZ Janusz \& BARTMIŃSKI Jerzy (red.)},
\newblock 1998,
\newblock {\em Język a~Kultura, Tom~12: Stereotyp jako przedmiot lingwistyki: Teoria, metodologia, analizy empiryczne},
\newblock Wrocław: Towarzystwo Przyjaciół Polonistyki Wrocławskiej.
%\newblock s.~109--114.

\bibitem[\protect\citename{Barlow \& Kemmer, }1999]{barkem99}
%{\sc Barlow Michael \& Kemmer Suzanne (red.)},
{BARLOW Michael \& KEMMER Suzanne (red.)},
\newblock 1999,
\newblock {\em Usage Based Models of Language},
\newblock Stanford: CSLI Publications.
%\newblock s.~12.

\bibitem[\protect\citename{Bartmiński, }1988]{bar88}
%{\sc Bartmiński Jerzy (red.)},
{BARTMIŃSKI Jerzy (red.)},
\newblock 1988,
\newblock {\em Konotacja},
\newblock Lublin: UMCS.
%\newblock s.~9--34.

\bibitem[\protect\citename{Bartmiński, }1992]{bar92}
%{\sc Bartmiński Jerzy (red.)},
{BARTMIŃSKI Jerzy (red.)},
\newblock 1992,
\newblock {\em Etnolingwistyka, Tom~4},
\newblock Lublin: UMCS.
%\newblock s.~33.

\bibitem[\protect\citename{Bartmiński, }1993]{bar93}
%{\sc Bartmiński Jerzy (red.)},
{BARTMIŃSKI Jerzy (red.)},
\newblock 1993,
\newblock {\em Encyklopedia Kultury Polskiej XX wieku, Tom~2: Współ- czesny język polski},
\newblock Wrocław: Wiedza o~Kulturze.
%\newblock s. 363-387.

\bibitem[\protect\citename{Bartmiński, }1985]{bar85}
%{\sc Bartmiński Jerzy},
{BARTMIŃSKI Jerzy},
\newblock 1985,
\newblock ,,Stereotyp jako przedmiot lingwistyki'',
\newblock {[w:]} {BASAJ Mieczy- sław \& RYTEL-SCHWARZ Danuta (red.)}, \newblock {\em Z~Problemów Frazeologii Polskiej i~Słowiańskiej, Tom~3},
\newblock Wrocław: PAN,
\newblock s.~25--53.

\bibitem[\protect\citename{Bartmiński \& Panasiuk, }1993]{barpan93}
{BARTMIŃSKI Jerzy \& PANASIUK Jolanta},
\newblock 1993,
\newblock ,,Stereotypy językowe'',
\newblock {[w:]} {BARTMIŃ- SKI Jerzy (red.)},
\newblock {\em Encyklopedia Kultury Polskiej XX wieku, Tom~2: Współczesny język polski},
\newblock Wrocław: Wiedza o~Kulturze,
\newblock s.~363--387.

\bibitem[\protect\citename{Basaj \& Rytel-Schwarz, }1982]{basryt82}
%{\sc Basaj Mieczysław \& Rytel-Schwarz Danuta (red.)},
{BASAJ Mieczysław \& RYTEL-SCHWARZ Danuta (red.)},
\newblock 1982,
\newblock {\em Z~Problemów Frazeologii Polskiej i~Słowiańskiej, Tom~1},
\newblock Wrocław: PAN.
%\newblock s. 33--47.

\bibitem[\protect\citename{Basaj \& Rytel-Schwarz, }1985]{basryt85}
%{\sc Basaj Mieczysław \& Rytel-Schwarz Danuta (red.)},
{BASAJ Mieczysław \& RYTEL-SCHWARZ Danuta (red.)},
\newblock 1985,
\newblock {\em Z~Problemów Frazeologii Polskiej i~Słowiańskiej, Tom~3},
\newblock Wrocław: PAN.
%\newblock s. 25--53.

\bibitem[\protect\citename{Blaeser, }1996]{bla96}
%{\sc Blaeser Kimberley},
{BLAESER Kimberley},
\newblock 1996,
\newblock {\em Gerald Vizenor: Writing in the Oral Tradition},
\newblock Norman: The University of Oklahoma Press.

\bibitem[\protect\citename{Bolton, }1962]{bol62}
%{\sc Bolton James David},
{BOLTON James David},
\newblock 1962,
\newblock {\em Aristeas of Proconnesus},
\newblock Oxford: Clarendon Press.

\bibitem[\protect\citename{Brinton, }1868]{bri68}
%{\sc Brinton Daniel Garrison},
{BRINTON Daniel Garrison},
\newblock 1868,
\newblock {\em The Myths of the New World: A~Treatise on the Sym- \linebreak bolism and Mythology of the Red Race of America},
\newblock New York: Leypoldt \& Holt.

\bibitem[\protect\citename{Brown, }1990]{bro90}
%{\sc Brown Cecil},
{BROWN Cecil},
\newblock 1990,
\newblock ,,A survey of category types in natural language'',
%\newblock {[w:]} {\sc Tsohatzidis Savas (red.)},
\newblock {[w:]} {TSOHA- TZIDIS Savas (red.)},
\newblock {\em Meanings and Prototypes: Studies in Linguistic Categorization},
\newblock London: Routledge,
\newblock s.~17--47.
\newpage
\bibitem[\protect\citename{Chomsky, }1982]{cho82}
{CHOMSKY Noam},
\newblock 1982,
\newblock {\em Some Concepts and Consequences of the Theory of Government and Binding},
\newblock Cambridge: The MIT Press.

\bibitem[\protect\citename{Chomsky, }1986a]{cho86a}
{CHOMSKY Noam},
\newblock 1986a,
\newblock {\em Barriers},
\newblock Cambridge: The MIT Press.

\bibitem[\protect\citename{Chomsky, }1986b]{cho86b}
{CHOMSKY Noam},
\newblock 1986b,
\newblock {\em Knowledge of Language: Its Nature, Origin, and Use},
\newblock New York: Praeger.

\bibitem[\protect\citename{Coleman \& Kay, }1981]{colkay81}
{COLEMAN Linda \& KAY Paul},
\newblock 1981,
\newblock {\em Prototype Semantics: The English Word ,,Lie''}, \linebreak
\newblock {[w:]} ,,Language 57'',
\newblock nr~1, s.~26--44.

\bibitem[\protect\citename{Colum, }1920]{col20}
{COLUM Padraic},
\newblock 1920,
\newblock {\em The Children of Odin},
\newblock New York: The Macmillan Company.

\bibitem[\protect\citename{Davis, }2000]{dav00}
{DAVIS Erik},
\newblock 2000,
\newblock {\em Trikster na rozdrożach},
\newblock tłum. {MISIUNA Dariusz},
\newblock {[w:]} ,,ALBO Albo: Trikster'',
\newblock nr~4, s.~45--56.

\bibitem[\protect\citename{Dąbrowska \& Kubiński, }2003]{dabkub03}
{DĄBROWSKA Ewa \& KUBIŃSKI Wojciech (red.)},
\newblock 2003,
\newblock {\em Akwizycja języka w~świetle języ- koznawstwa kognitywnego},
\newblock Kraków: Universitas.
%\newblock s.~30.

\bibitem[\protect\citename{Diamond, }2010]{dia10}
{DIAMOND Stanley},
\newblock 2010,
\newblock ,,Wstęp: Hiob i~trickster'',
\newblock {[do:]} {RADIN Paul},
\newblock {\em Trickster: Studium mitologii Indian północnoamerykańskich},
\newblock tłum. {TOPCZEWSKA Anna},
\newblock Warszawa: WUW,
\newblock s.~9--20.

\bibitem[\protect\citename{Doty \& Hynes, }1993]{dothyn93}
{DOTY William \& HYNES William},
\newblock 1993,
\newblock {\em Mythical Trickster Figures: Contours, Contexts, and Criticisms},
\newblock Tuscaloosa: The University of Alabama Press.

\bibitem[\protect\citename{Duch, }1998]{duc98}
{DUCH Włodzisław},
\newblock 1998,
\newblock {\em Czym jest kognitywistyka?},
\newblock {[w:]} ,,Kognitywistyka i~Media w~Edukacji'',
\newblock nr~1, s.~9--50.

\bibitem[\protect\citename{Edda, }1986]{edd86}
{\emph{Edda Poetycka}},
\newblock 1986,
\newblock tłum. {ZAŁUSKA-STRÖMBERG Apolonia},
\newblock Warszawa: Ossolineum.

\bibitem[\protect\citename{Fairclough, }2005]{fai05}
{FAIRCLOUGH Matthew},
\newblock 2005,
\newblock {\em A~Creative Healing Aspect within the Arts?},
\newblock {[w:]} ,,Trickster’s Way 4'',
\newblock nr~1, s.~1--25.

\bibitem[\protect\citename{Golka, }2009]{gol09}
{GOLKA Maria},
\newblock 2009,
\newblock ,,Teoria zbiorów rozmytych jako formalizacja teorii prototypu: Problemy z~kompozycjonalnością'',
\newblock {[w:]} {ŁUPKOWSKI Paweł, MOZYRKO Bartosz \& POROWSKA Agnieszka (red.)},
\newblock {\em Poznańskie Forum Kognitywistyczne: Teksty Pokonferencyjne, Nr~3},
\newblock Poznań: Zakład Logiki i~Kognitywistyki, Instytut Psychologii UAM,
\newblock s.~31--39.

\bibitem[\protect\citename{Griffiths, }1982]{gri82}
{GRIFFITHS Bede},
\newblock 1982,
\newblock {\em Zaślubiny Wschodu z~Zachodem},
\newblock tłum. {JASIŃSKI Grzegorz},
\newblock Poznań: Zysk i~S-ka.

\bibitem[\protect\citename{Grzegorczykowa, }1998]{grz98}
{GRZEGORCZYKOWA Renata},
\newblock 1998,
\newblock ,,O~rozumieniu prototypu i~stereotypu we współczesnych teoriach semantycznych'',
\newblock {[w:]} {ANUSIEWICZ Janusz \& BARTMIŃSKI Jerzy (red.)},
\newblock {\em Język a~Kultura, Tom~12: Stereotyp jako przedmiot lingwistyki: Teoria, metodologia, analizy empiryczne},
\newblock Wrocław: Towarzystwo Przyjaciół Polonistyki Wrocławskiej,
\newblock s.~109--114.

\bibitem[\protect\citename{Haavio, }1979]{haa79}
{HAAVIO Martti},
\newblock 1979,
\newblock {\em Mitologia fińska},
\newblock tłum. {LITWINIUK Jerzy},
\newblock Warszawa: PIW.

\bibitem[\protect\citename{Hiltebeitel, }1989]{hil89a}
{HILTEBEITEL Alf}, %\(red.)
\newblock 1989,
\newblock {\em Criminal Gods and Demon Devotees: Essays on the Guardians \linebreak of Popular Hinduism},
\newblock Albany: State University of New York Press.

%\bibitem[\protect\citename{Hiltebeitel, }1989]{hil89b}
%{\sc Hiltebeitel Alf},
%\newblock 1989,
%\newblock ,,Introduction'',
%\newblock {[do:]} {\sc Tenże (red.)},
%\newblock {\em Cryminal gods and demon devotees: Essays on the guardians of popular Hinduism},
%\newblock Albany: State University of New York Press,
%\newblock s.~1--18.

\bibitem[\protect\citename{Hyde, }1998/2008]{hyd9808}
{HYDE Lewis},
\newblock 1998,
\newblock {\em Trickster Makes this World: How Disruptive Imagination Creates Culture},
\newblock Edinburgh: Canongate.

\newpage
\bibitem[\protect\citename{Jackendoff, }1983]{jac83}
{JACKENDOFF Ray},
\newblock 1983,
\newblock {\em Semantics and Cognition},
\newblock Cambridge: The MIT Press.

\bibitem[\protect\citename{Jackendoff, }1987]{jac87}
{JACKENDOFF Ray},
\newblock 1987,
\newblock {\em Consciousness and the Computational Mind},
\newblock Cambridge: The MIT Press.

\bibitem[\protect\citename{Johnson, }1987]{joh87}
{JOHNSON Mark},
\newblock 1987,
\newblock {\em The Body in the Mind: The Bodily Basis of Meaning, Imagination, and Reason},
\newblock Chicago: University of Chicago Press.

\bibitem[\protect\citename{Jones, }1987]{jon87}
{JONES Lindsay (red.)},
\newblock 1987,
\newblock {\em Encyclopedia of Religion, Volume~15},
\newblock New York: Macmillan Reference USA.

\bibitem[\protect\citename{Jordanskaja \& Mielczuk, }1988]{jormie88}
{JORDANSKAJA Lidia \& MIELCZUK Igor},
\newblock 1988,
\newblock ,,Konotacja w~semantyce lingwistycznej i~leksykografii'',
\newblock {[w:]} {BARTMIŃSKI Jerzy (red.)},
\newblock {\em Konotacja},
\newblock Lublin: UMCS,
\newblock s.~9--34.

\bibitem[\protect\citename{Jung, }1954]{jun54}
{JUNG Carl Gustav},
\newblock 1954,
\newblock ,,On the Psychology of the Trickster-figure'',
\newblock {[w:]} {RADIN Paul},
\newblock 1956,
\newblock {\em The Trickster: A~Study in American Indian Mythology},
\newblock London: Routledge \linebreak \& Kegan Paul,
\newblock s.~195--211.

\bibitem[\protect\citename{Jung, }1958a]{jun58a}
{JUNG Carl Gustav},
\newblock 1958a,
\newblock {\em The Collected Works of C. G. Jung, Volume~18},
\newblock New York: Pan- theon Books.

\bibitem[\protect\citename{Jung, }1958b]{jun58b}
{JUNG Carl Gustav},
\newblock 1958b,
\newblock ,,Jung and Religious Belief'',
\newblock {[w:]} {TENŻE},
\newblock {\em The Collected Works of C. G. Jung, Volume~18},
\newblock New York: Pantheon Books,
\newblock s.~702--744.

\bibitem[\protect\citename{Kalevala, }1964]{kal64}
{\emph{Kalevala}},
\newblock 1964,
\newblock Helsinki: Suomalaisen Kirjallisuuden Seura.

\bibitem[\protect\citename{Kalevala, }1998]{kal98}
{\emph{Kalevala}},
\newblock 1998,
\newblock tłum. {LITWINIUK Jerzy},
\newblock Warszawa: PIW.

%\bibitem[\protect\citename{Kalevala, }2012]{kal12}
%{\sc \emph{Kalevala}},
%\newblock 2012,
%\newblock tłum. {\sc Lubarskaja Aleksandra},
%\newblock Moskwa: Nigma.

\bibitem[\protect\citename{Kardela, }1992]{kar92}
{KARDELA Henryk},
\newblock 1992,
\newblock ,,Gramatyka kognitywna jako globalna teoria języka'', 
\newblock {[w:]} {NO- \linebreak WAKOWSKA-KEMPNA Iwona (red.)},
\newblock {\em Język a~Kultura, Tom~8: Podstawy metodologiczne semantyki współczesnej},
\newblock Wrocław: Wiedza o~Kulturze,
\newblock s.~9--22.

\bibitem[\protect\citename{Kardela, }2006]{kar06}
{KARDELA Henryk},
\newblock 2006,
\newblock ,,Schemat i prototyp w morfologii kognitywnej'',
\newblock {[w:]} {TENŻE, MUSZYŃSKI Zbysław \& RAJEWSKI Maciej (red.)},
\newblock {\em Kognitywistyka~2: Podobieństwo},
\newblock Lublin: UMCS,
\newblock s.~179--208.

\bibitem[\protect\citename{Kardela, Muszyński \& Rajewski, }2006]{karmusraj06}
{KARDELA Henryk, MUSZYŃSKI Zbysław \& RAJEWSKI Maciej (red.)},
\newblock 2006,
\newblock {\em Kognitywistyka~2: Podobieństwo},
\newblock Lublin: UMCS.
%\newblock s.~180.

\bibitem[\protect\citename{Kerényi, }2010]{ker10}
{KERÉNYI Karl},
\newblock 2010,
\newblock ,,Trickster a~mitologia grecka'',
\newblock {[w:]} {RADIN Paul},
\newblock {\em Trickster: Studium mitologii Indian północnoamerykańskich},
\newblock tłum. {TOPCZEWSKA Anna},
\newblock Warszawa: WUW,
\newblock s.~195--217.

\bibitem[\protect\citename{Klawiter, }2006]{kla06}
{KLAWITER Andrzej},
\newblock 2006,
\newblock {\em Kognitywistyka} [online], dostęp: 20.03.2014, pod adresem: {\small \texttt{<\url{http://www.staff.amu.edu.pl/~klawiter/Kogni-notka.htm}>}}.

\bibitem[\protect\citename{Kleiber, }1990]{kle90}
{KLEIBER Georges},
\newblock 1990,
\newblock {\em La sémantique du prototype: Catégories et sens lexical},
\newblock Paris: Presses Universitaires de France.
%\newblock s.~59.

\bibitem[\protect\citename{Krzeszowski, }1988]{krz88}
{KRZESZOWSKI Tomasz},
\newblock 1988,
\newblock ,,Wstęp do wydania polskiego'',
\newblock {[do:]} {LAKOFF George \linebreak \& JOHNSON Mark},
\newblock {\em Metafory w naszym życiu},
\newblock tłum. {KRZESZOWSKI Tomasz},
\newblock War- \linebreak szawa:~PIW,
\newblock s.~5--18.

\bibitem[\protect\citename{Kudzinowski, }1988]{kud88}
{KUDZINOWSKI Czesław},
\newblock 1988,
\newblock {\em Słownik fińsko-polski},
\newblock Poznań: UAM.
\newpage
\bibitem[\protect\citename{Lakoff, }1987]{lak87}
{LAKOFF George},
\newblock 1987,
\newblock {\em Women, Fire, and Dangerous Things: What Categories Reveal About the Mind},
\newblock Chicago: University of Chicago Press.
%\newblock s.~74.

\bibitem[\protect\citename{Lakoff \& Johnson, }1988]{lakjoh88}
{LAKOFF George \& JOHNSON Mark},
\newblock 1988,
\newblock {\em Metafory w naszym życiu},
\newblock tłum. {KRZESZOW- \linebreak SKI Tomasz},
\newblock Warszawa: PIW.
%\newblock s.~15.

\bibitem[\protect\citename{Langacker, }1987]{lan87}
{LANGACKER Ronald},
\newblock 1987,
\newblock {\em Foundations of Cognitive Grammar, Volume~1: Theoretical prerequisites},
\newblock California: Stanford University Press.
%\newblock s.~371.

\bibitem[\protect\citename{Langacker, }1988a]{lan88a}
{LANGACKER Ronald},
\newblock 1988a,
\newblock ,,A~usage-based model'',
\newblock {[w:]} {RUDZKA-OSTYN Brygida (red.)},
\newblock {\em Topics in Cognitive Linguistics},
\newblock Amsterdam: John Benjamins, %\linebreak
\newblock s.~127--161.

\bibitem[\protect\citename{Langacker, }1988b]{lan88b}
{LANGACKER Ronald},
\newblock 1988b,
\newblock ,,A~view of linguistic semantics'',
\newblock {[w:]} {RUDZKA-OSTYN Brygida (red.)},
\newblock {\em Topics in Cognitive Linguistics},
\newblock Amsterdam: John Benjamins,
\newblock s.~49--90.

\bibitem[\protect\citename{Langacker, }1991]{lan91}
{LANGACKER Ronald},
\newblock 1991,
\newblock {\em Concept, Image, and Symbol: The Cognitive Basis of Grammar},
\newblock Berlin: Mouton de Gruyter.
%\newblock s.~109.

\bibitem[\protect\citename{Langacker, }1999]{lan99}
{LANGACKER Ronald},
\newblock 1999,
\newblock ,,A~dynamic usage-based model'',
\newblock {[w:]} {BARLOW Michael \linebreak \& KEMMER Suzanne (red.)},
\newblock {\em Usage Based Models of Language},
\newblock Stanford: CSLI~Publications,
\newblock s.~1--63.

\bibitem[\protect\citename{Langacker, }2003]{lan03}
{LANGACKER Ronald},
\newblock 2003,
\newblock ,,Model dynamiczny oparty na uzusie językowym'',
\newblock tłum. \linebreak {KUBIŃSKI Wojciech},
\newblock {[w:]} {DĄBROWSKA Ewa \& KUBIŃSKI Wojciech (red.)},
\newblock {\em Akwi- \linebreak zycja języka w świetle językoznawstwa kognitywnego},
\newblock Kraków: Universitas, 
\newblock s.~30--117.

\bibitem[\protect\citename{Langacker, }2005]{lan05}
{LANGACKER Ronald},
\newblock 2005,
\newblock {\em Wykłady z gramatyki kognitywnej},
\newblock tłum. {KARDELA Henryk \& ŁOZOWSKI Przemysław},
\newblock Lublin: UMCS.
%\newblock s.~16--17.

\bibitem[\protect\citename{Leeming, }1990]{lee90}
{LEEMING David Adams},
\newblock 1990,
\newblock {\em The World of Myth: An Anthology},
\newblock New York: Oxford University Press.

\bibitem[\protect\citename{Lewicki, }1982]{lew82}
{LEWICKI Andrzej Maria},
\newblock 1982,
\newblock ,,O motywacji frazeologizmów'',
\newblock {[w:]} {BASAJ Mieczysław \& RYTEL-SCHWARZ Danuta (red.)},
\newblock {\em Z~Problemów Frazeologii Polskiej i~Słowiańskiej, Tom~1},
\newblock Wrocław: PAN,
\newblock s.~33--47.

\bibitem[\protect\citename{Łupkowski, Mozyrko \& Porowska, }2009]{lupmozpor09}
{ŁUPKOWSKI Paweł, MOZYRKO Bartosz \& POROWSKA Agnieszka (red.)},
\newblock 2009, %\linebreak
\newblock {\em Poznańskie Forum Kognitywistyczne: Teksty Pokonferencyjne, Nr~3},
\newblock Poznań: Zakład Logiki i~Kognitywistyki, Instytut Psychologii UAM.
%\newblock s.~31--39.

\bibitem[\protect\citename{Mendecki, }2000]{men00}
{MENDECKI Bartłomiej},
\newblock 2000,
\newblock {\em Boski posłaniec Hermes},
\newblock {[w:]} ,,ALBO Albo: Trikster'',
\newblock nr~4, \linebreak s.~133--135.

\bibitem[\protect\citename{Mosiołek-Kłosińska, }1992]{mos92}
{MOSIOŁEK-KŁOSIŃSKA Katarzyna},
\newblock 1992,
\newblock {\em Stereotypy psa zawarte w języku polskim}, \linebreak
\newblock {[w:]} ,,Poradnik Językowy'',
\newblock nr~4, s.~301--304.

\bibitem[\protect\citename{Murray, }1895]{mur95}
{MURRAY Alexander Stuart},
\newblock 1895,
\newblock {\em Manual of Mythology},
\newblock Philadelphia: David McKay.

\bibitem[\protect\citename{Nowakowska-Kempna, }1992]{now92}
{NOWAKOWSKA-KEMPNA Iwona (red.)},
\newblock 1992,
\newblock {\em Język a~Kultura, Tom~8: Podstawy me- todologiczne semantyki współczesnej},
\newblock Wrocław: Wiedza o~Kulturze.
%\newblock s.~10--11.

\bibitem[\protect\citename{Osherson \& Smith, }1981]{oshsmi81}
{OSHERSON Daniel \& SMITH Edward},
\newblock 1981,
\newblock {\em On the adequacy of prototype theory as a~theory of concepts},
\newblock {[w:]} ,,Cognition'',
\newblock nr~9, s.~59--72.

\bibitem[\protect\citename{Otto, }2001]{ott01}
{OTTO Beatrice},
\newblock 2001,
\newblock {\em Fools Are Everywhere: The Court Jester around the World},
\newblock Chicago: University of Chicago Press.

\bibitem[\protect\citename{Pelton, }1980]{pel80}
{PELTON Robert},
\newblock 1980,
\newblock {\em The Trickster in West Africa: A Study of Mythic Irony and Sacred Delight},
\newblock Berkeley: University of California Press.

\bibitem[\protect\citename{Pisarkowa, }1976]{pis76}
{PISARKOWA Krystyna},
\newblock 1976,
\newblock {\em Konotacja semantyczna nazw narodowości},
\newblock {[w:]} ,,Zeszyty Prasoznawcze'',
\newblock nr~1, s.~5--26.

\bibitem[\protect\citename{Biblia, }2003]{bib03}
{\emph{Pismo święte Starego i~Nowego Testamentu: Biblia Tysiąclecia}},
\newblock 2003,
\newblock Poznań: Pallottinum.

\bibitem[\protect\citename{Putnam, }1975]{put75}
{PUTNAM Hilary},
\newblock 1975,
\newblock {\em Philosophical Papers, Volume~2: Mind, Language and Reality},
\newblock London: Cambridge University Press.
%\newblock no.~12, s.~249.

\bibitem[\protect\citename{Radin, }1956]{rad56}
{RADIN Paul},
\newblock 1956,
\newblock {\em The Trickster: A~Study in American Indian Mythology},
\newblock London: Rout- \linebreak ledge \& Kegan Paul.

\bibitem[\protect\citename{Radin, }2010]{rad10}
{RADIN Paul},
\newblock 2010,
\newblock {\em Trickster: Studium mitologii Indian północnoamerykańskich},
\newblock tłum. {TOPCZEWSKA Anna},
\newblock Warszawa: WUW.

\bibitem[\protect\citename{Ricketts, }1987]{ric87}
{RICKETTS Mac Linscott},
\newblock 1987,
\newblock ,,North American Tricksters'',
\newblock {[w:]} {JONES Lindsay (red.)},
\newblock {\em Encyclopedia of Religion, Volume 15},
\newblock New York: Macmillan Reference USA,
\newblock s.~48--51.

\bibitem[\protect\citename{Rosch, }1977]{ros77}
{ROSCH Eleanor},
\newblock 1977,
\newblock ,,Human Categorization'',
\newblock {[w:]} {WARREN Neil (red.)},
\newblock {\em Studies in Cross-Cultural Psychology, Volume~1},
\newblock London: Academic Press,
\newblock s.~1--72.

\bibitem[\protect\citename{Rudzka-Ostyn, }1988]{rud88}
{RUDZKA-OSTYN Brygida (red.)},
\newblock 1988,
\newblock {\em Topics in Cognitive Linguistics},
\newblock Amsterdam: John Benjamins.
%\newblock s.~49--50, 63--64, 127--161.

\bibitem[\protect\citename{Sieradzan, }2005]{sie05}
{SIERADZAN Jacek},
\newblock 2005,
\newblock {\em Szaleństwo w~religiach świata},
\newblock Kraków: Wydawnictwo Wanda.

\bibitem[\protect\citename{Stokłosa, }2012]{sto12}
{STOKŁOSA Bożenna},
\newblock 2012,
\newblock {\em Trikster --- postać mityczna, archetyp i figura twórcy}, \linebreak
\newblock {[w:]} ,,Magazyn Sztuki'',
\newblock nr~3, s.~119--130.

\bibitem[\protect\citename{Sullivan, }1987]{sul87}
{SULLIVAN Lawrence},
\newblock 1987,
\newblock ,,Tricksters: An Overview'',
\newblock {[w:]} {JONES Lindsay (red.)},
\newblock {\em Encyclopedia of Religion, Volume 15},
\newblock New York: Macmillan Reference USA,
\newblock s.~45--46.

\bibitem[\protect\citename{Tabakowska, }2001]{tab01}
{TABAKOWSKA Elżbieta},
\newblock 2001,
\newblock {\em Kognitywne podstawy języka i~językoznawstwa},
\newblock Kraków: Universitas.
%\newblock s.~54.

\bibitem[\protect\citename{Taylor, }2001]{tay01}
{TAYLOR John},
\newblock 2001,
\newblock {\em Kategoryzacja w języku: Prototypy w teorii językoznawczej},
\newblock tłum. {SKUCIŃSKA Anna},
\newblock Kraków: Universitas.
%\newblock s.~131.

\bibitem[\protect\citename{Tsohatzidis, }1990]{tso90}
{TSOHATZIDIS Savas (red.)},
\newblock 1990,
\newblock {\em Meanings and Prototypes: Studies in Linguistic Cat- \linebreak egorization},
\newblock London: Routledge.
%\newblock s.~17--47, 347--367.

\bibitem[\protect\citename{Turner, }1987]{tur87}
{TURNER Mark},
\newblock 1987,
\newblock {\em Death Is The Mother of Beauty: Mind, Metaphor, Criticism},
\newblock Chicago: University of Chicago Press.

\bibitem[\protect\citename{Wacewicz, }2006]{wac06}
{WACEWICZ Sławomir},
\newblock 2006,
\newblock {\em Kognitywistyka} [online], dostęp: 20.03.2014, pod adresem:
{\small \texttt{<\url{http://www.kognitywistyka.net/kognitywistyka/prezentacja.html}>}}.

\bibitem[\protect\citename{Wagner, }1910]{wag10}
{WAGNER Richard},
\newblock 1910,
\newblock {\em The Rhinegold and the Valkyrie},
\newblock London: William Heinemann.

\bibitem[\protect\citename{Warren, }1977]{war77}
{WARREN Neil (red.)},
\newblock 1977,
\newblock {\em Studies in Cross-Cultural Psychology, Volume~1},
\newblock London: Academic Press.
%\newblock s.~1--72.

\bibitem[\protect\citename{Wasilewski, }0000]{was00}
{WASILEWSKI Jerzy Sławomir},
\newblock w~druku,
\newblock {\em Niepoważność: Od mitologii śmiechu do antropologii komizmu}.

\bibitem[\protect\citename{Welch, }2011]{wel11}
{WELCH Lynda},
\newblock 2011,
\newblock {\em Goddess of the North},
\newblock Newburyport: Red Wheel.

\bibitem[\protect\citename{Wierzbicka, }1989]{wie89}
{WIERZBICKA Anna},
\newblock 1989,
\newblock {\em Prototypes in semantic and pragmatic: Explicating attitudinal meanings in terms of prototypes},
\newblock {[w:]} ,,Linguistics'',
\newblock nr~27, s.~731--769.

\bibitem[\protect\citename{Wierzbicka, }1990]{wie90}
{WIERZBICKA Anna},
\newblock 1990,
\newblock ,,»Prototype save«: On the uses and abuses of the notion of »prototype« in linguistics and related fields'',
\newblock {[w:]} {TSOHATZIDIS Savas (red.)},
\newblock {\em Meanings and Prototypes: Studies in Linguistic Categorization},
\newblock London: Routledge, \linebreak
\newblock s.~347--367.

\bibitem[\protect\citename{Wierzbicka, }1992]{wie92}
{WIERZBICKA Anna},
\newblock 1992,
\newblock ,,Uniwersalne pojęcia ludzkie i~ich konfiguracje w różnych kulturach'',
\newblock {[w:]} {BARTMIŃSKI Jerzy (red.)},
\newblock {\em Etnolingwistyka, Tom~4},
\newblock Lublin: UMCS, \linebreak
\newblock s.~7--40.

\bibitem[\protect\citename{Wittgenstein, }1953]{wit53}
{WITTGENSTEIN Ludwig},
\newblock 1953,
\newblock {\em Philosophical Investigations},
\newblock Oxford: Blackwell.

\bibitem[\protect\citename{Wittgenstein, }1972]{wit72}
{WITTGENSTEIN Ludwig},
\newblock 1972,
\newblock {\em Dociekania filozoficzne},
\newblock tłum. {WOLNIEWICZ Bogusław},
\newblock Warszawa: PWN.

\bibitem[\protect\citename{Zadeh, }1982]{zad82}
{ZADEH Lotfi Askar},
\newblock 1982,
\newblock {\em A~note on prototype theory and fuzzy sets},
\newblock {[w:]} ,,Cognition'',
\newblock nr~12, s.~291--297.

\bibitem[\protect\citename{Ziemiański, }2012]{zie12}
{ZIEMIAŃSKI Andrzej},
\newblock 2012,
\newblock {\em Pomnik Cesarzowej Achai, Tom 1},
\newblock Lublin: Fabryka Słów.

\end{thebibliography}

%\vspace*{13\baselineskip}

%\begin{center}
%\textbf{Summary}\\
%\end{center}
%\indent
%{\selectlanguage{english}
%The approach of cognitive linguistics gives us a~possibility of carrying out research on concepts and phenomena, which are often ambiguous and fuzzy. One of them is \emph{trick\-ster} --- a~divine madman, whose presence in the Kalevala is clearly visible, al\-though this subject has not yet been studied in the field of linguistics. A~thorough analysis of trickster domains, preceded by a~detailed presentation of cognitive science and complicated nature of the fool, proves that the Kalevala is full of trickster-figures, which are characterized mainly by magic, stupidity and a~link to a~deity. The tendencies in the distribution of domains allow us to mark the kalevalic tricksters as \emph{major} and \emph{minor} ones and to create a~concept of a~\emph{new kalevalic trickster} using their additional characteristics.\\

{
\setlength{\voffset}{0cm}
\setlength{\hoffset}{0cm}
\includepdf[pages=-,scale=0.75,offset=25 0]{ocena1.pdf}
\includepdf[pages=-,scale=0.75,offset=25 0]{ocena2.pdf}
}
%\setlength{\voffset}{-2.54cm}
%\setlength{\hoffset}{-2.54cm}
\end{document}
«»

\begin{description}
  \item[TEZA A:] ,,Struktura semantyczna jest redukowalna do struktury kognitywnej (czyli mentalnego doświadczenia)'' --- struktura semantyczna jest ,,zanurzona'' w~wiedzy o~świecie, a~znaczenia ,,zależne są w~sposób oczywisty od tej wiedzy, czyli od szeroko pojętego kontekstu, tak lingwistycznego, jak i~»pozajęzykowego«'' (por. Langacker~1991:~109).
  \item[TEZA B:] ,,Często używane wyrażenie przywołuje sieć powiązanych ze sobą znaczeń (sensów)'' --- wyrażenia języka cechuje polisemia oraz stanowią one węzeł dostępu do danej sieci znaczeń. Elementy kategorii, które odchodzą od prototypu dają podwaliny tzw. \textit{kategoriom radialnym} (ang. \textit{radial categories}), razem z~elementami prototypowymi (centralnymi) oraz peryferyjnymi (por. Tabakowska 2001:~54).
  \item[TEZA C:] ,,Struktury semantyczne określane są względem domen kognitywnych'' --- dana kategoria jest tym bardziej prototypowa, im więcej modeli z~tzw. \textit{wiązki modeli kognitywnych} (ang. \textit{cluster models}) jest przywoływanych dla zdefiniowania znaczenia danego słowa. Wyidealizowanymi modelami kognitywnymi (ang. \textit{Idealized Cognitive Models}, \textit{ICMs}) --- na przykładzie modelu matki (m.in. jako ‘najbliższego przodka płci żeńskiej’ w~modelu genealogicznym, ‘żony ojca’ w~modelu małżeńskim czy też ‘dorosłej kobiety, która wychowuje i~jest żywicielką dziecka’ w~myśl modelu żywieniowego) --- zajmuje się m.in. Lakoff (1987: 74).
  \item[TEZA D:] ,,Struktura semantyczna uzyskuje swój status w~wyniku nałożenia profilu (desygnatu) na bazę'' --- wyrażenie zyskuje swoje znaczenie poprzez \emph{profilowanie}, nazywane również \emph{perspektywizacją} (Taylor 2001:~131), stanowiące wyodrębnianie określonych struktur z~tzw. \textit{bazy kognitywnej}. Dobrym przykładem jest proces definiowania ,,przeciwprostokątnej'', który nie ma prawa bytu bez odwołania się do pojęcia figury ,,trójkąta''.
  \item[TEZA E:] ,,Struktury semantyczne zawierają \emph{skonwencjonalizowane obrazowanie} (ang. \emph{conventional imagery})'' --- według Langackera (1988b:~63--64) scena (sytuacja) jest konstruowana na wiele różnych sposobów, przy czym zwraca się uwagę na aspekty takie jak: poziom uszczegółowienia, założenia dotyczące tła i~oczekiwań, skala i~zakres predykacji, perspektywa oraz subiektywizacja i~sekundarna aktywizacja. Pierwszy aspekt dotyczy precyzji definicji, widoczny np. w~ciągu \emph{zwierzę} $\textrm{\ding{221}}$ \emph{ssak} $\textrm{\ding{221}}$ \emph{gryzoń} $\textrm{\ding{221}}$ \emph{chomik} $\textrm{\ding{221}}$ \emph{chomik syryjski}, w którym wyrażenie ostatnie jest najdokładniejsze.
\end{description}





